\input{preamble}

% OK, start here.
%
\begin{document}

\title{Faisceaux sur les espaces}


\maketitle

\phantomsection
\label{section-phantom}

\tableofcontents

\section{Introduction}
\label{section-introduction}

\noindent
Les propriétés fondamentales des faisceaux sur les espaces topologiques
seront exposées dans ce chapitre. On pourra consulter \cite{Godement}.

\medskip\noindent
La discussion des faisceaux sur les sites donnée plus loin dans ces documents
remplacera cet exposé. Mais il est peut-être utile de définir brièvement ici
certaines des notions en jeu.











\section{Notions de base}
\label{section-sheaves-basic}

\noindent
Voici une liste de notions fondamentales de topologie.

\begin{enumerate}
\item Soit $X$ un espace topologique. L'expression : « Soit
$U = \bigcup_{i \in I} U_i$ un recouvrement ouvert » signifie ce qui
suit : $I$ est un ensemble et, pour chaque $i \in I$, on se donne
une partie ouverte $U_i \subset X$ telle que $U$ soit la
réunion des $U_i$. On autorise $I = \emptyset$ ;
dans ce cas, il n'y a aucun $U_i$ et $U = \emptyset$.
Lorsque $I \not = \emptyset$, il est également permis que
certains, voire tous, les $U_i$ soient vides.
\item etc., etc.
\end{enumerate}












\section{Préfaisceaux}
\label{section-presheaves}


\begin{definition}
\label{definition-presheaf}
Soit $X$ un espace topologique.
\begin{enumerate}
\item Un {\it préfaisceau d'ensembles $\mathcal{F}$ sur $X$} est une règle qui
associe à chaque ouvert $U \subset X$ un ensemble $\mathcal{F}(U)$ et
à chaque inclusion $V \subset U$ une application
$\rho^U_V : \mathcal{F}(U) \to \mathcal{F}(V)$ telle que
$\rho^U_U = \text{id}_{\mathcal{F}(U)}$ et que,
si $W \subset V \subset U$, on ait
$\rho^U_W = \rho^V_W \circ \rho ^U_V$.
\item Un {\it morphisme $\varphi : \mathcal{F} \to \mathcal{G}$
de préfaisceaux d'ensembles sur $X$} est une règle qui associe à chaque
ouvert $U \subset X$ une application d'ensembles $\varphi : \mathcal{F}(U)
\to \mathcal{G}(U)$ compatible avec les applications de restriction,
c'est-à-dire que, lorsque $V \subset U \subset X$ sont ouverts, le
diagramme
$$
\xymatrix{
\mathcal{F}(U) \ar[r]^\varphi \ar[d]^{\rho^U_V} &
\mathcal{G}(U) \ar[d]^{\rho^U_V} \\
\mathcal{F}(V) \ar[r]^\varphi & \mathcal{G}(V)
}
$$
est commutatif.
\item La catégorie des préfaisceaux d'ensembles sur $X$ sera notée
$\textit{PSh}(X)$.
\end{enumerate}
\end{definition}

\noindent
Les éléments de l'ensemble $\mathcal{F}(U)$ sont appelés
les {\it sections} de $\mathcal{F}$ sur $U$.
Pour tout $V \subset U$, l'application
$\rho^U_V : \mathcal{F}(U) \to \mathcal{F}(V)$
est appelée l'{\it application de restriction}. Nous emploierons la
notation $s|_V := \rho^U_V(s)$ si $s\in \mathcal{F}(U)$.
Cette notation est compatible avec la notion topologique de restriction
des fonctions, car si $W \subset V \subset U$
et si $s$ est une section de $\mathcal{F}$ sur $U$, alors
$s|_W = (s|_V)|_W$, par la propriété des applications de restriction
exprimée dans la définition ci-dessus.

\medskip\noindent
Une autre notation souvent employée consiste à désigner les sections
sur un ouvert $U$ par le symbole $\Gamma(U, -)$ ou par
$H^0(U, -)$. Autrement dit, les égalités suivantes
sont tautologiques :
$$
\Gamma(U, \mathcal{F}) = \mathcal{F}(U) = H^0(U, \mathcal{F}).
$$
Nous n'emploierons pas cette notation dans ce chapitre, mais nous
l'emploierons dans d'autres chapitres.

\begin{definition}
\label{definition-constant-presheaf}
Soit $X$ un espace topologique. Soit $A$ un ensemble.
Le {\it préfaisceau constant de valeur $A$} est le
préfaisceau qui associe l'ensemble $A$ à tout ouvert
$U \subset X$ et dont toutes les applications de restriction
sont $\text{id}_A$.
\end{definition}

\section{Préfaisceaux abéliens}
\label{section-abelian-presheaves}

\noindent
Dans cette section, nous indiquons brièvement quelques propriétés de la
catégorie des préfaisceaux qui permettent de définir les préfaisceaux
de groupes abéliens.

\begin{example}
\label{example-singleton-presheaf}
Soit $X$ un espace topologique. Considérons une règle $\mathcal{F}$ qui
associe un singleton à toute partie ouverte de $X$. Comme tout ensemble admet
une unique application vers un singleton, les applications de restriction
$\rho^U_V$ existent et sont uniques. La structure ainsi obtenue est un
préfaisceau d'ensembles sur $X$. C'est un objet final de la catégorie des
préfaisceaux d'ensembles sur $X$, en vertu de la propriété des singletons
mentionnée ci-dessus. Il est donc
également unique à isomorphisme unique près. Nous désignerons parfois ce
préfaisceau par $*$.
\end{example}

\begin{lemma}
\label{lemma-product-presheaves}
Soit $X$ un espace topologique. La catégorie des préfaisceaux d'ensembles
sur $X$ admet des produits (voir
Catégories, Définition \ref{categories-definition-product}).
De plus, l'ensemble des
sections du produit $\mathcal{F} \times \mathcal{G}$
sur un ouvert $U$ est le produit des ensembles de sections de
$\mathcal{F}$ et $\mathcal{G}$ sur $U$.
\end{lemma}

\begin{proof}
En effet, supposons que $\mathcal{F}$ et $\mathcal{G}$ soient des
préfaisceaux d'ensembles sur l'espace topologique $X$.
Considérons la règle $U \mapsto \mathcal{F}(U) \times \mathcal{G}(U)$,
notée $\mathcal{F} \times \mathcal{G}$. Si $V \subset U \subset X$
sont ouverts, définissons l'application de restriction
$$
(\mathcal{F} \times \mathcal{G})(U)
\longrightarrow
(\mathcal{F} \times \mathcal{G})(V)
$$
par la formule $(s, t) \mapsto (s|_V, t|_V)$. Il est alors immédiatement
clair que $\mathcal{F} \times \mathcal{G}$ est un préfaisceau.
Il existe en outre des morphismes de projection
$p : \mathcal{F} \times \mathcal{G} \to \mathcal{F}$
et
$q : \mathcal{F} \times \mathcal{G} \to \mathcal{G}$.
Nous laissons au lecteur le soin de montrer que,
pour tout autre préfaisceau $\mathcal{H}$, on a
$\Mor(\mathcal{H}, \mathcal{F} \times \mathcal{G})
= \Mor(\mathcal{H}, \mathcal{F}) \times
\Mor(\mathcal{H}, \mathcal{G})$.
\end{proof}

\noindent
Rappelons que si $(A, + : A \times A \to A, - : A \to A, 0\in A)$
est un groupe abélien, alors le zéro et l'application de passage à l'opposé
sont déterminés de manière unique par la loi d'addition. Autrement dit, il
est légitime de dire : « Soit $(A, +)$ un groupe abélien. »

\begin{lemma}
\label{lemma-abelian-presheaves}
Soit $X$ un espace topologique.
Soit $\mathcal{F}$ un préfaisceau d'ensembles.
Considérons les types de structures suivants sur $\mathcal{F}$ :
\begin{enumerate}
\item Pour tout ouvert $U$, une structure de groupe abélien
sur $\mathcal{F}(U)$ telle que toutes les applications de restriction soient des
homomorphismes de groupes abéliens.
\item Un morphisme de préfaisceaux
$+ : \mathcal{F} \times \mathcal{F} \to \mathcal{F}$,
un morphisme de préfaisceaux $- : \mathcal{F} \to \mathcal{F}$
et un morphisme $0 : * \to \mathcal{F}$
(voir l'Exemple \ref{example-singleton-presheaf})
satisfaisant tous les axiomes vérifiés par $+, -, 0$ dans un
groupe abélien usuel.
\item Un morphisme de préfaisceaux
$+ : \mathcal{F} \times \mathcal{F} \to \mathcal{F}$,
un morphisme de préfaisceaux $- : \mathcal{F} \to \mathcal{F}$
et un morphisme $0 : * \to \mathcal{F}$
tels que, pour chaque ouvert $U \subset X$, le quadruplet
$(\mathcal{F}(U), +, -, 0)$ soit un groupe abélien,
\item Un morphisme de préfaisceaux $+ : \mathcal{F} \times \mathcal{F}
\to \mathcal{F}$ tel que, pour tout ouvert $U \subset X$,
l'application $+ : \mathcal{F}(U) \times \mathcal{F}(U) \to \mathcal{F}(U)$
définisse une structure de groupe abélien.
\end{enumerate}
Il existe des bijections naturelles entre les collections de
données des types (1) à (4) ci-dessus.
\end{lemma}

\begin{proof}
Omis.
\end{proof}

\noindent
Le lemme affirme que munir $\mathcal{F}$ d'une structure d'objet en groupes
abéliens dans la catégorie des préfaisceaux revient à se donner un préfaisceau
d'ensembles
$\mathcal{F}$ tel que tous les ensembles $\mathcal{F}(U)$ soient munis
d'une structure de groupe abélien et que toutes les applications de restriction
soient des homomorphismes de groupes. Pour la plupart des structures algébriques,
nous procéderons de même pour les (pré)faisceaux de tels objets : autrement dit,
nous définirons un (pré)faisceau de tels objets comme un (pré)faisceau
$\mathcal{F}$ d'ensembles dont tous les ensembles de sections $\mathcal{F}(U)$
sont munis de cette structure d'une manière compatible avec les applications
de restriction.

\begin{definition}
\label{definition-abelian-presheaves}
Soit $X$ un espace topologique.
\begin{enumerate}
\item Un {\it préfaisceau de groupes abéliens sur $X$}, ou un
{\it préfaisceau abélien sur $X$},
est un préfaisceau d'ensembles $\mathcal{F}$ tel que, pour chaque ouvert
$U \subset X$, l'ensemble $\mathcal{F}(U)$ soit muni
d'une structure de groupe abélien et que toutes les applications de restriction
$\rho^U_V$ soient des homomorphismes de groupes abéliens ; voir le
Lemme \ref{lemma-abelian-presheaves} ci-dessus.
\item Un {\it morphisme de préfaisceaux abéliens sur $X$}
$\varphi : \mathcal{F} \to \mathcal{G}$ est un morphisme de préfaisceaux
d'ensembles qui induit
un homomorphisme de groupes abéliens $\mathcal{F}(U) \to \mathcal{G}(U)$
pour tout ouvert $U \subset X$.
\item La catégorie des préfaisceaux de groupes abéliens sur $X$ est notée
$\textit{PAb}(X)$.
\end{enumerate}
\end{definition}

\begin{example}
\label{example-direct-sum-points}
Soit $X$ un espace topologique. Pour chaque $x \in X$, soit
$M_x$ un groupe abélien. Pour $U \subset X$ ouvert,
posons
$$
\mathcal{F}(U) = \bigoplus\nolimits_{x \in U} M_x.
$$
Un élément de ce groupe abélien sera noté sous la forme
$\sum_{i = 1}^n m_{x_i}$, où $x_i \in U$ et $m_{x_i} \in M_{x_i}$.
(Bien entendu, nous pouvons toujours choisir cette écriture de telle sorte que
$x_1, \ldots, x_n$ soient deux à deux distincts.)
Pour des ouverts $V \subset U \subset X$, nous définissons une application de
restriction $\mathcal{F}(U) \to \mathcal{F}(V)$ qui
envoie un élément $s = \sum_{i = 1}^n m_{x_i}$
sur l'élément $s|_V = \sum_{x_i \in V} m_{x_i}$.
Nous laissons au lecteur le soin de vérifier qu'il s'agit d'un
préfaisceau de groupes abéliens.
\end{example}

\section{Préfaisceaux de structures algébriques}
\label{section-presheaves-structures}

\noindent
Précisons la définition
des préfaisceaux de structures algébriques.
Supposons que $\mathcal{C}$ soit une catégorie et
que $F : \mathcal{C} \to \textit{Sets}$ soit
un foncteur fidèle. En général, $F$ est un « foncteur
d'oubli ». Pour un objet $M \in \Ob(\mathcal{C})$,
nous appelons souvent $F(M)$ l'{\it ensemble sous-jacent} de
l'objet $M$. Si $M \to M'$ est un morphisme de $\mathcal{C}$,
nous appelons $F(M) \to F(M')$ l'{\it application ensembliste sous-jacente}.
En fait, nous ne distinguerons souvent pas un objet
de son ensemble sous-jacent, et de même pour les morphismes.
Nous dirons donc qu'une application d'ensembles $F(M) \to F(M')$
est un {\it morphisme de structures algébriques} si elle est
égale à $F(f)$ pour un certain morphisme $f : M \to M'$
de $\mathcal{C}$.

\medskip\noindent
Par analogie avec la définition \ref{definition-abelian-presheaves}
ci-dessus, un « préfaisceau d'objets de $\mathcal{C}$ » pourrait être
défini par les données suivantes :
\begin{enumerate}
\item un préfaisceau d'ensembles $\mathcal{F}$, et
\item pour tout ouvert $U \subset X$, le choix
d'un objet $A(U) \in \Ob(\mathcal{C})$,
\end{enumerate}
soumises aux conditions suivantes (avec la terminologie ci-dessus) :
\begin{enumerate}
\item pour tout ouvert $U \subset X$, l'ensemble $\mathcal{F}(U)$
est l'ensemble sous-jacent de $A(U)$, et
\item pour tous ouverts $V \subset U \subset X$,
l'application d'ensembles $\rho_V^U: \mathcal{F}(U) \to \mathcal{F}(V)$
est un morphisme de structures algébriques.
\end{enumerate}
Autrement dit, pour tous ouverts $V \subset U$ de $X$,
l'application de restriction $\rho^U_V$ est l'image
$F(\alpha^U_V)$ d'un unique morphisme
$\alpha^U_V : A(U) \to A(V)$ de la catégorie $\mathcal{C}$.
L'unicité découle de ce que $F$ est
fidèle ; elle implique également que
$\alpha^U_W = \alpha^V_W \circ \alpha^U_V$
lorsque $W \subset V \subset U$ sont des ouverts de $X$.
Le système $(A(-), \alpha^U_V)$ est ce que nous définirons comme un
préfaisceau à valeurs dans $\mathcal{C}$ sur $X$ ; comparer
Sites, définition \ref{sites-definition-presheaf}.
Nous retrouvons notre préfaisceau d'ensembles $(\mathcal{F}, \rho_V^U)$
au moyen des règles $\mathcal{F}(U) = F(A(U))$ et
$\rho_V^U = F(\alpha_V^U)$.

\begin{definition}
\label{definition-presheaf-values-in-category}
Soit $X$ un espace topologique.
Soit $\mathcal{C}$ une catégorie.
\begin{enumerate}
\item Un {\it préfaisceau $\mathcal{F}$ sur $X$ à valeurs dans $\mathcal{C}$}
est la donnée d'une règle qui associe à tout ouvert $U \subset X$
un objet $\mathcal{F}(U)$ de $\mathcal{C}$
et à toute inclusion $V \subset U$
un morphisme $\rho_V^U : \mathcal{F}(U) \to \mathcal{F}(V)$
de $\mathcal{C}$ tel que, lorsque $W \subset V \subset U$,
on ait $\rho_W^U = \rho_W^V \circ \rho_V^U$.
\item Un {\it morphisme $\varphi : \mathcal{F} \to \mathcal{G}$
de préfaisceaux à valeurs dans $\mathcal{C}$} est donné par un
morphisme $\varphi : \mathcal{F}(U) \to \mathcal{G}(U)$
de $\mathcal{C}$ compatible avec les morphismes de restriction.
\end{enumerate}
\end{definition}

\begin{definition}
\label{definition-underlying-presheaf-sets}
Soit $X$ un espace topologique. Soit $\mathcal{C}$ une catégorie.
Soit $F : \mathcal{C} \to \textit{Sets}$ un foncteur fidèle.
Soit $\mathcal{F}$ un préfaisceau sur $X$ à valeurs dans $\mathcal{C}$.
Le préfaisceau d'ensembles $U \mapsto F(\mathcal{F}(U))$
est appelé le {\it préfaisceau d'ensembles sous-jacent à $\mathcal{F}$}.
\end{definition}

\noindent
Il est d'usage d'employer la même lettre $\mathcal{F}$ pour désigner
le préfaisceau d'ensembles sous-jacent, et cela est
justifié par la discussion qui précède la
définition \ref{definition-presheaf-values-in-category}.
En particulier, l'expression « soit $s \in \mathcal{F}(U)$ »
ou « soit $s$ une section de $\mathcal{F}$ au-dessus de $U$ » signifie
que $s \in F(\mathcal{F}(U))$.

\medskip\noindent
Cette notation et ces définitions s'appliquent en particulier aux
{\it préfaisceaux de groupes (non nécessairement abéliens), d'anneaux, de modules
sur un anneau fixé, d'espaces vectoriels sur un corps fixé, } etc., ainsi qu'aux
{\it morphismes entre ceux-ci}.

\section{Préfaisceaux de modules}
\label{section-presheaves-modules}

\noindent
Supposons que $\mathcal{O}$ soit un préfaisceau d'anneaux sur $X$.
Nous voudrions définir la notion de préfaisceau de
$\mathcal{O}$-modules sur $X$. Par analogie avec la définition
\ref{definition-abelian-presheaves}, nous serions tentés de le définir
comme un préfaisceau d'ensembles $\mathcal{F}$ tel que, pour tout ouvert
$U \subset X$, l'ensemble $\mathcal{F}(U)$ soit muni d'une structure
de $\mathcal{O}(U)$-module compatible avec les applications de restriction
(de $\mathcal{F}$ et de $\mathcal{O}$). Il est cependant d'usage
(et équivalent) d'adopter la définition suivante.

\begin{definition}
\label{definition-presheaf-modules}
Soit $X$ un espace topologique, et soit $\mathcal{O}$
un préfaisceau d'anneaux sur $X$.
\begin{enumerate}
\item Un {\it préfaisceau de $\mathcal{O}$-modules}
est la donnée d'un préfaisceau abélien $\mathcal{F}$ et d'une
application de préfaisceaux d'ensembles
$$
\mathcal{O} \times \mathcal{F} \longrightarrow \mathcal{F}
$$
telle que, pour tout ouvert $U \subset X$, l'application
$\mathcal{O}(U) \times \mathcal{F}(U) \to \mathcal{F}(U)$
définisse une structure de $\mathcal{O}(U)$-module
sur le groupe abélien $\mathcal{F}(U)$.
\item Un {\it morphisme $\varphi : \mathcal{F} \to \mathcal{G}$
de préfaisceaux de $\mathcal{O}$-modules} est un morphisme de préfaisceaux abéliens
$\varphi : \mathcal{F} \to \mathcal{G}$ tel que
le diagramme
$$
\xymatrix{
\mathcal{O} \times \mathcal{F} \ar[r] \ar[d]_{\text{id} \times \varphi} &
\mathcal{F} \ar[d]^{\varphi} \\
\mathcal{O} \times \mathcal{G} \ar[r] &
\mathcal{G}
}
$$
soit commutatif.
\item L'ensemble des morphismes de $\mathcal{O}$-modules ci-dessus est
noté $\Hom_\mathcal{O}(\mathcal{F}, \mathcal{G})$.
\item La catégorie des préfaisceaux de $\mathcal{O}$-modules
est notée $\textit{PMod}(\mathcal{O})$.
\end{enumerate}
\end{definition}

\noindent
Supposons que $\mathcal{O}_1 \to \mathcal{O}_2$ soit un
morphisme de préfaisceaux d'anneaux sur $X$. Dans ce cas,
si $\mathcal{F}$ est un préfaisceau de $\mathcal{O}_2$-modules,
nous pouvons considérer $\mathcal{F}$ comme un préfaisceau de
$\mathcal{O}_1$-modules en utilisant la composée
$$
\mathcal{O}_1 \times \mathcal{F}
\to
\mathcal{O}_2 \times \mathcal{F}
\to
\mathcal{F}.
$$
Nous notons parfois ce préfaisceau $\mathcal{F}_{\mathcal{O}_1}$
pour indiquer la restriction des scalaires. Nous appelons ceci
la {\it restriction de $\mathcal{F}$}. Nous obtenons le
foncteur de restriction
$$
\textit{PMod}(\mathcal{O}_2)
\longrightarrow
\textit{PMod}(\mathcal{O}_1)
$$

\medskip\noindent
D'autre part, étant donné un préfaisceau de $\mathcal{O}_1$-modules
$\mathcal{G}$,
nous pouvons construire un préfaisceau de $\mathcal{O}_2$-modules
$\mathcal{O}_2 \otimes_{p, \mathcal{O}_1} \mathcal{G}$
par la règle
$$
\left(\mathcal{O}_2 \otimes_{p, \mathcal{O}_1} \mathcal{G}\right)(U)
=
\mathcal{O}_2(U) \otimes_{\mathcal{O}_1(U)} \mathcal{G}(U)
$$
L'indice $p$ signifie « préfaisceau » et non « point ».
Ce préfaisceau est appelé le {\it préfaisceau produit tensoriel}. Nous obtenons
le foncteur d'{\it extension des scalaires}
$$
\textit{PMod}(\mathcal{O}_1)
\longrightarrow
\textit{PMod}(\mathcal{O}_2)
$$

\begin{lemma}
\label{lemma-adjointness-tensor-restrict-presheaves}
Avec $X$, $\mathcal{O}_1$, $\mathcal{O}_2$, $\mathcal{F}$ et
$\mathcal{G}$ comme ci-dessus, il existe une bijection canonique
$$
\Hom_{\mathcal{O}_1}(\mathcal{G}, \mathcal{F}_{\mathcal{O}_1})
=
\Hom_{\mathcal{O}_2}(
\mathcal{O}_2 \otimes_{p, \mathcal{O}_1} \mathcal{G},
\mathcal{F}
)
$$
Autrement dit, les foncteurs de restriction et d'extension des scalaires
sont adjoints.
\end{lemma}

\begin{proof}
Cela résulte du fait que, pour un homomorphisme d'anneaux
$A \to B$, le foncteur de restriction et le foncteur
d'extension des scalaires sont adjoints.
\end{proof}





\section{Faisceaux}
\label{section-sheaves}

\noindent
Dans cette section, nous expliquons la condition de faisceau.

\begin{definition}
\label{definition-sheaf}
Soit $X$ un espace topologique.
\begin{enumerate}
\item Un {\it faisceau $\mathcal{F}$ d'ensembles sur $X$} est un préfaisceau
d'ensembles qui satisfait la propriété supplémentaire suivante : étant donnés
un recouvrement ouvert quelconque $U = \bigcup_{i \in I} U_i$ et une famille
de sections $s_i \in \mathcal{F}(U_i)$, $i \in I$, telle que
$\forall i, j\in I$
$$
s_i|_{U_i \cap U_j} = s_j|_{U_i \cap U_j}
$$
il existe une unique section $s \in \mathcal{F}(U)$ telle que
$s_i = s|_{U_i}$ pour tout $i \in I$.
\item Un {\it morphisme de faisceaux d'ensembles} est simplement un
morphisme de préfaisceaux d'ensembles.
\item La catégorie des faisceaux d'ensembles sur $X$ est notée
$\Sh(X)$.
\end{enumerate}
\end{definition}

\begin{remark}
\label{remark-confusion}
Il subsiste toujours une certaine confusion quant à la
nécessité de dire quelque chose sur l'ensemble des sections d'un
faisceau au-dessus de l'ensemble vide $\emptyset \subset X$.
Cela est nécessaire, et nous l'avons déjà fait si vous lisez
correctement la définition. En effet, remarquez que l'ensemble vide est
recouvert par le recouvrement ouvert vide, si bien que la « famille
de sections $s_i$ » de la définition ci-dessus forme en réalité
un élément du produit vide, qui est l'objet final
de la catégorie dans laquelle le faisceau prend ses valeurs. Autrement dit,
si vous lisez correctement la définition, vous en déduisez automatiquement
que $\mathcal{F}(\emptyset) = \textit{un objet final}$,
ce qui, dans le cas d'un faisceau d'ensembles, est un singleton.
Si cet argument ne vous plaît pas, vous pouvez simplement exiger
que $\mathcal{F}(\emptyset) = \{*\}$.

\medskip\noindent
En particulier, cette condition garantira alors que, si
$U, V \subset X$ sont des ouverts {\it disjoints}, alors
$$
\mathcal{F}(U \cup V) = \mathcal{F}(U) \times \mathcal{F}(V).
$$
(En effet, le produit fibré au-dessus d'un objet final est un produit.)
\end{remark}

\begin{example}
\label{example-basic-continuous-maps}
Soient $X$, $Y$ des espaces topologiques.
Considérons la règle $\mathcal{F}$ qui associe à
l'ouvert $U \subset X$ l'ensemble
$$
\mathcal{F}(U) = \{ f : U \to Y \mid f \text{ est continue}\}
$$
muni des applications de restriction évidentes. Nous affirmons que
$\mathcal{F}$ est un faisceau. Pour le voir, supposons que
$U = \bigcup_{i\in I} U_i$ soit un recouvrement ouvert, et que
l'on se donne $f_i \in \mathcal{F}(U_i)$, $i\in I$, avec
$f_i |_{U_i \cap U_j} = f_j|_{U_i \cap U_j}$ pour tous $i, j \in I$.
Dans ce cas, définissons $f : U \to Y$ en prenant $f(u)$
égal à la valeur de $f_i(u)$ pour n'importe quel $i \in I$ tel que
$u \in U_i$. Cette application est bien définie par hypothèse. De plus,
$f : U \to Y$ est une application dont la restriction à $U_i$
coïncide avec l'application continue $f_i$. Par conséquent, $f$ est
manifestement continue !
\end{example}

\noindent
Nous pouvons utiliser le résultat de l'exemple pour définir les
faisceaux constants. Soit donc $A$ un ensemble. Munissons $A$ de
la topologie discrète. Soit $U \subset X$ un ouvert.
Nous avons alors
$$
\{ f : U \to A \mid f\text{ continue}\}
=
\{ f : U \to A \mid f\text{ localement constante}\}.
$$
Ainsi, la règle qui associe à un ouvert l'ensemble des applications
localement constantes à valeurs dans $A$ est un faisceau.

\begin{definition}
\label{definition-constant-sheaf}
Soit $X$ un espace topologique. Soit $A$ un ensemble.
Le {\it faisceau constant de valeur $A$}, noté $\underline{A}$ ou
$\underline{A}_X$, est le faisceau qui associe à un ouvert $U \subset X$
l'ensemble de toutes les applications localement constantes $U \to A$, les applications
de restriction étant données par les restrictions de fonctions.
\end{definition}

\begin{example}
\label{example-sheaf-product-pointwise}
Soit $X$ un espace topologique. Soit $(A_x)_{x \in X}$
une famille d'ensembles $A_x$ indexée par les points $x \in X$. Nous
allons construire à partir de ces données un faisceau d'ensembles $\Pi$.
Pour tout ouvert $U \subset X$,
$$
\Pi(U) = \prod\nolimits_{x \in U} A_x.
$$
Pour tous ouverts $V \subset U \subset X$, définissons
une application de restriction par la règle suivante : un élément
$s = (a_x)_{x\in U} \in \Pi(U)$ a pour restriction
$s|_V = (a_x)_{x \in V}$. Il est clair que cela
définit un préfaisceau d'ensembles. Nous affirmons que ce préfaisceau est un faisceau.
Soit en effet $U = \bigcup U_i$ un recouvrement ouvert.
Supposons que les $s_i \in \Pi(U_i)$
soient tels que $s_i$ et $s_j$ coïncident sur $U_i \cap U_j$. Écrivons
$s_i = (a_{i, x})_{x\in U_i}$. La condition de compatibilité implique que
$a_{i, x} = a_{j, x}$ dans l'ensemble $A_x$ lorsque $x \in U_i \cap U_j$.
Il existe donc un unique élément $s = (a_x)_{x\in U}$
de $\Pi(U) = \prod_{x\in U} A_x$ tel que
$a_x = a_{i, x}$ dès que $x \in U_i$ pour un certain $i$. Bien entendu, cet
élément $s$ vérifie $s|_{U_i} = s_i$ pour tout $i$.
\end{example}

\begin{example}
\label{example-direct-sum-points-not-sheaf}
Soit $X$ un espace topologique.
Supposons que, pour tout $x\in X$, un groupe abélien $M_x$ soit donné.
Considérons le préfaisceau $\mathcal{F} : U \mapsto \bigoplus_{x \in U} M_x$
défini dans l'exemple \ref{example-direct-sum-points}. Ce
n'est pas un faisceau en général. Par exemple, si $X$ est un ensemble infini
muni de la topologie discrète, alors la condition de faisceau
impliquerait que $\mathcal{F}(X) = \prod_{x\in X} \mathcal{F}(\{x\})$,
mais, par définition, nous avons $\mathcal{F}(X)
= \bigoplus_{x \in X} M_x = \bigoplus_{x \in X} \mathcal{F}(\{x\})$.
Or une somme directe infinie diffère en général d'un
produit direct infini.

\medskip\noindent
Cependant, si $X$ est un espace topologique tel que tout ouvert
de $X$ soit quasi-compact, alors $\mathcal{F}$ {\it est} un faisceau.
Ceci est laissé en exercice au lecteur.
\end{example}

\section{Faisceaux abéliens}
\label{section-abelian-sheaves}

\begin{definition}
\label{definition-abelian-sheaf}
Soit $X$ un espace topologique.
\begin{enumerate}
\item Un {\it faisceau abélien sur $X$} ou un
{\it faisceau de groupes abéliens sur $X$}
est un préfaisceau abélien sur $X$ tel que le préfaisceau
d'ensembles sous-jacent soit un faisceau.
\item La catégorie des faisceaux de groupes abéliens
est notée $\textit{Ab}(X)$.
\end{enumerate}
\end{definition}

\noindent
Soit $X$ un espace topologique.
Dans le cas d'un préfaisceau abélien $\mathcal{F}$, la condition de
faisceau relative à un recouvrement ouvert $U = \bigcup U_i$
s'exprime souvent en disant que le complexe de groupes abéliens
$$
0 \to \mathcal{F}(U)
\to \prod\nolimits_i \mathcal{F}(U_i)
\to \prod\nolimits_{(i_0, i_1)} \mathcal{F}(U_{i_0} \cap U_{i_1})
$$
est exact. La première application est l'application usuelle, tandis que la seconde
envoie l'élément $(s_i)_{i \in I}$ sur l'élément
$$
(
s_{i_0}|_{U_{i_0} \cap U_{i_1}} -
s_{i_1}|_{U_{i_0} \cap U_{i_1}}
)_{(i_0, i_1)}
\in \prod\nolimits_{(i_0, i_1)} \mathcal{F}(U_{i_0} \cap U_{i_1})
$$

\section{Faisceaux de structures algébriques}
\label{section-sheaves-structures}

\noindent
Précisons la définition des faisceaux de certaines espèces de structures.
Commençons par reformuler la condition de faisceau. Supposons donc que
$\mathcal{F}$ soit un préfaisceau d'ensembles sur l'espace topologique $X$.
La condition de faisceau peut se reformuler comme suit. Soit
$U = \bigcup_{i\in I} U_i$ un recouvrement ouvert. Considérons le
diagramme
$$
\xymatrix{
\mathcal{F}(U) \ar[r]
&
\prod\nolimits_{i\in I}
\mathcal{F}(U_i)
\ar@<1ex>[r] \ar@<-1ex>[r]
&
\prod\nolimits_{(i_0, i_1) \in I \times I}
\mathcal{F}(U_{i_0} \cap U_{i_1})
}
$$
Ici, l'application de gauche est définie par la règle
$s \mapsto \prod_{i \in I} s|_{U_i}$. Les deux applications
de droite sont les applications
$$
\prod\nolimits_i s_i
\mapsto
\prod\nolimits_{(i_0, i_1)} s_{i_0}|_{U_{i_0} \cap U_{i_1}}
\text{ resp. }
\prod\nolimits_i s_i
\mapsto
\prod\nolimits_{(i_0, i_1)} s_{i_1}|_{U_{i_0} \cap U_{i_1}}.
$$
La condition de faisceau signifie exactement que la flèche de gauche
est l'égalisateur des deux flèches de droite. Cette reformulation s'étend
immédiatement aux préfaisceaux à valeurs dans une
catégorie, pourvu que celle-ci admette des produits.

\begin{definition}
\label{definition-sheaf-values-in-category}
Soit $X$ un espace topologique. Soit $\mathcal{C}$ une
catégorie admettant des produits. Un préfaisceau $\mathcal{F}$ à
valeurs dans $\mathcal{C}$ sur $X$ est un {\it faisceau}
si, pour tout recouvrement ouvert, le diagramme
$$
\xymatrix{
\mathcal{F}(U) \ar[r]
&
\prod\nolimits_{i\in I}
\mathcal{F}(U_i)
\ar@<1ex>[r] \ar@<-1ex>[r]
&
\prod\nolimits_{(i_0, i_1) \in I \times I}
\mathcal{F}(U_{i_0} \cap U_{i_1})
}
$$
est un diagramme égalisateur dans la catégorie $\mathcal{C}$.
\end{definition}

\noindent
Supposons que $\mathcal{C}$ soit une catégorie et que
$F : \mathcal{C} \to \textit{Sets}$ soit un foncteur fidèle.
Un bon exemple à garder à l'esprit est le cas où $\mathcal{C}$
est la catégorie des groupes abéliens et $F$ le foncteur d'oubli.
Considérons un préfaisceau $\mathcal{F}$ à valeurs dans $\mathcal{C}$ sur $X$.
Nous voudrions reformuler la condition précédente en termes
du préfaisceau d'ensembles sous-jacent
(Définition \ref{definition-underlying-presheaf-sets}).
Remarquons que le
préfaisceau d'ensembles sous-jacent est un faisceau d'ensembles si et seulement si tous les
diagrammes
$$
\xymatrix{
F(\mathcal{F}(U)) \ar[r]
&
\prod\nolimits_{i\in I}
F(\mathcal{F}(U_i))
\ar@<1ex>[r] \ar@<-1ex>[r]
&
\prod\nolimits_{(i_0, i_1) \in I \times I}
F(\mathcal{F}(U_{i_0} \cap U_{i_1}))
}
$$
d'ensembles -- après application du foncteur d'oubli $F$ -- sont des
diagrammes égalisateurs ! Nous voudrions donc que $\mathcal{C}$ admette
des produits et des égalisateurs et que $F$ commute avec
ces constructions. Cela équivaut à demander que $\mathcal{C}$
admette des limites et que $F$ commute avec elles ; voir
Catégories, Lemme \ref{categories-lemma-limits-products-equalizers}.
Mais cela ne suffit pas encore
(voir l'Exemple \ref{example-sheaves-topological-spaces}) ;
il faut aussi que $F$ {\it reflète les isomorphismes}.
Cette propriété signifie que, pour tout morphisme
$f : A \to A'$ de $\mathcal{C}$, $f$ est
un isomorphisme si (et seulement si) $F(f)$ est une bijection.

\begin{lemma}
\label{lemma-sheaves-structure}
Supposons que la catégorie $\mathcal{C}$ et
le foncteur $F : \mathcal{C} \to \textit{Sets}$
possèdent les propriétés suivantes :
\begin{enumerate}
\item $F$ est fidèle,
\item $\mathcal{C}$ admet des limites et $F$ commute avec elles, et
\item le foncteur $F$ reflète les isomorphismes.
\end{enumerate}
Soit $X$ un espace topologique. Soit $\mathcal{F}$
un préfaisceau à valeurs dans $\mathcal{C}$.
Alors $\mathcal{F}$ est un faisceau si et seulement si le
préfaisceau d'ensembles sous-jacent est un faisceau.
\end{lemma}

\begin{proof}
Supposons que $\mathcal{F}$ soit un faisceau. Alors
$\mathcal{F}(U)$ est l'égalisateur du diagramme
ci-dessus et, par hypothèse, $F(\mathcal{F}(U))$
est l'égalisateur du diagramme correspondant
d'ensembles. Ainsi, $F(\mathcal{F})$ est un faisceau d'ensembles.

\medskip\noindent
Supposons que $F(\mathcal{F})$ soit un faisceau.
Soit $E \in \Ob(\mathcal{C})$
l'égalisateur des deux flèches parallèles de la
Définition \ref{definition-sheaf-values-in-category}.
Nous obtenons un morphisme canonique $\mathcal{F}(U) \to E$,
simplement parce que $\mathcal{F}$ est un préfaisceau.
Par hypothèse, l'application induite $F(\mathcal{F}(U)) \to F(E)$
est un isomorphisme, puisque $F(E)$ est l'égalisateur
du diagramme correspondant d'ensembles. Nous voyons donc que
$\mathcal{F}(U) \to E$ est un isomorphisme
par la condition (3) du lemme.
\end{proof}

\noindent
Le lemme s'applique en particulier aux
{\it faisceaux de groupes, d'anneaux, d'algèbres sur un anneau fixé, de modules
sur un anneau fixé, d'espaces vectoriels sur un corps fixé, } etc.
Autrement dit, ce sont les préfaisceaux de groupes, d'anneaux,
de modules sur un anneau fixé, d'espaces vectoriels sur un corps fixé, etc.,
dont le préfaisceau d'ensembles sous-jacent est un faisceau.

\begin{example}
\label{example-C0-sheaf-rings}
Soit $X$ un espace topologique. Pour tout ouvert $U \subset X$, considérons
la $\mathbf{R}$-algèbre
$\mathcal{C}^{0}(U) = \{ f : U \to \mathbf{R} \mid f\text{ est continue}\}$.
Les applications de restriction évidentes en font un
préfaisceau de $\mathbf{R}$-algèbres sur $X$.
D'après l'Exemple \ref{example-basic-continuous-maps}, c'est un faisceau d'ensembles.
Le Lemme \ref{lemma-sheaves-structure} montre donc que c'est un faisceau de
$\mathbf{R}$-algèbres sur $X$.
\end{example}

\begin{example}
\label{example-sheaves-topological-spaces}
Considérons la catégorie des espaces topologiques $\textit{Top}$.
Il existe un foncteur fidèle naturel $\textit{Top} \to \textit{Sets}$
qui commute aux produits et aux égalisateurs. Mais il ne
reflète pas les isomorphismes. En fait, il s'avère que
l'analogue du Lemme \ref{lemma-sheaves-structure} est faux.
En effet, supposons $X = \mathbf{N}$ muni de la topologie
discrète. Donnons-nous des espaces topologiques discrets $A_i$, pour
$i \in \mathbf{N}$. Pour toute partie $U \subset \mathbf{N}$,
définissons $\mathcal{F}(U) = \prod_{i\in U} A_i$, muni de la
topologie discrète. On obtient alors un préfaisceau d'espaces
topologiques dont le préfaisceau d'ensembles sous-jacent est un faisceau ; voir
l'Exemple \ref{example-sheaf-product-pointwise}.
Cependant, si chaque $A_i$ possède au moins deux éléments, ce
préfaisceau n'est pas un faisceau d'espaces topologiques
au sens de la Définition \ref{definition-sheaf-values-in-category}.
Le lecteur pourra vérifier qu'en munissant chaque
$\mathcal{F}(U) = \prod_{i\in U} A_i$ de la {\it topologie produit}, on
obtient bien un faisceau
d'espaces topologiques sur $X$.
\end{example}


\section{Faisceaux de modules}
\label{section-sheaves-modules}

\begin{definition}
\label{definition-sheaf-modules}
Soit $X$ un espace topologique.
Soit $\mathcal{O}$ un faisceau d'anneaux sur $X$.
\begin{enumerate}
\item Un {\it faisceau de $\mathcal{O}$-modules} est un préfaisceau
de $\mathcal{O}$-modules $\mathcal{F}$,
voir la Définition \ref{definition-presheaf-modules},
tel que le préfaisceau de groupes abéliens sous-jacent $\mathcal{F}$
soit un faisceau.
\item Un {\it morphisme de faisceaux de $\mathcal{O}$-modules}
est un morphisme de préfaisceaux de $\mathcal{O}$-modules.
\item Étant donnés des faisceaux de $\mathcal{O}$-modules
$\mathcal{F}$ et $\mathcal{G}$, nous notons
$\Hom_\mathcal{O}(\mathcal{F}, \mathcal{G})$
l'ensemble des morphismes de faisceaux de $\mathcal{O}$-modules.
\item La catégorie des faisceaux de $\mathcal{O}$-modules
est notée $\textit{Mod}(\mathcal{O})$.
\end{enumerate}
\end{definition}

\noindent
Cette définition garde un certain sens même si $\mathcal{O}$ n'est qu'un
préfaisceau d'anneaux, bien que nous ne connaissions aucun exemple où
cela soit utile ; nous éviterons d'employer la terminologie
« faisceaux de $\mathcal{O}$-modules »
lorsque $\mathcal{O}$ n'est pas un faisceau d'anneaux.






\section{Fibres}
\label{section-stalks}

\noindent
Soit $X$ un espace topologique. Soit $x \in X$ un point.
Soit $\mathcal{F}$ un préfaisceau d'ensembles sur $X$.
La {\it fibre de $\mathcal{F}$ en $x$} est l'ensemble
$$
\mathcal{F}_x
=
\colim_{x\in U} \mathcal{F}(U)
$$
où la limite inductive est prise sur l'ensemble des voisinages ouverts
$U$ de $x$ dans $X$. L'ensemble des voisinages ouverts est
partiellement ordonné par l'inclusion inverse :
nous posons $U \geq U' \Leftrightarrow U \subset U'$.
Les morphismes de transition du système sont
donnés par les applications de restriction de $\mathcal{F}$.
Voir Catégories, Section \ref{categories-section-posets-limits}
pour les notations et la terminologie relatives aux (co)limites de systèmes.
Remarquons qu'il s'agit d'une limite inductive filtrante.
Il est donc facile de décrire $\mathcal{F}_x$. Plus précisément,
$$
\mathcal{F}_x
=
\{
(U, s)
\mid
x\in U, s\in \mathcal{F}(U)
\}/\sim
$$
où la relation d'équivalence est donnée par $(U, s) \sim (U', s')$ si et seulement s'il
existe un ouvert $U'' \subset U \cap U'$ tel que $x \in U''$ et
$s|_{U''} = s'|_{U''}$. Étant donné un couple $(U, s)$, nous notons parfois
$s_x$ l'élément de $\mathcal{F}_x$ correspondant à la classe d'équivalence
de $(U, s)$. Nous employons parfois l'expression
« image de $s$ dans $\mathcal{F}_x$ » pour désigner $s_x$.
Par exemple, étant donnés deux couples $(U, s)$ et $(U', s')$, nous disons parfois
« $s$ est égal à $s'$ dans $\mathcal{F}_x$ » pour indiquer
que $s_x = s'_x$. D'autres auteurs emploient la terminologie
« germe de $s$ en $x$ ».

\medskip\noindent
Une conséquence évidente de cette définition est que,
pour tout ouvert $U \subset X$, il existe une application canonique
$$
\mathcal{F}(U)
\longrightarrow
\prod\nolimits_{x \in U} \mathcal{F}_x
$$
définie par $s \mapsto \prod_{x \in U} (U, s)$. À méditer !

\begin{lemma}
\label{lemma-sheaf-subset-stalks}
Soit $\mathcal{F}$ un faisceau d'ensembles sur l'espace topologique $X$.
Pour tout ouvert $U \subset X$, l'application
$$
\mathcal{F}(U)
\longrightarrow
\prod\nolimits_{x \in U} \mathcal{F}_x
$$
est injective.
\end{lemma}

\begin{proof}
Supposons que $s, s' \in \mathcal{F}(U)$ aient la même image
dans chaque fibre $\mathcal{F}_x$ pour tout $x \in U$. Cela signifie que,
pour tout $x \in U$, il existe un ouvert $V^x \subset U$,
$x \in V^x$, tel que $s|_{V^x} = s'|_{V^x}$. Mais alors
$U = \bigcup_{x \in U} V^x$ est un recouvrement ouvert. L'unicité
dans la condition de faisceau entraîne donc que $s = s'$.
\end{proof}

\begin{definition}
\label{definition-separated}
Soit $X$ un espace topologique.
Un préfaisceau d'ensembles $\mathcal{F}$ sur $X$ est {\it séparé}
si, pour tout ouvert $U \subset X$, l'application
$\mathcal{F}(U) \to \prod_{x \in U} \mathcal{F}_x$ est
injective.
\end{definition}

\noindent
Observons aussi que la construction de la fibre
$\mathcal{F}_x$ est fonctorielle par rapport au préfaisceau $\mathcal{F}$.
Autrement dit, elle définit un foncteur
$$
\textit{PSh}(X) \longrightarrow \textit{Sets},
\ \mathcal{F} \longmapsto \mathcal{F}_x.
$$
Ce foncteur est appelé le {\it foncteur fibre}.
Plus précisément, si $\varphi : \mathcal{F} \to \mathcal{G}$ est
un morphisme de préfaisceaux, nous définissons
$\varphi_x : \mathcal{F}_x \to \mathcal{G}_x$
par la règle $(U, s) \mapsto (U, \varphi(s))$.
Pour vérifier que cette définition a un sens, il faut montrer que,
si $(U, s) = (U', s')$ dans $\mathcal{F}_x$, alors
$(U, \varphi(s)) = (U', \varphi(s'))$ dans $\mathcal{G}_x$.
C'est clair, puisque $\varphi$ est compatible avec les
applications de restriction.

\begin{example}
\label{example-stalk-constant-presheaf}
Soit $X$ un espace topologique. Soit $A$ un ensemble.
Notons provisoirement $A_p$ le préfaisceau constant
de valeur $A$ ($p$ pour préfaisceau -- non pour point).
Il existe un morphisme canonique de préfaisceaux
$A_p \to \underline{A}$ vers le faisceau constant de
valeur $A$. En tout point, nous avons des
bijections canoniques $A = (A_p)_x = \underline{A}_x$, où
la seconde application est induite, par fonctorialité, par
l'application $A_p \to \underline{A}$.
\end{example}



\begin{example}
\label{example-germs-functions}
Supposons que $X = \mathbf{R}^n$ soit muni de la topologie euclidienne.
Considérons le préfaisceau des fonctions $\mathcal{C}^\infty$
sur $X$, noté $\mathcal{C}^\infty_{\mathbf{R}^n}$.
Autrement dit, $\mathcal{C}^\infty_{\mathbf{R}^n}(U)$ est l'ensemble
des fonctions de classe $\mathcal{C}^\infty$, $f : U \to \mathbf{R}$.
Comme dans l'Exemple \ref{example-basic-continuous-maps},
il est facile de montrer que c'est un faisceau. C'est même
un faisceau de $\mathbf{R}$-espaces vectoriels.

\medskip\noindent
Soit ensuite $x \in X = \mathbf{R}^n$ un point. Comment
se représenter un élément de la fibre $\mathcal{C}^\infty_{\mathbf{R}^n, x}$ ?
Un tel élément est représenté par une fonction $\mathcal{C}^\infty$
$f$ dont le domaine contient $x$. Deux telles
fonctions $f$, $g$ déterminent
le même élément de la fibre si elles coïncident sur un voisinage
de $x$. Autrement dit, un élément de $\mathcal{C}^\infty_{\mathbf{R}^n, x}$
n'est autre que ce qu'on appelle parfois
un {\it germe de fonction $\mathcal{C}^\infty$ en $x$}.
\end{example}

\begin{example}
\label{example-sheaf-product-pointwise-stalk}
Soit $X$ un espace topologique. Soit $A_x$ un ensemble pour tout $x \in X$.
Considérons le faisceau $\mathcal{F} : U \mapsto \prod_{x\in U} A_x$ de l'Exemple
\ref{example-sheaf-product-pointwise}. Nous voulons simplement signaler
ici que la fibre $\mathcal{F}_x$ de $\mathcal{F}$ en $x$
n'est en général {\it pas} égale à l'ensemble $A_x$. Il existe bien entendu
une application $\mathcal{F}_x \to A_x$, mais on ne peut en général rien dire
de plus. Supposons par exemple que $x = \lim x_n$ avec $x_n \not = x_m$ pour
tous $n \not = m$ et que $A_y = \{0, 1\}$ pour tout $y \in X$. Alors
$\mathcal{F}_x$ admet une surjection sur l'ensemble (infini) des queues de suites
de $0$ et de $1$. En effet, tout voisinage ouvert de $x$ contient presque
tous les $x_n$. En revanche, si tout voisinage de $x$ contient
un point $y$ tel que $A_y = \emptyset$, alors $\mathcal{F}_x = \emptyset$.
\end{example}



\section{Fibres des préfaisceaux abéliens}
\label{section-stalks-abelian-presheaves}

\noindent
Nous traitons d'abord le cas des groupes abéliens, qui
servira de modèle au cas général.

\begin{lemma}
\label{lemma-stalk-abelian-presheaf}
Soit $X$ un espace topologique. Soit $\mathcal{F}$ un préfaisceau
de groupes abéliens sur $X$. Il existe une unique structure de
groupe abélien sur $\mathcal{F}_x$ telle que, pour tout
$U \subset X$ ouvert avec $x\in U$, l'application $\mathcal{F}(U) \to \mathcal{F}_x$
soit un homomorphisme de groupes. De plus,
$$
\mathcal{F}_x
=
\colim_{x\in U} \mathcal{F}(U)
$$
dans la catégorie des groupes abéliens.
\end{lemma}

\begin{proof}
Nous définissons la somme des deux éléments
$(U, s)$ et $(V, t)$ comme le couple $(U \cap V, s|_{U\cap V} +
t|_{U \cap V})$. Le reste se vérifie aisément.
\end{proof}

\noindent
Le point crucial de la démonstration précédente est que l'ensemble partiellement ordonné des
voisinages ouverts est un ensemble dirigé (Catégories, Définition
\ref{categories-definition-directed-set}). En effet, le coproduit
de deux groupes abéliens $A, B$ est la somme directe $A \oplus B$, tandis que
le coproduit dans la catégorie des ensembles est la réunion
disjointe $A \amalg B$, ce qui montre que les limites inductives dans la catégorie
des groupes abéliens ne coïncident pas en général avec les limites inductives dans la
catégorie des ensembles.


\section{Fibres des préfaisceaux de structures algébriques}
\label{section-stalks-presheaves-structures}

\noindent
La démonstration du Lemme \ref{lemma-stalk-abelian-presheaf} vaut
pour toute espèce de structure algébrique pour laquelle le foncteur
d'oubli commute aux limites inductives filtrantes.

\begin{lemma}
\label{lemma-stalk-presheaf-values-in-category}
Soit $\mathcal{C}$ une catégorie. Soit $F : \mathcal{C} \to \textit{Sets}$
un foncteur. Supposons que
\begin{enumerate}
\item $F$ soit fidèle, et
\item les limites inductives filtrantes existent dans $\mathcal{C}$ et que $F$ commute avec
elles.
\end{enumerate}
Soit $X$ un espace topologique. Soit $x \in X$. Soit $\mathcal{F}$
un préfaisceau à valeurs dans $\mathcal{C}$.
Alors
$$
\mathcal{F}_x = \colim_{x\in U} \mathcal{F}(U)
$$
existe dans $\mathcal{C}$. Son ensemble sous-jacent est égal à la
fibre du préfaisceau d'ensembles sous-jacent à $\mathcal{F}$.
En outre, la construction $\mathcal{F} \mapsto \mathcal{F}_x$
est un foncteur de la catégorie des préfaisceaux à valeurs dans
$\mathcal{C}$ vers $\mathcal{C}$.
\end{lemma}

\begin{proof}
Omis.
\end{proof}

\noindent
Par définition même, tous les morphismes $\mathcal{F}(U)
\to \mathcal{F}_x$ sont des morphismes de la catégorie $\mathcal{C}$
qui, après application du foncteur d'oubli $F$, deviennent
les applications correspondantes du faisceau d'ensembles sous-jacent.
Comme d'habitude, nous ne distinguerons pas le morphisme
de $\mathcal{C}$ de l'application d'ensembles sous-jacente, ce qui
est permis puisque $F$ est fidèle.

\medskip\noindent
Ce lemme s'applique en particulier aux
{\it préfaisceaux de groupes (non nécessairement abéliens), d'anneaux, de modules
sur un anneau fixé et d'espaces vectoriels sur un corps fixé}.

\section{Fibres des préfaisceaux de modules}
\label{section-stalk-presheaves-modules}

\begin{lemma}
\label{lemma-stalk-module}
Soit $X$ un espace topologique.
Soit $\mathcal{O}$ un préfaisceau d'anneaux sur $X$.
Soit $\mathcal{F}$ un préfaisceau de $\mathcal{O}$-modules.
Soit $x \in X$.
L'application canonique $\mathcal{O}_x \times \mathcal{F}_x
\to \mathcal{F}_x$ provenant de l'application de multiplication
$\mathcal{O} \times \mathcal{F} \to \mathcal{F}$ définit
une structure de $\mathcal{O}_x$-module sur le groupe abélien
$\mathcal{F}_x$.
\end{lemma}

\begin{proof}
Omis.
\end{proof}

\begin{lemma}
\label{lemma-stalk-tensor-presheaf-modules}
Soit $X$ un espace topologique.
Soit $\mathcal{O} \to \mathcal{O}'$ un morphisme de
préfaisceaux d'anneaux sur $X$.
Soit $\mathcal{F}$ un préfaisceau de $\mathcal{O}$-modules.
Soit $x \in X$. Nous avons
$$
\mathcal{F}_x \otimes_{\mathcal{O}_x} \mathcal{O}'_x
=
(\mathcal{F} \otimes_{p, \mathcal{O}} \mathcal{O}')_x
$$
en tant que $\mathcal{O}'_x$-modules.
\end{lemma}

\begin{proof}
Omis.
\end{proof}
\section{Structures algébriques}
\label{section-algebraic-structures}

\noindent
Dans cette section, nous formalisons légèrement les notions rencontrées dans
les sections précédentes.

\begin{definition}
\label{definition-algebraic-structure}
Une {\it espèce de structure algébrique} est la donnée d'une catégorie
$\mathcal{C}$ et d'un foncteur $F : \mathcal{C} \to \textit{Sets}$ ayant les
propriétés suivantes :
\begin{enumerate}
\item $F$ est fidèle,
\item $\mathcal{C}$ admet des limites et $F$ commute aux limites,
\item $\mathcal{C}$ admet des limites inductives filtrantes et $F$ commute avec
celles-ci, et
\item $F$ reflète les isomorphismes.
\end{enumerate}
\end{definition}

\noindent
Nous introduisons cette définition afin de mettre en évidence les propriétés
que nous emploierons dans plusieurs arguments ci-dessous. Nous n'étudierons
cependant pas cette notion en détail, car nos conventions nous interdisent
d'étudier les catégories « grandes », à l'exception de celles énumérées dans
Catégories, remarque \ref{categories-remark-big-categories}.
Parmi celles-ci, les catégories suivantes possèdent les propriétés requises.

\begin{lemma}
\label{lemma-list-algebraic-structures}
Les catégories suivantes, munies du foncteur d'oubli évident, définissent des
espèces de structures algébriques :
\begin{enumerate}
\item La catégorie des ensembles pointés.
\item La catégorie des groupes abéliens.
\item La catégorie des groupes.
\item La catégorie des monoïdes.
\item La catégorie des anneaux.
\item La catégorie des $R$-modules, pour un anneau fixé $R$.
\item La catégorie des algèbres de Lie sur un corps fixé.
\end{enumerate}
\end{lemma}

\begin{proof}
Omis.
\end{proof}

\noindent
Désormais, nous considérerons les (pré)faisceaux de structures algébriques et
leurs fibres au moyen des (pré)faisceaux d'ensembles sous-jacents. Cela est
permis par les lemmes \ref{lemma-sheaves-structure} et
\ref{lemma-stalk-presheaf-values-in-category}.

\medskip\noindent
Dans le reste de cette section, nous signalons quelques résultats sur les
structures algébriques qui seront utiles par la suite.

\begin{lemma}
\label{lemma-properties-algebraic-structures}
Soit $(\mathcal{C}, F)$ une espèce de structure algébrique.
\begin{enumerate}
\item $\mathcal{C}$ admet un objet final $0$ et $F(0) = \{ * \}$.
\item $\mathcal{C}$ admet des produits et $F(\prod A_i) = \prod F(A_i)$.
\item $\mathcal{C}$ admet des produits fibrés et
$F(A \times_B C) = F(A)\times_{F(B)}F(C)$.
\item $\mathcal{C}$ admet des égalisateurs et, si $E \to A$ est
l'égalisateur de $a, b : A \to B$, alors $F(E) \to F(A)$ est l'égalisateur
de $F(a), F(b) : F(A) \to F(B)$.
\item $A \to B$ est un monomorphisme si et seulement si
$F(A) \to F(B)$ est injective.
\item si $F(a) : F(A) \to F(B)$ est surjective, alors $a$ est un
épimorphisme.
\item étant donnée une suite $A_1 \to A_2 \to A_3 \to \ldots$, la limite
inductive $\colim A_i$ existe et $F(\colim A_i) = \colim F(A_i)$ ; plus
généralement, il en va de même pour toute limite inductive filtrante.
\end{enumerate}
\end{lemma}

\begin{proof}
Omis. La seule assertion qui mérite commentaire est (5). Elle résulte du fait
que $A \to B$ est un monomorphisme si et seulement si
$A \to A \times_B A$ est un isomorphisme, puis de ce que $F$ reflète les
isomorphismes.
\end{proof}

\begin{lemma}
\label{lemma-image-contained-in}
Soit $(\mathcal{C}, F)$ une espèce de structure algébrique.
Supposons que $A, B, C \in \Ob(\mathcal{C})$.
Soient $f : A \to B$ et $g : C \to B$ des morphismes de $\mathcal{C}$.
Si $F(g)$ est injective et si $\Im(F(f)) \subset \Im(F(g))$, alors $f$ se
factorise sous la forme $f = g \circ t$ pour un certain morphisme
$t : A \to C$.
\end{lemma}

\begin{proof}
Considérons $A \times_B C$. Les hypothèses impliquent que
$F(A \times_B C) = F(A) \times_{F(B)} F(C) = F(A)$.
Ainsi, $A = A \times_B C$, puisque $F$ reflète les isomorphismes.
Le résultat s'ensuit.
\end{proof}

\begin{example}
\label{example-application-lemma-image-contained-in}
Nous appliquerons souvent le lemme dans la situation suivante.
Supposons donné un diagramme
$$
\xymatrix{
A \ar[r] & B \ar[d] \\
C \ar[r] & D
}
$$
dans $\mathcal{C}$. Supposons que $C \to D$ soit injective sur les ensembles
sous-jacents et que l'image, sur les ensembles sous-jacents, de la composée
$A \to B \to D$ soit contenue dans l'image de $C \to D$.
On obtient alors un diagramme commutatif
$$
\xymatrix{
A \ar[r] \ar[d] & B \ar[d] \\
C \ar[r] & D
}
$$
dans $\mathcal{C}$.
\end{example}

\begin{example}
\label{example-sheaf-product-pointwise-algebraic-structure}
Soit $F : \mathcal{C} \to \textit{Sets}$ une espèce de structure algébrique.
Soit $X$ un espace topologique. Supposons que, pour tout $x \in X$, un objet
$A_x \in \Ob(\mathcal{C})$ soit donné. Considérons le préfaisceau $\Pi$ à
valeurs dans $\mathcal{C}$ sur $X$, défini par la règle
$\Pi(U) = \prod_{x \in U} A_x$ (avec les applications de restriction
évidentes). Remarquons que le préfaisceau d'ensembles associé
$U \mapsto F(\Pi(U)) = \prod_{x \in U} F(A_x)$ est un faisceau d'après
l'exemple \ref{example-sheaf-product-pointwise}.
Ainsi, $\Pi$ est un faisceau de structures algébriques de type
$(\mathcal{C} , F)$. On obtient de cette manière de nombreux exemples de
faisceaux de groupes abéliens, de groupes, d'anneaux, etc.
\end{example}












\section{Exactitude et points}
\label{section-exactness-points}

\noindent
Dans toute catégorie, on dispose des notions d'épimorphisme, de
monomorphisme, d'isomorphisme, etc.

\begin{lemma}
\label{lemma-points-exactness}
Soit $X$ un espace topologique. Soit $\varphi : \mathcal{F} \to \mathcal{G}$
un morphisme de faisceaux d'ensembles sur $X$.
\begin{enumerate}
\item L'application $\varphi$ est un monomorphisme dans la catégorie des
faisceaux si et seulement si, pour tout $x \in X$, l'application
$\varphi_x : \mathcal{F}_x \to \mathcal{G}_x$ est injective.
\item L'application $\varphi$ est un épimorphisme dans la catégorie des
faisceaux si et seulement si, pour tout $x \in X$, l'application
$\varphi_x : \mathcal{F}_x \to \mathcal{G}_x$ est surjective.
\item L'application $\varphi$ est un isomorphisme dans la catégorie des
faisceaux si et seulement si, pour tout $x \in X$, l'application
$\varphi_x : \mathcal{F}_x \to \mathcal{G}_x$ est bijective.
\end{enumerate}
\end{lemma}

\begin{proof}
Omis.
\end{proof}

\noindent
Il en résulte que, dans la catégorie des faisceaux d'ensembles, les notions
d'épimorphisme et de monomorphisme se décrivent comme suit.

\begin{definition}
\label{definition-injective-surjective}
Soit $X$ un espace topologique.
\begin{enumerate}
\item Un préfaisceau $\mathcal{F}$ est appelé un {\it sous-préfaisceau} d'un
préfaisceau $\mathcal{G}$ si
$\mathcal{F}(U) \subset \mathcal{G}(U)$ pour tout ouvert $U \subset X$ et
si les applications de restriction de $\mathcal{G}$ induisent celles de
$\mathcal{F}$. Si $\mathcal{F}$ et $\mathcal{G}$ sont des faisceaux, on dit
que $\mathcal{F}$ est un {\it sous-faisceau} de $\mathcal{G}$. Nous
l'indiquons parfois par la notation $\mathcal{F} \subset \mathcal{G}$.
\item Un morphisme de préfaisceaux d'ensembles
$\varphi : \mathcal{F} \to \mathcal{G}$ sur $X$ est dit {\it injectif} si
et seulement si $\mathcal{F}(U) \to \mathcal{G}(U)$ est injective pour tout
ouvert $U$ de $X$.
\item Un morphisme de préfaisceaux d'ensembles
$\varphi : \mathcal{F} \to \mathcal{G}$ sur $X$ est dit {\it surjectif} si
et seulement si $\mathcal{F}(U) \to \mathcal{G}(U)$ est surjective pour tout
ouvert $U$ de $X$.
\item Un morphisme de faisceaux d'ensembles
$\varphi : \mathcal{F} \to \mathcal{G}$ sur $X$ est dit {\it injectif} si
et seulement si $\mathcal{F}(U) \to \mathcal{G}(U)$ est injective pour tout
ouvert $U$ de $X$.
\item Un morphisme de faisceaux d'ensembles
$\varphi : \mathcal{F} \to \mathcal{G}$ sur $X$ est dit {\it surjectif} si
et seulement si, pour tout ouvert $U$ de $X$ et toute section $s$ de
$\mathcal{G}(U)$, il existe un recouvrement ouvert
$U = \bigcup U_i$ tel que $s|_{U_i}$ appartienne à l'image de
$\mathcal{F}(U_i) \to \mathcal{G}(U_i)$ pour tout $i$.
\end{enumerate}
\end{definition}


\begin{lemma}
\label{lemma-characterize-epi-mono}
Soit $X$ un espace topologique.
\begin{enumerate}
\item Les épimorphismes (resp.\ les monomorphismes) dans la catégorie des
préfaisceaux sont exactement les morphismes surjectifs (resp.\ injectifs) de
préfaisceaux.
\item Les épimorphismes (resp.\ les monomorphismes) dans la catégorie des
faisceaux sont exactement les morphismes surjectifs (resp.\ injectifs) de
faisceaux ; ce sont également exactement les morphismes surjectifs
(resp.\ injectifs) sur toutes les fibres.
\end{enumerate}
\end{lemma}

\begin{proof}
Omis.
\end{proof}


\begin{lemma}
\label{lemma-check-homomorphism-stalks}
Soit $X$ un espace topologique.
Soit $(\mathcal{C}, F)$ une espèce de structure algébrique.
Supposons que $\mathcal{F}$ et $\mathcal{G}$ soient des faisceaux sur $X$ à
valeurs dans $\mathcal{C}$.
Soit $\varphi : \mathcal{F} \to \mathcal{G}$ une application entre les
faisceaux d'ensembles sous-jacents.
Si, pour tout point $x \in X$, l'application
$\mathcal{F}_x \to \mathcal{G}_x$ est un morphisme de structures
algébriques, alors $\varphi$ est un morphisme de faisceaux de structures
algébriques.
\end{lemma}

\begin{proof}
Soit $U$ un ouvert de $X$. Considérons le diagramme d'ensembles sous-jacents
$$
\xymatrix{
\mathcal{F}(U) \ar[r] \ar[d] &
\prod_{x \in U} \mathcal{F}_x \ar[d] \\
\mathcal{G}(U) \ar[r] &
\prod_{x \in U} \mathcal{G}_x
}
$$
Par hypothèse et d'après les résultats précédents, toutes les flèches sauf la
flèche verticale de gauche sont des morphismes de structures algébriques. De
plus, la flèche horizontale inférieure est injective ; voir le
lemme \ref{lemma-sheaf-subset-stalks}.
La conclusion résulte donc du lemme \ref{lemma-image-contained-in} ; voir
également l'exemple \ref{example-application-lemma-image-contained-in}.
\end{proof}

\noindent
Les suites exactes courtes de faisceaux abéliens, entre autres, seront étudiées
dans le chapitre consacré aux faisceaux de modules. Voir
Modules, section \ref{modules-section-kernels}.



\section{Faisceau associé}
\label{section-sheafification}

\noindent
Dans cette section, nous expliquons comment construire le faisceau associé
à un préfaisceau sur un espace topologique. Nous utiliserons les fibres
pour décrire ici le faisceau associé. Cette construction diffère
du procédé général décrit dans Sites, Section
\ref{sites-section-sheafification}, et elle est peut-être un peu
plus facile à comprendre.

\medskip\noindent
La construction fondamentale est la suivante. Soit $\mathcal{F}$ un préfaisceau
d'ensembles sur un espace topologique $X$.
Pour tout ouvert $U \subset X$, nous posons
$$
\mathcal{F}^{\#}(U)
=
\{
(s_u) \in \prod\nolimits_{u \in U} \mathcal{F}_u
\text{ tels que }(*)
\}
$$
où $(*)$ désigne la propriété suivante :
\begin{itemize}
\item[(*)] Pour tout $u \in U$, il existe un voisinage ouvert
$u \in V \subset U$ et une section $\sigma \in \mathcal{F}(V)$
tels que, pour tout $v \in V$, on ait $s_v = (V, \sigma)$
dans $\mathcal{F}_v$.
\end{itemize}
Remarquons que $(*)$ est une condition imposée en chaque $u \in U$
et que, $u \in U$ étant fixé, sa validité
ne dépend que des valeurs $s_v$ pour $v$ dans un voisinage ouvert quelconque
de $u$. Il est donc clair que,
si $V \subset U \subset X$ sont ouverts, les applications de projection
$$
\prod\nolimits_{u \in U} \mathcal{F}_u
\longrightarrow
\prod\nolimits_{v \in V} \mathcal{F}_v
$$
envoient les éléments de $\mathcal{F}^{\#}(U)$ dans $\mathcal{F}^{\#}(V)$.
En prenant ces applications pour applications de restriction, on munit $\mathcal{F}^\#$
d'une structure de préfaisceau d'ensembles sur $X$.

\medskip\noindent
De plus, l'application $\mathcal{F}(U) \to \prod_{u \in U} \mathcal{F}_u$
décrite dans la Section \ref{section-stalks} est manifestement à valeurs
dans $\mathcal{F}^{\#}(U)$. Par ailleurs, si $V \subset U \subset X$ sont
ouverts, on a le diagramme commutatif suivant :
$$
\xymatrix{
\mathcal{F}(U) \ar[r] \ar[d] &
\mathcal{F}^{\#}(U) \ar[r] \ar[d] &
\prod_{u\in U} \mathcal{F}_u \ar[d] \\
\mathcal{F}(V) \ar[r] &
\mathcal{F}^{\#}(V) \ar[r] &
\prod_{v\in V} \mathcal{F}_v
}
$$
où les flèches verticales sont induites par les
applications de restriction. On voit donc qu'il existe
un morphisme canonique de préfaisceaux
$\mathcal{F} \to \mathcal{F}^{\#}$.

\medskip\noindent
Dans l'Exemple \ref{example-sheaf-product-pointwise}, nous avons vu
que la règle $\Pi(\mathcal{F}) : U \mapsto \prod_{u\in U} \mathcal{F}_u$
définit un faisceau, muni des applications de restriction évidentes. Par construction,
$\mathcal{F}^{\#}$ en est un sous-préfaisceau. Autrement dit, on
a des morphismes de préfaisceaux
$$
\mathcal{F} \to \mathcal{F}^\# \to \Pi(\mathcal{F}).
$$
En outre, la règle qui associe à $\mathcal{F}$
la suite ci-dessus est manifestement fonctorielle en le préfaisceau $\mathcal{F}$.
Cette notation sera
utilisée dans les démonstrations des lemmes qui suivent.

\begin{lemma}
\label{lemma-sheafification-sheaf}
Le préfaisceau $\mathcal{F}^{\#}$ est un faisceau.
\end{lemma}

\begin{proof}
Il vaut sans doute mieux que le lecteur trouve lui-même une explication
plutôt que de lire la démonstration ci-dessous. En fait, le lemme est vrai
pour la même raison que le préfaisceau des fonctions
continues est un faisceau, voir l'Exemple \ref{example-basic-continuous-maps}
(et cette analogie peut être précisée au moyen de l'« espace \'etal\'e »).

\medskip\noindent
Quoi qu'il en soit, soit $U = \bigcup U_i$ un recouvrement ouvert.
Supposons donnés des $s_i = (s_{i, u})_{u \in U_i} \in \mathcal{F}^{\#}(U_i)$
tels que $s_i$ et $s_j$ coïncident sur $U_i \cap U_j$.
Comme $\Pi(\mathcal{F})$ est un faisceau,
il existe un élément $s = (s_u)_{u\in U}$ de $\prod_{u\in U} \mathcal{F}_u$
qui se restreint en $s_i$ sur $U_i$. Il reste à vérifier la propriété $(*)$.
Soit $u \in U$. Alors $u \in U_i$ pour un certain $i$. La propriété $(*)$ appliquée à $s_i$
fournit donc un ouvert $V$, avec $u \in V \subset U_i$,
et un $\sigma \in \mathcal{F}(V)$
tels que $s_{i, v} = (V, \sigma)$ dans $\mathcal{F}_v$
pour tout $v \in V$. Puisque $s_{i, v} = s_v$, la propriété $(*)$ vaut pour $s$.
\end{proof}

\begin{lemma}
\label{lemma-stalk-sheafification}
Soit $X$ un espace topologique.
Soit $\mathcal{F}$ un préfaisceau d'ensembles sur $X$.
Soit $x \in X$. Alors $\mathcal{F}_x = \mathcal{F}^\#_x$.
\end{lemma}

\begin{proof}
L'application $\mathcal{F}_x \to \mathcal{F}^\#_x$
est injective, puisque l'application
$\mathcal{F}_x \to \Pi(\mathcal{F})_x$ l'est déjà.
En effet, il existe une application canonique $\Pi(\mathcal{F})_x \to \mathcal{F}_x$
qui est un inverse à gauche de l'application $\mathcal{F}_x \to \Pi(\mathcal{F})_x$,
voir l'Exemple \ref{example-sheaf-product-pointwise-stalk}.
Pour montrer qu'elle est surjective, soit
$\overline{s} \in \mathcal{F}^\#_x$.
On peut trouver un voisinage ouvert $U$ de $x$ tel que
$\overline{s}$ soit la classe d'équivalence de $(U, s)$,
où $s \in \mathcal{F}^\#(U)$.
Par définition, cela signifie qu'il existe un voisinage ouvert
$V \subset U$ de $x$ et une section $\sigma \in \mathcal{F}(V)$
tels que $s|_V$ soit l'image de $\sigma$ dans $\Pi(\mathcal{F})(V)$.
La classe de $(V, \sigma)$ définit manifestement un élément de
$\mathcal{F}_x$ qui s'envoie sur $\overline{s}$.
\end{proof}

\begin{lemma}
\label{lemma-sheafify-universal}
Soit $\mathcal{F}$ un préfaisceau d'ensembles sur $X$.
Tout morphisme $\mathcal{F} \to \mathcal{G}$ vers un faisceau d'ensembles
se factorise de manière unique sous la forme
$\mathcal{F} \to \mathcal{F}^\# \to \mathcal{G}$.
\end{lemma}

\begin{proof}
On a manifestement un diagramme commutatif
$$
\xymatrix{
\mathcal{F} \ar[r] \ar[d] &
\mathcal{F}^\# \ar[r] \ar[d] &
\Pi(\mathcal{F}) \ar[d] \\
\mathcal{G} \ar[r] &
\mathcal{G}^\# \ar[r] &
\Pi(\mathcal{G}) \\
}
$$
Il suffit donc de montrer que $\mathcal{G} = \mathcal{G}^\#$.
D'après le Lemme \ref{lemma-points-exactness}, il suffit pour cela de montrer,
pour tout point $x \in X$, que l'application
$\mathcal{G}_x \to \mathcal{G}^\#_x$ est bijective. C'est précisément
le Lemme \ref{lemma-stalk-sheafification} ci-dessus.
\end{proof}

\noindent
Ce lemme affirme en réalité que les foncteurs
$i : \Sh(X) \to \textit{PSh}(X)$
(inclusion) et $\# : \textit{PSh}(X) \to \Sh(X)$
(faisceau associé) forment un couple de foncteurs adjoints. La formule est
$$
\Mor_{\textit{PSh}(X)}(\mathcal{F}, i(\mathcal{G}))
=
\Mor_{\Sh(X)}(\mathcal{F}^\#, \mathcal{G})
$$
ce qui signifie que le foncteur faisceau associé est adjoint à gauche du
foncteur d'inclusion. Voir Categories, Section
\ref{categories-section-adjoint}.

\begin{example}
\label{example-sheafify-constant}
Pour les notations, voir l'Exemple \ref{example-stalk-constant-presheaf}.
L'application $A_p \to \underline{A}$ induit une application
$A_p^\# \to \underline{A}$. On voit aisément que celle-ci
est un isomorphisme. Autrement dit, le faisceau associé
au préfaisceau constant de valeur $A$ est le
faisceau constant de valeur $A$.
\end{example}

\begin{lemma}
\label{lemma-separated-presheaf-into-sheaf}
Soit $X$ un espace topologique.
Un préfaisceau $\mathcal{F}$ est séparé (voir la
Définition \ref{definition-separated}) si et seulement si
l'application canonique $\mathcal{F} \to \mathcal{F}^\#$ est injective.
\end{lemma}

\begin{proof}
Cela résulte immédiatement de la construction de $\mathcal{F}^\#$ donnée dans cette
section.
\end{proof}

\begin{lemma}
\label{lemma-sheafify-epi-mono}
Soit $X$ un espace topologique. Le morphisme de faisceaux associés à un morphisme surjectif
(resp.\ injectif) de préfaisceaux d'ensembles est surjectif
(resp.\ injectif).
\end{lemma}

\begin{proof}
Omis.
\end{proof}


\section{Faisceau associé à un préfaisceau abélien}
\label{section-sheafify-abelian-presheaves}

\noindent
Le lemme suivant, d'aspect quelque peu étrange, est vraisemblablement superflu, mais
très commode pour traiter le passage au faisceau associé des préfaisceaux
de structures algébriques.

\begin{lemma}
\label{lemma-diagram-fibre-product}
Soit $X$ un espace topologique. Soit $\mathcal{F}$
un préfaisceau d'ensembles sur $X$. Soit $U \subset X$ un ouvert.
On a un diagramme cartésien canonique
$$
\xymatrix{
\mathcal{F}^\#(U) \ar[d] \ar[r] &
\Pi(\mathcal{F})(U) \ar[d] \\
\prod_{x \in U} \mathcal{F}_x
\ar[r] &
\prod_{x \in U} \Pi(\mathcal{F})_x
}
$$
dont les flèches sont les suivantes :
\begin{enumerate}
\item Les composantes de la flèche verticale de gauche sont les applications
$\mathcal{F}^\#(U) \to \mathcal{F}^\#_x = \mathcal{F}_x$
où l'égalité résulte du Lemme \ref{lemma-stalk-sheafification}.
\item La flèche horizontale supérieure provient du
morphisme de préfaisceaux $\mathcal{F} \to \Pi(\mathcal{F})$ décrit
dans la Section \ref{section-sheafification}.
\item Les composantes évidentes de la flèche verticale de droite sont les
applications $\Pi(\mathcal{F})(U) \to \Pi(\mathcal{F})_x$.
\item Les composantes de la flèche horizontale inférieure sont les applications
$\mathcal{F}_x \to \Pi(\mathcal{F})_x$
qui proviennent du morphisme de préfaisceaux
$\mathcal{F} \to \Pi(\mathcal{F})$ décrit
dans la Section \ref{section-sheafification}.
\end{enumerate}
\end{lemma}

\begin{proof}
Le diagramme est manifestement commutatif. Il faut montrer
qu'il est cartésien. La flèche horizontale inférieure
est injective puisque toutes les applications $\mathcal{F}_x \to \Pi(\mathcal{F})_x$
sont injectives (voir le début de la démonstration du
Lemme \ref{lemma-stalk-sheafification}).
Une section $s \in \Pi(\mathcal{F})(U)$ appartient à $\mathcal{F}^\#$ si et
seulement si elle vérifie $(*)$. Or $(*)$ affirme qu'au voisinage de tout point
la section $s$ provient d'une section de $\mathcal{F}$. Par définition
des foncteurs fibres, cela équivaut à dire que
la valeur de $s$ dans toute fibre $\Pi(\mathcal{F})_x$ provient
d'un élément de la fibre $\mathcal{F}_x$. D'où le lemme.
\end{proof}


\begin{lemma}
\label{lemma-sheafify-abelian-presheaf}
Soit $X$ un espace topologique.
Soit $\mathcal{F}$ un préfaisceau abélien sur $X$.
Il existe alors une unique structure de
faisceau abélien sur $\mathcal{F}^\#$ telle que
$\mathcal{F} \to \mathcal{F}^\#$ soit un morphisme
de préfaisceaux abéliens. En outre, on a la propriété
d'adjonction suivante :
$$
\Mor_{\textit{PAb}(X)}(\mathcal{F}, i(\mathcal{G}))
=
\Mor_{\textit{Ab}(X)}(\mathcal{F}^\#, \mathcal{G}).
$$
\end{lemma}

\begin{proof}
Rappelons le faisceau d'ensembles $\Pi(\mathcal{F})$ défini dans la
Section \ref{section-sheafification}. Toutes les fibres
$\mathcal{F}_x$ sont des groupes abéliens, voir le
Lemme \ref{lemma-stalk-abelian-presheaf}.
Par conséquent, $\Pi(\mathcal{F})$ est un faisceau de groupes abéliens d'après
l'Exemple \ref{example-sheaf-product-pointwise-algebraic-structure}.
Il est également clair que l'application $\mathcal{F} \to \Pi(\mathcal{F})$
est un morphisme de préfaisceaux abéliens. Si l'on montre que
la condition $(*)$ de la Section \ref{section-sheafification} définit un sous-groupe
de $\Pi(\mathcal{F})(U)$ pour tout ouvert $U \subset X$,
alors $\mathcal{F}^\#$ hérite canoniquement d'une structure de faisceau abélien.
Cela se vérifie aisément directement, et nous laissons au
lecteur le soin d'en trouver un argument simple. Voici l'argument que nous employons,
qui se généralise aux préfaisceaux de structures algébriques :
le Lemme \ref{lemma-diagram-fibre-product} montre que
$\mathcal{F}^\#(U)$ est le produit fibré d'un diagramme
de groupes abéliens. Ainsi $\mathcal{F}^\#$ est bien le
sous-groupe abélien voulu.

\medskip\noindent
Remarquons qu'à ce stade $\mathcal{F}^\#_x$ est un groupe
abélien d'après le Lemme \ref{lemma-stalk-abelian-presheaf},
et que $\mathcal{F}_x \to \mathcal{F}^\#_x$ est à la fois une
bijection (Lemme \ref{lemma-stalk-sheafification})
et un homomorphisme de groupes abéliens. Par conséquent,
$\mathcal{F}_x \to \mathcal{F}^\#_x$ est un isomorphisme
de groupes abéliens. Nous l'utiliserons ci-dessous sans autre mention.

\medskip\noindent
Pour démontrer la propriété d'adjonction, nous utilisons celle du
foncteur faisceau associé pour les préfaisceaux d'ensembles. Par exemple,
si $\psi : \mathcal{F} \to i(\mathcal{G})$ est un morphisme de préfaisceaux,
on obtient un morphisme de faisceaux
$\psi' : \mathcal{F}^\# \to \mathcal{G}$. Il reste à vérifier
qu'il s'agit d'un morphisme de faisceaux abéliens. On peut le faire,
par exemple, en constatant que c'est vrai sur les fibres,
d'après le Lemme \ref{lemma-stalk-sheafification}, puis en appliquant le
Lemme \ref{lemma-check-homomorphism-stalks} ci-dessus.
\end{proof}



\section{Faisceau associé à un préfaisceau de structures algébriques}
\label{section-sheafification-presheaves-structures}


\begin{lemma}
\label{lemma-sheafify-presheaf-structures}
Soit $X$ un espace topologique.
Soit $(\mathcal{C}, F)$ une espèce de structure algébrique.
Soit $\mathcal{F}$ un préfaisceau à valeurs dans $\mathcal{C}$
sur $X$. Il existe alors un faisceau $\mathcal{F}^\#$ à valeurs
dans $\mathcal{C}$ et un morphisme $\mathcal{F} \to \mathcal{F}^\#$
de préfaisceaux à valeurs dans $\mathcal{C}$ possédant les
propriétés suivantes :
\begin{enumerate}
\item L'application $\mathcal{F} \to \mathcal{F}^\#$ identifie
le faisceau d'ensembles sous-jacent à $\mathcal{F}^\#$ au
faisceau associé au préfaisceau d'ensembles sous-jacent à $\mathcal{F}$.
\item Pour tout morphisme $\mathcal{F} \to \mathcal{G}$, où
$\mathcal{G}$ est un faisceau à valeurs dans $\mathcal{C}$, il existe
une unique factorisation $\mathcal{F} \to \mathcal{F}^\# \to \mathcal{G}$.
\end{enumerate}
\end{lemma}

\begin{proof}
La démonstration est la même que celle du
Lemme \ref{lemma-sheafify-abelian-presheaf},
en appliquant à plusieurs reprises le
Lemme \ref{lemma-image-contained-in} (voir aussi
l'Exemple \ref{example-application-lemma-image-contained-in}).
L'idée principale consiste toutefois à définir $\mathcal{F}^\#(U)$
comme le produit fibré dans $\mathcal{C}$ du diagramme
$$
\xymatrix{
 &
\Pi(\mathcal{F})(U) \ar[d] \\
\prod_{x \in U} \mathcal{F}_x
\ar[r] &
\prod_{x \in U} \Pi(\mathcal{F})_x
}
$$
à comparer au Lemme \ref{lemma-diagram-fibre-product}.
\end{proof}

\section{Faisceau associé à un préfaisceau de modules}
\label{section-sheafification-presheaves-modules}

\begin{lemma}
\label{lemma-sheafification-presheaf-modules}
Soit $X$ un espace topologique.
Soit $\mathcal{O}$ un préfaisceau d'anneaux sur $X$.
Soit $\mathcal{F}$ un préfaisceau de $\mathcal{O}$-modules.
Soit $\mathcal{O}^\#$ le faisceau associé à $\mathcal{O}$.
Soit $\mathcal{F}^\#$ le faisceau associé à $\mathcal{F}$
considéré comme préfaisceau de groupes abéliens. Il existe un morphisme de
faisceaux d'ensembles
$$
\mathcal{O}^\# \times \mathcal{F}^\#
\longrightarrow
\mathcal{F}^\#
$$
qui rend commutatif le diagramme
$$
\xymatrix{
\mathcal{O} \times \mathcal{F} \ar[r] \ar[d] &
\mathcal{F} \ar[d] \\
\mathcal{O}^\# \times \mathcal{F}^\# \ar[r] &
\mathcal{F}^\#
}
$$
et qui munit $\mathcal{F}^\#$ d'une structure de faisceau
de $\mathcal{O}^\#$-modules. En outre, si $\mathcal{G}$
est un faisceau de $\mathcal{O}^\#$-modules, tout morphisme
de préfaisceaux de $\mathcal{O}$-modules $\mathcal{F} \to \mathcal{G}$
(où $\mathcal{G}$ est considéré, par restriction des scalaires, comme $\mathcal{O}$-module)
se factorise de manière unique sous la forme $\mathcal{F} \to \mathcal{F}^\# \to \mathcal{G}$,
où $\mathcal{F}^\# \to \mathcal{G}$ est un morphisme de
$\mathcal{O}^\#$-modules.
\end{lemma}

\begin{proof}
Omis.
\end{proof}

\noindent
Cela signifie en fait que le foncteur
$i : \textit{Mod}(\mathcal{O}^\#) \to \textit{PMod}(\mathcal{O})$
(qui combine la restriction des scalaires et l'inclusion des faisceaux dans les préfaisceaux)
et le foncteur faisceau associé du lemme
${}^\# : \textit{PMod}(\mathcal{O}) \to \textit{Mod}(\mathcal{O}^\#)$
sont adjoints. En formule :
$$
\Mor_{\textit{PMod}(\mathcal{O})}(\mathcal{F}, i\mathcal{G})
=
\Mor_{\textit{Mod}(\mathcal{O}^\#)}(\mathcal{F}^\#, \mathcal{G})
$$

\medskip\noindent
Soit $X$ un espace topologique.
Soit $\mathcal{O}_1 \to \mathcal{O}_2$ un
morphisme de faisceaux d'anneaux sur $X$.
Dans la Section \ref{section-presheaves-modules},
nous avons défini un foncteur de restriction des scalaires
et un foncteur d'extension des scalaires sur les préfaisceaux de modules
associés à cette situation.

\medskip\noindent
Si $\mathcal{F}$ est un faisceau de $\mathcal{O}_2$-modules,
alors la restriction des scalaires $\mathcal{F}_{\mathcal{O}_1}$
de $\mathcal{F}$ est manifestement un faisceau
de $\mathcal{O}_1$-modules. On obtient le foncteur de restriction des scalaires
$$
\textit{Mod}(\mathcal{O}_2)
\longrightarrow
\textit{Mod}(\mathcal{O}_1)
$$

\medskip\noindent
Inversement, étant donné un faisceau de $\mathcal{O}_1$-modules
$\mathcal{G}$, le préfaisceau de $\mathcal{O}_2$-modules
$\mathcal{O}_2 \otimes_{p, \mathcal{O}_1} \mathcal{G}$
ne constitue pas en général un faisceau. Nous définissons donc le
{\it faisceau produit tensoriel}
$\mathcal{O}_2 \otimes_{\mathcal{O}_1} \mathcal{G}$
par la formule
$$
\mathcal{O}_2 \otimes_{\mathcal{O}_1} \mathcal{G}
=
(\mathcal{O}_2 \otimes_{p, \mathcal{O}_1} \mathcal{G})^\#
$$
comme le faisceau associé à la construction faite pour les préfaisceaux.
On obtient ainsi le foncteur {\it d'extension des scalaires}
$$
\textit{Mod}(\mathcal{O}_1)
\longrightarrow
\textit{Mod}(\mathcal{O}_2)
$$

\begin{lemma}
\label{lemma-adjointness-tensor-restrict}
Avec $X$, $\mathcal{O}_1$, $\mathcal{O}_2$, $\mathcal{F}$ et
$\mathcal{G}$ comme ci-dessus, on a une bijection canonique
$$
\Hom_{\mathcal{O}_1}(\mathcal{G}, \mathcal{F}_{\mathcal{O}_1})
=
\Hom_{\mathcal{O}_2}(
\mathcal{O}_2 \otimes_{\mathcal{O}_1} \mathcal{G},
\mathcal{F}
)
$$
Autrement dit, les foncteurs de restriction et d'extension des scalaires
sont adjoints l'un de l'autre.
\end{lemma}

\begin{proof}
Cela résulte du
Lemme \ref{lemma-adjointness-tensor-restrict-presheaves}
et du fait que
$\Hom_{\mathcal{O}_2}(
\mathcal{O}_2 \otimes_{\mathcal{O}_1} \mathcal{G},
\mathcal{F}
)
=
\Hom_{\mathcal{O}_2}(
\mathcal{O}_2 \otimes_{p, \mathcal{O}_1} \mathcal{G},
\mathcal{F}
)$
puisque $\mathcal{F}$ est un faisceau.
\end{proof}

\begin{lemma}
\label{lemma-stalk-tensor-sheaf-modules}
Soit $X$ un espace topologique.
Soit $\mathcal{O} \to \mathcal{O}'$ un morphisme de
faisceaux d'anneaux sur $X$.
Soit $\mathcal{F}$ un faisceau de $\mathcal{O}$-modules.
Soit $x \in X$. On a
$$
\mathcal{F}_x \otimes_{\mathcal{O}_x} \mathcal{O}'_x
=
(\mathcal{F} \otimes_\mathcal{O} \mathcal{O}')_x
$$
en tant que $\mathcal{O}'_x$-modules.
\end{lemma}

\begin{proof}
Cela résulte directement du Lemme \ref{lemma-stalk-tensor-presheaf-modules}
et du fait que le passage aux fibres commute au passage au faisceau associé.
\end{proof}













\section{Applications continues et faisceaux}
\label{section-presheaves-functorial}

\noindent
Soit $f : X \to Y$ une application continue d'espaces topologiques.
Nous allons définir les foncteurs image directe et image inverse pour les
préfaisceaux et les faisceaux.

\medskip\noindent
Soit $\mathcal{F}$ un préfaisceau d'ensembles sur $X$. Nous définissons
l'{\it image directe} de $\mathcal{F}$ par la règle
$$
f_*\mathcal{F}(V) = \mathcal{F}(f^{-1}(V))
$$
pour tout ouvert $V \subset Y$.
Étant donnés des ouverts $V_1 \subset V_2 \subset Y$, l'application de restriction
est définie par la commutativité du diagramme
$$
\xymatrix{
f_*\mathcal{F}(V_2) \ar[d] \ar@{=}[r] &
\mathcal{F}(f^{-1}(V_2)) \ar[d]^{\text{restriction de }\mathcal{F}} \\
f_*\mathcal{F}(V_1) \ar@{=}[r] &
\mathcal{F}(f^{-1}(V_1))
}
$$
Il est clair que ceci définit un préfaisceau d'ensembles. La construction
est manifestement fonctorielle en $\mathcal{F}$ ; nous
obtenons donc un foncteur
$$
f_* : \textit{PSh}(X) \longrightarrow \textit{PSh}(Y).
$$

\begin{lemma}
\label{lemma-pushforward-sheaf}
Soit $f : X \to Y$ une application continue.
Soit $\mathcal{F}$ un faisceau d'ensembles sur $X$.
Alors $f_*\mathcal{F}$ est un faisceau sur $Y$.
\end{lemma}

\begin{proof}
Cela résulte immédiatement du fait que,
si $V = \bigcup V_j$ est un recouvrement ouvert de $Y$,
alors $f^{-1}(V) = \bigcup f^{-1}(V_j)$ est un recouvrement ouvert de $X$.
\end{proof}

\noindent
Nous obtenons par conséquent un foncteur
$$
f_* : \Sh(X) \longrightarrow \Sh(Y).
$$
Celui-ci est compatible avec la composition au sens
fort suivant.

\begin{lemma}
\label{lemma-pushforward-composition}
Soient $f : X \to Y$ et $g : Y \to Z$ des applications continues
d'espaces topologiques. Les foncteurs $(g \circ f)_*$
et $g_* \circ f_*$ sont égaux (tant sur les préfaisceaux
que sur les faisceaux d'ensembles).
\end{lemma}

\begin{proof}
Cela vient de ce que $(g \circ f)_*\mathcal{F}(W) =
\mathcal{F}((g \circ f)^{-1}W)$ et
$(g_* \circ f_*)\mathcal{F}(W) = \mathcal{F}(f^{-1} g^{-1} W)$,
et que $(g \circ f)^{-1}W = f^{-1} g^{-1} W$.
\end{proof}

\noindent
Soit $\mathcal{G}$ un préfaisceau d'ensembles sur $Y$.
Le {\it préfaisceau image inverse} $f_p\mathcal{G}$
d'un préfaisceau donné $\mathcal{G}$ est défini comme l'adjoint à gauche
de l'image directe $f_*$ sur les préfaisceaux. Autrement dit, ce
devrait être un préfaisceau $f_p \mathcal{G}$ sur $X$ tel que
$$
\Mor_{\textit{PSh}(X)}(f_p\mathcal{G}, \mathcal{F})
=
\Mor_{\textit{PSh}(Y)}(\mathcal{G}, f_*\mathcal{F}).
$$
Par le lemme de Yoneda, cette propriété détermine l'image inverse de manière unique.
Il se trouve qu'elle existe effectivement.

\begin{lemma}
\label{lemma-pullback-presheaves}
Soit $f : X \to Y$ une application continue.
Il existe un foncteur
$f_p : \textit{PSh}(Y) \to \textit{PSh}(X)$
qui est adjoint à gauche de $f_*$. Pour un préfaisceau
$\mathcal{G}$, il est déterminé par la règle
$$
f_p\mathcal{G}(U) = \colim_{f(U) \subset V} \mathcal{G}(V)
$$
où la limite inductive porte sur la collection des voisinages ouverts
$V$ de $f(U)$ dans $Y$. Les limites inductives sont prises sur des
ensembles partiellement ordonnés filtrants.
(Les applications de restriction de $f_p\mathcal{G}$ sont expliquées dans la preuve.)
\end{lemma}

\begin{proof}
La limite inductive porte sur l'ensemble partiellement ordonné formé des ouverts
$V \subset Y$ qui contiennent $f(U)$, ordonné par l'inclusion opposée.
Cet ensemble partiellement ordonné est filtrant puisque, si $V, V'$ lui appartiennent,
il en va de même de $V \cap V'$. En outre, si
$U_1 \subset U_2$, alors tout voisinage ouvert de $f(U_2)$
est un voisinage ouvert de $f(U_1)$. Le système qui définit
$f_p\mathcal{G}(U_2)$ est donc un sous-système de celui qui définit
$f_p\mathcal{G}(U_1)$, et nous obtenons une application de restriction (par
exemple en appliquant les généralités de Catégories,
lemme \ref{categories-lemma-functorial-colimit}).

\medskip\noindent
Remarquons que la construction de la limite inductive est manifestement fonctorielle
en $\mathcal{G}$, de même que les applications de restriction.
Nous avons donc défini $f_p$ comme un foncteur.

\medskip\noindent
Une petite remarque utile est qu'il existe
une application canonique $\mathcal{G}(U) \to f_p\mathcal{G}(f^{-1}(U))$,
car le système des voisinages ouverts
de $f(f^{-1}(U))$ contient l'élément $U$. Celle-ci est compatible
avec les applications de restriction. Autrement dit, il existe un
morphisme canonique $i_\mathcal{G} : \mathcal{G} \to f_* f_p \mathcal{G}$.

\medskip\noindent
Soit $\mathcal{F}$ un préfaisceau d'ensembles sur $X$.
Supposons que $\psi : f_p\mathcal{G} \to \mathcal{F}$
soit un morphisme de préfaisceaux d'ensembles. Le morphisme correspondant
$\mathcal{G} \to f_*\mathcal{F}$ est le morphisme
$f_*\psi \circ i_\mathcal{G} :
\mathcal{G} \to f_* f_p \mathcal{G} \to f_* \mathcal{F}$.

\medskip\noindent
Une autre petite remarque utile est qu'il existe un
morphisme canonique $c_\mathcal{F} : f_p f_* \mathcal{F} \to \mathcal{F}$.
Soit en effet $U \subset X$ un ouvert.
Pour tout voisinage ouvert $V \supset f(U)$ de $Y$,
il existe une application
$f_*\mathcal{F}(V) = \mathcal{F}(f^{-1}(V))\to \mathcal{F}(U)$,
à savoir l'application de restriction de $\mathcal{F}$. Celle-ci est compatible
avec les applications de restriction entre les valeurs de $\mathcal{F}$
sur les images par $f^{-1}$ des différents ouverts contenant $f(U)$. Nous obtenons donc
une application canonique $f_p f_* \mathcal{F}(U) \to \mathcal{F}(U)$.
Une autre vérification immédiate montre que ces applications sont compatibles
avec les applications de restriction et définissent un morphisme $c_\mathcal{F}$
de préfaisceaux d'ensembles.

\medskip\noindent
Supposons que $\varphi : \mathcal{G} \to f_*\mathcal{F}$
soit un morphisme de préfaisceaux d'ensembles. Considérons $f_p\varphi :
f_p \mathcal{G} \to f_p f_* \mathcal{F}$.
En composant ensuite avec $c_\mathcal{F}$, nous obtenons le morphisme cherché
$c_\mathcal{F} \circ f_p\varphi : f_p\mathcal{G} \to \mathcal{F}$.
Nous omettons la vérification que cette construction est inverse
de la construction dans l'autre sens donnée ci-dessus.
\end{proof}

\begin{lemma}
\label{lemma-stalk-pullback-presheaf}
Soit $f : X \to Y$ une application continue.
Soit $x \in X$. Soit $\mathcal{G}$ un préfaisceau d'ensembles sur $Y$.
Il existe une bijection canonique entre les fibres
$(f_p\mathcal{G})_x = \mathcal{G}_{f(x)}$.
\end{lemma}

\begin{proof}
On peut le voir comme suit :
\begin{eqnarray*}
(f_p\mathcal{G})_x
& = &
\colim_{x \in U} f_p\mathcal{G}(U) \\
& = &
\colim_{x \in U} \colim_{f(U) \subset V} \mathcal{G}(V) \\
& = &
\colim_{f(x) \in V} \mathcal{G}(V) \\
& = &
\mathcal{G}_{f(x)}
\end{eqnarray*}
Nous avons utilisé ici
Catégories, lemme \ref{categories-lemma-colimits-commute},
ainsi que le fait que tout ouvert $V$ de $Y$ contenant $f(x)$
intervient dans la troisième description ci-dessus. Nous omettons les détails.
\end{proof}

\noindent
Soit $\mathcal{G}$ un faisceau d'ensembles sur $Y$.
Le {\it faisceau image inverse} $f^{-1}\mathcal{G}$ est défini
par la formule
$$
f^{-1}\mathcal{G} = (f_p\mathcal{G})^\# .
$$
Le foncteur image inverse $f^{-1}$ est adjoint à gauche du
foncteur image directe sur les faisceaux. Autrement dit,
$$
\Mor_{\Sh(X)}(f^{-1}\mathcal{G}, \mathcal{F})
=
\Mor_{\Sh(Y)}(\mathcal{G}, f_*\mathcal{F}).
$$
En effet, nous avons
\begin{eqnarray*}
\Mor_{\Sh(X)}(f^{-1}\mathcal{G}, \mathcal{F})
& = &
\Mor_{\textit{PSh}(X)}(f_p\mathcal{G}, \mathcal{F}) \\
& = &
\Mor_{\textit{PSh}(Y)}(\mathcal{G}, f_*\mathcal{F}) \\
& = &
\Mor_{\Sh(Y)}(\mathcal{G}, f_*\mathcal{F})
\end{eqnarray*}
Pour la première égalité, nous utilisons que le foncteur « faisceau associé » est l'adjoint à gauche
du foncteur d'inclusion des faisceaux dans les préfaisceaux. Pour la deuxième égalité,
nous utilisons que $f_p$ est adjoint à gauche de $f_*$ sur les préfaisceaux. Nous reviendrons
sur cet énoncé dans la preuve du lemme \ref{lemma-f-map}.

\begin{lemma}
\label{lemma-stalk-pullback}
Soit $x \in X$. Soit $\mathcal{G}$ un faisceau d'ensembles sur $Y$.
Il existe une bijection canonique entre les fibres
$(f^{-1}\mathcal{G})_x = \mathcal{G}_{f(x)}$.
\end{lemma}

\begin{proof}
C'est une combinaison des lemmes \ref{lemma-stalk-sheafification}
et \ref{lemma-stalk-pullback-presheaf}.
\end{proof}

\begin{lemma}
\label{lemma-pullback-composition}
Soient $f : X \to Y$ et $g : Y \to Z$ des applications continues
d'espaces topologiques. Les foncteurs $(g \circ f)^{-1}$
et $f^{-1} \circ g^{-1}$ sont canoniquement isomorphes.
De même, $(g \circ f)_p \cong f_p \circ g_p$ sur les
préfaisceaux.
\end{lemma}

\begin{proof}
Pour le voir,
utilisons le fait que les foncteurs adjoints sont uniques à isomorphisme unique près,
ainsi que le lemme \ref{lemma-pushforward-composition}.
\end{proof}

\begin{definition}
\label{definition-f-map}
Soit $f : X \to Y$ une application continue.
Soit $\mathcal{F}$ un faisceau d'ensembles sur $X$ et
soit $\mathcal{G}$ un faisceau d'ensembles sur $Y$.
Un {\it $f$-morphisme $\xi : \mathcal{G} \to \mathcal{F}$}
est une famille d'applications
$\xi_V : \mathcal{G}(V) \to \mathcal{F}(f^{-1}(V))$
indexée par les ouverts $V \subset Y$ telle que
$$
\xymatrix{
\mathcal{G}(V) \ar[r]_{\xi_V} \ar[d]_{\text{restriction de }\mathcal{G}} &
\mathcal{F}(f^{-1}V) \ar[d]^{\text{restriction de }\mathcal{F}} \\
\mathcal{G}(V') \ar[r]^{\xi_{V'}} &
\mathcal{F}(f^{-1}V')
}
$$
soit commutatif pour tous ouverts $V' \subset V \subset Y$.
\end{definition}

\noindent
Dans la littérature, on rencontre parfois une autre définition,
comme dans la partie (4) du lemme \ref{lemma-f-map}, mais le lemme montre qu'il
n'y a en réalité aucune différence.

\begin{lemma}
\label{lemma-f-map}
Soit $f : X \to Y$ une application continue.
Il existe des bijections entre les quatre ensembles suivants :
\begin{enumerate}
\item l'ensemble des morphismes $\mathcal{G} \to f_*\mathcal{F}$,
\item l'ensemble des morphismes $f^{-1}\mathcal{G} \to \mathcal{F}$,
\item l'ensemble des $f$-morphismes $\xi : \mathcal{G} \to \mathcal{F}$, et
\item l'ensemble de toutes les familles d'applications
$\xi_{U, V} : \mathcal{G}(V) \to \mathcal{F}(U)$ pour tous
ouverts $U \subset X$ et $V \subset Y$ tels que $f(U) \subset V$,
compatibles avec toutes les applications de restriction,
\end{enumerate}
fonctoriellement en $\mathcal{F} \in \Sh(X)$ et $\mathcal{G} \in \Sh(Y)$.
\end{lemma}

\begin{proof}
Un morphisme de faisceaux $a : \mathcal{G} \to f_*\mathcal{F}$
est par définition une règle qui associe à chaque ouvert $V$ de $Y$
une application $a_V : \mathcal{G}(V) \to f_*\mathcal{F}(V)$,
et nous avons $f_*\mathcal{F}(V) = \mathcal{F}(f^{-1}(V))$.
Ainsi, pour ce qui est des « données », on obtient exactement ce qu'il faut
pour un $f$-morphisme $\xi$ de $\mathcal{G}$ vers $\mathcal{F}$.
Pour montrer que les ensembles (1) et (3) sont en bijection, nous observons que
$a$ est un morphisme de faisceaux si et seulement si la famille d'applications $a_V$
correspondante satisfait la condition de la définition \ref{definition-f-map}.

\medskip\noindent
Rappelons que $f^{-1}\mathcal{G}$ est le faisceau associé au préfaisceau $f_p\mathcal{G}$.
Par la propriété universelle du faisceau associé, un morphisme de faisceaux
$b : f^{-1}\mathcal{G} \to \mathcal{F}$ équivaut à un morphisme de
préfaisceaux $b_p : f_p\mathcal{G} \to \mathcal{F}$, où $f_p$ est le
foncteur défini plus haut dans cette section. Pour donner un tel morphisme $b_p$,
il faut spécifier, pour chaque ouvert $U$ de $X$, une application
$$
b_{p, U} :
\colim_{f(U) \subset V} \mathcal{G}(V)
\longrightarrow
\mathcal{F}(U)
$$
compatible avec les applications de restriction. Nous regardons $b_{p, U}$
comme une famille d'applications $b_{p, U, V} : \mathcal{G}(V) \to \mathcal{F}(U)$
pour tout ouvert $V$ de $Y$ tel que $f(U) \subset V$. Ces applications doivent être
compatibles avec toutes les applications de restriction possibles.
Autrement dit, nous voyons que $b_p$ correspond à une famille d'applications
comme en (4). Réciproquement, une telle famille définit bien sûr un morphisme
$b_p$, puis un morphisme $b : f^{-1}\mathcal{G} \to \mathcal{F}$.

\medskip\noindent
Pour achever la preuve du lemme, il faut montrer que, par
« oubli de structure », la règle qui, à une famille $\xi_{U, V}$
comme en (4), associe le $f$-morphisme $\xi$ défini par $\xi_V = \xi_{f^{-1}(V), V}$
est bijective. Pour cela, si $\xi$ est un $f$-morphisme usuel, nous
définissons simplement $\tilde \xi_{U, V}$ comme la composée de
$\xi_V : \mathcal{G}(V) \to \mathcal{F}(f^{-1}(V))$
avec l'application de restriction $\mathcal{F}(f^{-1}(V)) \to \mathcal{F}(U)$,
qui a un sens précisément parce que $f(U) \subset V$, c'est-à-dire
$U \subset f^{-1}(V)$. Ceci achève la preuve.
\end{proof}

\noindent
Il est parfois commode de considérer les $f$-morphismes
plutôt que les morphismes entre faisceaux, soit sur $X$, soit sur $Y$.
Nous définissons comme suit la composition des $f$-morphismes.

\begin{definition}
\label{definition-composition-f-maps}
Supposons que $f : X \to Y$ et $g : Y \to Z$ soient des applications
continues d'espaces topologiques. Supposons que $\mathcal{F}$ soit
un faisceau sur $X$, que $\mathcal{G}$ soit un faisceau sur $Y$ et que
$\mathcal{H}$ soit un faisceau sur $Z$.
Soit $\varphi : \mathcal{G} \to \mathcal{F}$ un $f$-morphisme.
Soit $\psi : \mathcal{H} \to \mathcal{G}$ un $g$-morphisme.
La {\it composée de $\varphi$ et $\psi$} est le
$(g \circ f)$-morphisme $\varphi \circ \psi$ défini
par la commutativité du diagramme
$$
\xymatrix{
\mathcal{H}(W) \ar[rr]_{(\varphi \circ \psi)_W}
\ar[rd]_{\psi_W} & &
\mathcal{F}(f^{-1}g^{-1}W) \\
&
\mathcal{G}(g^{-1}W)
\ar[ru]_{\varphi_{g^{-1}W}}
}
$$
\end{definition}

\noindent
Nous laissons au lecteur le soin de vérifier que cela fonctionne.
Une autre manière de le voir consiste à considérer
$\varphi \circ \psi$ comme la composée
$$
\mathcal{H}
\xrightarrow{\psi}
g_*\mathcal{G}
\xrightarrow{g_*\varphi}
g_* f_* \mathcal{F} = (g \circ f)_* \mathcal{F}
$$
N'avez-vous pas maintenant l'impression que considérer les $f$-morphismes est
en quelque sorte plus simple ?

\medskip\noindent
Enfin, étant donnés une application continue $f : X \to Y$ et un
$f$-morphisme $\varphi : \mathcal{G} \to \mathcal{F}$, il existe
une application naturelle entre les fibres
$$
\varphi_x : \mathcal{G}_{f(x)} \longrightarrow \mathcal{F}_x
$$
pour tout $x \in X$. Un représentant $(V, s)$
d'un élément de $\mathcal{G}_{f(x)}$ est envoyé sur
l'élément de $\mathcal{F}_x$ ayant pour représentant $(f^{-1}V,
\varphi_V(s))$. Nous laissons au lecteur le soin de vérifier que cette application
est bien définie. Une autre formulation consiste à dire qu'il s'agit de
l'unique application pour laquelle tous les diagrammes
$$
\xymatrix{
\mathcal{F}(f^{-1}V) \ar[r] &
\mathcal{F}_x \\
\mathcal{G}(V) \ar[r] \ar[u]^{\varphi_V} &
\mathcal{G}_{f(x)} \ar[u]^{\varphi_x}
}
$$
(pour tout ouvert $f(x) \in V \subset Y$) sont commutatifs.

\begin{lemma}
\label{lemma-compose-f-maps-stalks}
Supposons que $f : X \to Y$ et $g : Y \to Z$ soient des applications
continues d'espaces topologiques. Supposons que $\mathcal{F}$ soit
un faisceau sur $X$, que $\mathcal{G}$ soit un faisceau sur $Y$ et que
$\mathcal{H}$ soit un faisceau sur $Z$.
Soit $\varphi : \mathcal{G} \to \mathcal{F}$ un $f$-morphisme.
Soit $\psi : \mathcal{H} \to \mathcal{G}$ un $g$-morphisme.
Soit $x \in X$ un point. L'application induite sur les fibres
$(\varphi \circ \psi)_x : \mathcal{H}_{g(f(x))}
\to \mathcal{F}_x$ est la composée
$$
\mathcal{H}_{g(f(x))}
\xrightarrow{\psi_{f(x)}}
\mathcal{G}_{f(x)}
\xrightarrow{\varphi_x}
\mathcal{F}_x
$$
\end{lemma}

\begin{proof}
Cela résulte immédiatement de la définition \ref{definition-composition-f-maps}
et de la définition de l'application sur les fibres ci-dessus.
\end{proof}
\section{Applications continues et faisceaux abéliens}
\label{section-abelian-presheaves-functorial}

\noindent
Soit $f : X \to Y$ une application continue.
Nous affirmons qu'il existe des foncteurs
\begin{eqnarray*}
f_* : \textit{PAb}(X) & \longrightarrow & \textit{PAb}(Y) \\
f_* : \textit{Ab}(X) & \longrightarrow & \textit{Ab}(Y) \\
f_p : \textit{PAb}(Y) & \longrightarrow & \textit{PAb}(X) \\
f^{-1} : \textit{Ab}(Y) & \longrightarrow & \textit{Ab}(X)
\end{eqnarray*}
qui possèdent des propriétés analogues à celles de leurs homologues de la
Section \ref{section-presheaves-functorial}.
Pour le voir, nous procédons comme suit.

\medskip\noindent
Chacun de ces foncteurs sera construit de la même
manière que le foncteur correspondant de la
Section \ref{section-presheaves-functorial}.
Cela fonctionne parce que toutes les limites inductives de cette section
sont filtrantes (mais nous allons donner les détails ci-dessous).

\medskip\noindent
Tout d'abord, étant donnés un préfaisceau abélien $\mathcal{F}$ sur $X$ et
un préfaisceau abélien $\mathcal{G}$ sur $Y$, nous définissons
\begin{eqnarray*}
f_*\mathcal{F}(V) & = & \mathcal{F}(f^{-1}(V)) \\
f_p\mathcal{G}(U) & = & \colim_{f(U) \subset V} \mathcal{G}(V)
\end{eqnarray*}
comme groupes abéliens. Les applications de restriction sont les mêmes que
pour les préfaisceaux d'ensembles (et ce sont
toutes des homomorphismes de groupes abéliens).

\medskip\noindent
Les affectations $\mathcal{F} \mapsto f_*\mathcal{F}$ et
$\mathcal{G} \to f_p\mathcal{G}$ sont des foncteurs sur
les catégories de préfaisceaux de groupes abéliens.
C'est clair : par exemple, un morphisme de préfaisceaux abéliens
$\mathcal{G}_1 \to \mathcal{G}_2$ induit un morphisme de
systèmes filtrants
$\{\mathcal{G}_1(V)\}_{f(U) \subset V} \to
\{\mathcal{G}_2(V)\}_{f(U) \subset V}$
dont toutes les composantes sont des homomorphismes ;
il induit donc un homomorphisme de groupes abéliens
$f_p\mathcal{G}_1(U) \to f_p\mathcal{G}_2(U)$.

\medskip\noindent
Les foncteurs $f_*$ et $f_p$ sont adjoints sur la
catégorie des préfaisceaux de groupes abéliens, c'est-à-dire que
$$
\Mor_{\textit{PAb}(X)}(f_p\mathcal{G}, \mathcal{F})
=
\Mor_{\textit{PAb}(Y)}(\mathcal{G}, f_*\mathcal{F}).
$$
Pour le démontrer, remarquons que le morphisme
$i_\mathcal{G} : \mathcal{G} \to f_* f_p\mathcal{G}$ de la preuve
du Lemme \ref{lemma-pullback-presheaves}
est un morphisme de préfaisceaux abéliens. Ainsi, si
$\psi : f_p\mathcal{G} \to \mathcal{F}$
est un morphisme de préfaisceaux abéliens, alors le morphisme correspondant
$\mathcal{G} \to f_*\mathcal{F}$, à savoir
$f_*\psi \circ i_\mathcal{G} :
\mathcal{G} \to f_* f_p \mathcal{G} \to f_* \mathcal{F}$,
est lui aussi un morphisme de préfaisceaux abéliens. Dans l'autre sens,
remarquons que le morphisme
$c_\mathcal{F} : f_p f_* \mathcal{F} \to \mathcal{F}$
de la preuve du Lemme \ref{lemma-pullback-presheaves} est également un
morphisme de préfaisceaux abéliens (car il est construit à partir des applications
de restriction de $\mathcal{F}$, qui sont toutes des homomorphismes). Ainsi,
étant donné un morphisme de préfaisceaux abéliens $\varphi : \mathcal{G} \to f_*\mathcal{F}$,
le morphisme
$c_\mathcal{F} \circ f_p\varphi : f_p\mathcal{G} \to \mathcal{F}$
est lui aussi un morphisme de préfaisceaux abéliens. Puisque ces constructions
$\psi \mapsto f_*\psi$ et $\varphi \mapsto c_\mathcal{F} \circ f_p\varphi$
sont inverses l'une de l'autre sur les morphismes de préfaisceaux d'ensembles,
elles le sont aussi sur les morphismes de préfaisceaux abéliens.

\medskip\noindent
Si $\mathcal{F}$ est un faisceau abélien sur $Y$, alors $f_*\mathcal{F}$
est un faisceau abélien sur $X$. Cela résulte de la définition
d'un faisceau abélien et du résultat correspondant pour les faisceaux d'ensembles ;
voir le Lemme \ref{lemma-pushforward-sheaf}. On obtient ainsi le foncteur
$f_*$ sur la catégorie des faisceaux abéliens.

\medskip\noindent
Nous définissons $f^{-1}\mathcal{G} = (f_p\mathcal{G})^\#$ comme précédemment.
L'adjonction de $f_*$ et $f^{-1}$ s'obtient formellement comme dans
le cas des préfaisceaux d'ensembles. Voici l'argument :
\begin{eqnarray*}
\Mor_{\textit{Ab}(X)}(f^{-1}\mathcal{G}, \mathcal{F})
& = &
\Mor_{\textit{PAb}(X)}(f_p\mathcal{G}, \mathcal{F}) \\
& = &
\Mor_{\textit{PAb}(Y)}(\mathcal{G}, f_*\mathcal{F}) \\
& = &
\Mor_{\textit{Ab}(Y)}(\mathcal{G}, f_*\mathcal{F})
\end{eqnarray*}

\begin{lemma}
\label{lemma-pullback-abelian-stalk}
Soit $f : X \to Y$ une application continue.
\begin{enumerate}
\item Soit $\mathcal{G}$ un préfaisceau abélien sur $Y$.
Soit $x \in X$. La bijection
$\mathcal{G}_{f(x)} \to (f_p\mathcal{G})_x$ du
Lemme \ref{lemma-stalk-pullback-presheaf} est un isomorphisme de groupes abéliens.
\item Soit $\mathcal{G}$ un faisceau abélien sur $Y$.
Soit $x \in X$. La bijection
$\mathcal{G}_{f(x)} \to (f^{-1}\mathcal{G})_x$ du
Lemme \ref{lemma-stalk-pullback} est un isomorphisme de groupes abéliens.
\end{enumerate}
\end{lemma}

\begin{proof}
Omis.
\end{proof}

\noindent
Étant donnés une application continue $f : X \to Y$ et des faisceaux de groupes
abéliens $\mathcal{F}$ sur $X$, $\mathcal{G}$ sur $Y$, la notion
d'{\it $f$-morphisme $\mathcal{G} \to \mathcal{F}$
de faisceaux de groupes abéliens} a un sens. Nous pouvons la définir
exactement comme dans la Définition \ref{definition-f-map} (en remplaçant les applications
d'ensembles par des homomorphismes de groupes abéliens), ou
dire simplement qu'il s'agit d'un morphisme de faisceaux abéliens
$\mathcal{G} \to f_*\mathcal{F}$. Nous emploierons librement cette
notion dans la suite. Le groupe des $f$-morphismes entre
$\mathcal{G}$ et $\mathcal{F}$ sera canoniquement en bijection
avec les groupes
$\Mor_{\textit{Ab}(X)}(f^{-1}\mathcal{G}, \mathcal{F})$
et
$\Mor_{\textit{Ab}(Y)}(\mathcal{G}, f_*\mathcal{F})$.

\medskip\noindent
La composition des $f$-morphismes se définit exactement de la
même manière que pour les $f$-morphismes de faisceaux
d'ensembles. De plus, étant donné un $f$-morphisme $\mathcal{G} \to \mathcal{F}$
comme ci-dessus, les applications induites sur les fibres
$$
\varphi_x : \mathcal{G}_{f(x)} \longrightarrow \mathcal{F}_x
$$
sont des homomorphismes de groupes abéliens.


\section{Applications continues et faisceaux de structures algébriques}
\label{section-presheaves-structures-functorial}

\noindent
Soit $(\mathcal{C}, F)$ une espèce de structure algébrique.
Pour un espace topologique $X$, introduisons les notations suivantes :
\begin{enumerate}
\item $\textit{PSh}(X, \mathcal{C})$ désignera la catégorie
des préfaisceaux à valeurs dans $\mathcal{C}$.
\item $\Sh(X, \mathcal{C})$ désignera la catégorie
des faisceaux à valeurs dans $\mathcal{C}$.
\end{enumerate}
Soit $f : X \to Y$ une application continue d'espaces topologiques.
Les mêmes arguments que dans la section précédente montrent
qu'il existe des foncteurs
\begin{eqnarray*}
f_* : \textit{PSh}(X, \mathcal{C}) &
\longrightarrow &
\textit{PSh}(Y, \mathcal{C}) \\
f_* : \Sh(X, \mathcal{C}) &
\longrightarrow &
\Sh(Y, \mathcal{C}) \\
f_p : \textit{PSh}(Y, \mathcal{C}) &
\longrightarrow &
\textit{PSh}(X, \mathcal{C}) \\
f^{-1} : \Sh(Y, \mathcal{C}) &
\longrightarrow &
\Sh(X, \mathcal{C})
\end{eqnarray*}
construits de la même manière et possédant les mêmes
propriétés que les foncteurs construits pour les
(pré)faisceaux abéliens. En particulier, on a des diagrammes commutatifs
$$
\xymatrix{
\textit{PSh}(X, \mathcal{C}) \ar[r]^{f_*} \ar[d]^F &
\textit{PSh}(Y, \mathcal{C}) \ar[d]^F &
\Sh(X, \mathcal{C}) \ar[r]^{f_*} \ar[d]^F &
\Sh(Y, \mathcal{C}) \ar[d]^F
\\
\textit{PSh}(X) \ar[r]^{f_*} &
\textit{PSh}(Y) &
\Sh(X) \ar[r]^{f_*} &
\Sh(Y)
\\
\textit{PSh}(Y, \mathcal{C}) \ar[r]^{f_p} \ar[d]^F &
\textit{PSh}(X, \mathcal{C}) \ar[d]^F &
\Sh(Y, \mathcal{C}) \ar[r]^{f^{-1}} \ar[d]^F &
\Sh(X, \mathcal{C}) \ar[d]^F
\\
\textit{PSh}(Y) \ar[r]^{f_p} &
\textit{PSh}(X) &
\Sh(Y) \ar[r]^{f^{-1}} &
\Sh(X)
}
$$

\medskip\noindent
Les principales formules à retenir sont les suivantes :
\begin{eqnarray*}
f_*\mathcal{F}(V) & = & \mathcal{F}(f^{-1}(V)) \\
f_p\mathcal{G}(U) & = & \colim_{f(U) \subset V} \mathcal{G}(V) \\
f^{-1}\mathcal{G} & = & (f_p\mathcal{G})^\# \\
(f_p\mathcal{G})_x & = & \mathcal{G}_{f(x)} \\
(f^{-1}\mathcal{G})_x & = & \mathcal{G}_{f(x)}
\end{eqnarray*}
Chacune de ces formules est valable dans la
catégorie $\mathcal{C}$ et redonne, après passage aux ensembles
sous-jacents, la formule correspondante pour les préfaisceaux d'ensembles.
Nous avons en outre les propriétés d'adjonction
\begin{eqnarray*}
\Mor_{\textit{PSh}(X, \mathcal{C})}(f_p\mathcal{G}, \mathcal{F})
& = &
\Mor_{\textit{PSh}(Y, \mathcal{C})}(\mathcal{G}, f_*\mathcal{F}) \\
\Mor_{\Sh(X, \mathcal{C})}(f^{-1}\mathcal{G}, \mathcal{F})
& = &
\Mor_{\Sh(Y, \mathcal{C})}(\mathcal{G}, f_*\mathcal{F}).
\end{eqnarray*}
Pour les démontrer, l'étape principale consiste à construire les morphismes
$$
i_\mathcal{G} : \mathcal{G} \longrightarrow f_*f_p\mathcal{G}
$$
et
$$
c_\mathcal{F} : f_p f_* \mathcal{F} \longrightarrow \mathcal{F}
$$
qui interviennent dans la preuve du Lemme \ref{lemma-pullback-presheaves} comme
morphismes de préfaisceaux à valeurs dans $\mathcal{C}$. Nous pouvons
laisser cette vérification au lecteur, car les constructions sont exactement les mêmes
que pour les préfaisceaux d'ensembles.

\medskip\noindent
Étant donnés une application continue $f : X \to Y$ et des faisceaux de structures
algébriques $\mathcal{F}$ sur $X$, $\mathcal{G}$ sur $Y$, la notion
d'{\it $f$-morphisme $\mathcal{G} \to \mathcal{F}$
de faisceaux de structures algébriques} a un sens. Nous pouvons la définir
exactement comme dans la Définition \ref{definition-f-map} (en remplaçant les applications
d'ensembles par des morphismes de $\mathcal{C}$), ou
dire simplement qu'il s'agit d'un morphisme de faisceaux de structures
algébriques $\mathcal{G} \to f_*\mathcal{F}$. Nous emploierons librement cette
notion dans la suite. L'ensemble des $f$-morphismes entre
$\mathcal{G}$ et $\mathcal{F}$ sera canoniquement en bijection
avec les ensembles
$\Mor_{\Sh(X, \mathcal{C})}(f^{-1}\mathcal{G}, \mathcal{F})$
et
$\Mor_{\Sh(Y, \mathcal{C})}(\mathcal{G}, f_*\mathcal{F})$.

\medskip\noindent
La composition des $f$-morphismes se définit exactement de la
même manière que pour les $f$-morphismes de faisceaux
d'ensembles. De plus, étant donné un $f$-morphisme $\mathcal{G} \to \mathcal{F}$
comme ci-dessus, les applications induites sur les fibres
$$
\varphi_x : \mathcal{G}_{f(x)} \longrightarrow \mathcal{F}_x
$$
sont des homomorphismes de structures algébriques.


\begin{lemma}
\label{lemma-f-map-sets-algebraic-structures}
Soit $f : X \to Y$ une application continue d'espaces topologiques.
Supposons donnés des faisceaux de structures algébriques
$\mathcal{F}$ sur $X$, $\mathcal{G}$ sur $Y$. Soit
$\varphi : \mathcal{G} \to \mathcal{F}$ un $f$-morphisme
des faisceaux d'ensembles sous-jacents. Si, pour tout ouvert $V \subset Y$,
l'application d'ensembles $\varphi_V : \mathcal{G}(V) \to \mathcal{F}(f^{-1}V)$
provient, sur les ensembles sous-jacents, d'un morphisme de $\mathcal{C}$,
alors $\varphi$ provient d'un unique $f$-morphisme entre
les faisceaux de structures algébriques.
\end{lemma}

\begin{proof}
Omis.
\end{proof}




\section{Applications continues et faisceaux de modules}
\label{section-presheaves-modules-functorial}

\noindent
Le cas des faisceaux de modules est plus compliqué.
En effet, le cadre naturel pour définir
les foncteurs image inverse et image directe est celui
des espaces annelés, que nous définirons plus bas. Nous
énonçons d'abord quelques lemmes évidents.

\begin{lemma}
\label{lemma-pushforward-presheaf-module}
Soit $f : X \to Y$ une application continue d'espaces topologiques.
Soit $\mathcal{O}$ un préfaisceau d'anneaux sur $X$. Soit
$\mathcal{F}$ un préfaisceau de $\mathcal{O}$-modules.
Il existe un morphisme naturel de préfaisceaux d'ensembles sous-jacents
$$
f_*\mathcal{O} \times f_*\mathcal{F}
\longrightarrow
f_*\mathcal{F}
$$
qui munit $f_*\mathcal{F}$ d'une structure de préfaisceau de
$f_*\mathcal{O}$-modules. Cette construction est
fonctorielle en $\mathcal{F}$.
\end{lemma}

\begin{proof}
Soit $V \subset Y$ un ouvert. Nous définissons l'application du lemme
comme l'application
$$
f_*\mathcal{O}(V) \times f_*\mathcal{F}(V)
=
\mathcal{O}(f^{-1}V) \times \mathcal{F}(f^{-1}V)
\to
\mathcal{F}(f^{-1}V)
=
f_*\mathcal{F}(V).
$$
Ici, la flèche du milieu est l'application de multiplication sur $X$.
Nous laissons au lecteur le soin de vérifier qu'elle est compatible avec les
applications de restriction et définit une structure de
$f_*\mathcal{O}$-module sur $f_*\mathcal{F}$.
\end{proof}

\begin{lemma}
\label{lemma-pullback-presheaf-module}
Soit $f : X \to Y$ une application continue d'espaces topologiques.
Soit $\mathcal{O}$ un préfaisceau d'anneaux sur $Y$. Soit
$\mathcal{G}$ un préfaisceau de $\mathcal{O}$-modules.
Il existe un morphisme naturel de préfaisceaux d'ensembles sous-jacents
$$
f_p\mathcal{O} \times f_p\mathcal{G}
\longrightarrow
f_p\mathcal{G}
$$
qui munit $f_p\mathcal{G}$ d'une structure de préfaisceau de $f_p\mathcal{O}$-modules.
Cette construction est fonctorielle en $\mathcal{G}$.
\end{lemma}

\begin{proof}
Soit $U \subset X$ un ouvert. Nous définissons l'application du lemme
comme l'application
\begin{eqnarray*}
f_p\mathcal{O}(U) \times f_p\mathcal{G}(U)
& = &
\colim_{f(U) \subset V} \mathcal{O}(V)
\times
\colim_{f(U) \subset V} \mathcal{G}(V) \\
& = &
\colim_{f(U) \subset V} (\mathcal{O}(V)\times \mathcal{G}(V)) \\
& \to &
\colim_{f(U) \subset V} \mathcal{G}(V) \\
& = &
f_p\mathcal{G}(U).
\end{eqnarray*}
Ici, la flèche du milieu est l'application de multiplication sur $Y$.
La seconde égalité vaut parce que les limites inductives filtrantes commutent
aux limites finies ; voir
Catégories, Lemme \ref{categories-lemma-directed-commutes}.
Nous laissons au lecteur le soin de vérifier que cette application est compatible avec les
applications de restriction et définit une structure de
$f_p\mathcal{O}$-module sur $f_p\mathcal{G}$.
\end{proof}

\noindent
Soit $f : X \to Y$ une application continue.
Soit $\mathcal{O}_X$ un préfaisceau d'anneaux sur $X$ et
soit $\mathcal{O}_Y$ un préfaisceau d'anneaux sur $Y$.
À ce stade, nous avons donc défini des foncteurs
\begin{eqnarray*}
f_* : \textit{PMod}(\mathcal{O}_X) &
\longrightarrow &
\textit{PMod}(f_*\mathcal{O}_X) \\
f_p : \textit{PMod}(\mathcal{O}_Y) &
\longrightarrow &
\textit{PMod}(f_p\mathcal{O}_Y)
\end{eqnarray*}
Ils vérifient les compatibilités suivantes.

\begin{lemma}
\label{lemma-adjoint-push-pull-presheaves-modules}
Soit $f : X \to Y$ une application continue d'espaces topologiques.
Soit $\mathcal{O}$ un préfaisceau d'anneaux sur $Y$.
Soit $\mathcal{G}$ un préfaisceau de $\mathcal{O}$-modules.
Soit $\mathcal{F}$ un préfaisceau de $f_p\mathcal{O}$-modules.
Alors
$$
\Mor_{\textit{PMod}(f_p\mathcal{O})}(f_p\mathcal{G}, \mathcal{F})
=
\Mor_{\textit{PMod}(\mathcal{O})}(\mathcal{G}, f_*\mathcal{F}).
$$
Nous utilisons ici
les Lemmes \ref{lemma-pullback-presheaf-module}
et \ref{lemma-pushforward-presheaf-module}, et considérons
$f_*\mathcal{F}$ comme un $\mathcal{O}$-module par restriction des scalaires le long de
$i_\mathcal{O} : \mathcal{O} \to f_*f_p\mathcal{O}$
(défini pour la première fois dans la preuve du Lemme \ref{lemma-pullback-presheaves}).
\end{lemma}

\begin{proof}
Remarquons que
$$
\Mor_{\textit{PAb}(X)}(f_p\mathcal{G}, \mathcal{F})
=
\Mor_{\textit{PAb}(Y)}(\mathcal{G}, f_*\mathcal{F}).
$$
Cette égalité vaut d'après la Section \ref{section-abelian-presheaves-functorial}.
Il reste donc à montrer que, sous cette correspondance, les
sous-ensembles formés des morphismes de modules se correspondent. De plus, la correspondance
est déterminée par la règle
$$
(\psi : f_p\mathcal{G} \to \mathcal{F})
\longmapsto
(f_*\psi \circ i_\mathcal{G} :
\mathcal{G} \to f_* \mathcal{F})
$$
et, dans l'autre sens, par la règle
$$
(\varphi : \mathcal{G} \to f_* \mathcal{F})
\longmapsto
(c_\mathcal{F} \circ f_p\varphi : f_p\mathcal{G} \to \mathcal{F})
$$
où $i_\mathcal{G}$ et $c_\mathcal{F}$ sont ceux de la
Section \ref{section-abelian-presheaves-functorial}.
Ainsi, par fonctorialité de $f_*$ et $f_p$, nous voyons
qu'il suffit de vérifier que les morphismes
$i_\mathcal{G} : \mathcal{G} \to f_* f_p \mathcal{G}$ et
$c_\mathcal{F} : f_p f_* \mathcal{F} \to \mathcal{F}$
sont compatibles avec les structures de modules, ce que nous laissons au lecteur.
\end{proof}

\begin{lemma}
\label{lemma-adjoint-pull-push-presheaves-modules}
Soit $f : X \to Y$ une application continue d'espaces topologiques.
Soit $\mathcal{O}$ un préfaisceau d'anneaux sur $X$.
Soit $\mathcal{F}$ un préfaisceau de $\mathcal{O}$-modules.
Soit $\mathcal{G}$ un préfaisceau de $f_*\mathcal{O}$-modules.
Alors
$$
\Mor_{\textit{PMod}(\mathcal{O})}(
\mathcal{O} \otimes_{p, f_pf_*\mathcal{O}} f_p\mathcal{G}, \mathcal{F})
=
\Mor_{\textit{PMod}(f_*\mathcal{O})}(\mathcal{G}, f_*\mathcal{F}).
$$
Nous utilisons ici
les Lemmes \ref{lemma-pullback-presheaf-module}
et \ref{lemma-pushforward-presheaf-module}, ainsi que
le morphisme $c_\mathcal{O} : f_pf_*\mathcal{O} \to \mathcal{O}$
dans la définition du produit tensoriel.
\end{lemma}

\begin{proof}
Cela résulte des égalités
\begin{eqnarray*}
\Mor_{\textit{PMod}(\mathcal{O})}(
\mathcal{O} \otimes_{p, f_pf_*\mathcal{O}} f_p\mathcal{G}, \mathcal{F})
& = &
\Mor_{\textit{PMod}(f_pf_*\mathcal{O})}(
f_p\mathcal{G}, \mathcal{F}_{f_pf_*\mathcal{O}}) \\
& = &
\Mor_{\textit{PMod}(f_*\mathcal{O})}(\mathcal{G},
f_*(\mathcal{F}_{f_pf_*\mathcal{O}})) \\
& = &
\Mor_{\textit{PMod}(f_*\mathcal{O})}(\mathcal{G}, f_*\mathcal{F}).
\end{eqnarray*}
La première égalité est le
Lemme \ref{lemma-adjointness-tensor-restrict-presheaves}.
La deuxième égalité est le
Lemme \ref{lemma-adjoint-push-pull-presheaves-modules}.
La troisième égalité est donnée par l'égalité
$f_*(\mathcal{F}_{f_pf_*\mathcal{O}}) = f_*\mathcal{F}$
de faisceaux abéliens qui est $f_*\mathcal{O}$-linéaire. En effet,
$\text{id}_{f_*\mathcal{O}}$ correspond à $c_\mathcal{O}$
par l'adjonction décrite dans la preuve du
Lemme \ref{lemma-pullback-presheaves} ;
ainsi
$\text{id}_{f_*\mathcal{O}} = f_*c_\mathcal{O} \circ i_{f_*\mathcal{O}}$.
\end{proof}

\begin{lemma}
\label{lemma-pushforward-module}
Soit $f : X \to Y$ une application continue d'espaces topologiques.
Soit $\mathcal{O}$ un faisceau d'anneaux sur $X$. Soit
$\mathcal{F}$ un faisceau de $\mathcal{O}$-modules.
L'image directe $f_*\mathcal{F}$, telle que définie dans le
Lemme \ref{lemma-pushforward-presheaf-module},
est un faisceau de $f_*\mathcal{O}$-modules.
\end{lemma}

\begin{proof}
C'est évident d'après la définition et le Lemme \ref{lemma-pushforward-sheaf}.
\end{proof}

\begin{lemma}
\label{lemma-pullback-module}
Soit $f : X \to Y$ une application continue d'espaces topologiques.
Soit $\mathcal{O}$ un faisceau d'anneaux sur $Y$. Soit
$\mathcal{G}$ un faisceau de $\mathcal{O}$-modules.
Il existe un morphisme naturel de préfaisceaux d'ensembles sous-jacents
$$
f^{-1}\mathcal{O} \times f^{-1}\mathcal{G}
\longrightarrow
f^{-1}\mathcal{G}
$$
qui munit $f^{-1}\mathcal{G}$ d'une structure de
faisceau de $f^{-1}\mathcal{O}$-modules.
\end{lemma}

\begin{proof}
Rappelons que $f^{-1}$ est défini comme la composée du
foncteur $f_p$ et du foncteur faisceau associé. Le lemme
résulte donc du Lemme \ref{lemma-pullback-presheaf-module}
et du Lemme \ref{lemma-sheafification-presheaf-modules}.
\end{proof}

\noindent
Soit $f : X \to Y$ une application continue.
Soit $\mathcal{O}_X$ un faisceau d'anneaux sur $X$ et
soit $\mathcal{O}_Y$ un faisceau d'anneaux sur $Y$.
Nous avons donc maintenant défini des foncteurs
\begin{eqnarray*}
f_* : \textit{Mod}(\mathcal{O}_X) &
\longrightarrow &
\textit{Mod}(f_*\mathcal{O}_X) \\
f^{-1} : \textit{Mod}(\mathcal{O}_Y) &
\longrightarrow &
\textit{Mod}(f^{-1}\mathcal{O}_Y)
\end{eqnarray*}
Ils vérifient les compatibilités suivantes.

\begin{lemma}
\label{lemma-adjoint-push-pull-modules}
Soit $f : X \to Y$ une application continue d'espaces topologiques.
Soit $\mathcal{O}$ un faisceau d'anneaux sur $Y$.
Soit $\mathcal{G}$ un faisceau de $\mathcal{O}$-modules.
Soit $\mathcal{F}$ un faisceau de $f^{-1}\mathcal{O}$-modules.
Alors
$$
\Mor_{\textit{Mod}(f^{-1}\mathcal{O})}(f^{-1}\mathcal{G}, \mathcal{F})
=
\Mor_{\textit{Mod}(\mathcal{O})}(\mathcal{G}, f_*\mathcal{F}).
$$
Nous utilisons ici
les Lemmes \ref{lemma-pullback-module}
et \ref{lemma-pushforward-module}, et considérons
$f_*\mathcal{F}$ comme un $\mathcal{O}$-module par restriction des scalaires le long de
$\mathcal{O} \to f_*f^{-1}\mathcal{O}$.
\end{lemma}

\begin{proof}
Le résultat découle des égalités
\begin{eqnarray*}
\Mor_{\textit{Mod}(f^{-1}\mathcal{O})}(f^{-1}\mathcal{G}, \mathcal{F})
& = &
\Mor_{\textit{Mod}(f_p\mathcal{O})}(f_p\mathcal{G}, \mathcal{F}) \\
& = &
\Mor_{\textit{Mod}(\mathcal{O})}(\mathcal{G}, f_*\mathcal{F}).
\end{eqnarray*}
La seconde égalité provient du
Lemme \ref{lemma-adjoint-push-pull-presheaves-modules}
et la première du Lemme \ref{lemma-sheafification-presheaf-modules}.
\end{proof}

\begin{lemma}
\label{lemma-adjoint-pull-push-modules}
Soit $f : X \to Y$ une application continue d'espaces topologiques.
Soit $\mathcal{O}$ un faisceau d'anneaux sur $X$.
Soit $\mathcal{F}$ un faisceau de $\mathcal{O}$-modules.
Soit $\mathcal{G}$ un faisceau de $f_*\mathcal{O}$-modules.
Alors
$$
\Mor_{\textit{Mod}(\mathcal{O})}(
\mathcal{O} \otimes_{f^{-1}f_*\mathcal{O}} f^{-1}\mathcal{G}, \mathcal{F})
=
\Mor_{\textit{Mod}(f_*\mathcal{O})}(\mathcal{G}, f_*\mathcal{F}).
$$
Nous utilisons ici
les Lemmes \ref{lemma-pullback-module}
et \ref{lemma-pushforward-module}, ainsi que
le morphisme canonique $f^{-1}f_*\mathcal{O} \to \mathcal{O}$
dans la définition du produit tensoriel.
\end{lemma}

\begin{proof}
Cela résulte des égalités
\begin{eqnarray*}
\Mor_{\textit{Mod}(\mathcal{O})}(
\mathcal{O} \otimes_{f^{-1}f_*\mathcal{O}} f^{-1}\mathcal{G}, \mathcal{F})
& = &
\Mor_{\textit{Mod}(f^{-1}f_*\mathcal{O})}(
f^{-1}\mathcal{G}, \mathcal{F}_{f^{-1}f_*\mathcal{O}}) \\
& = &
\Mor_{\textit{Mod}(f_*\mathcal{O})}(\mathcal{G}, f_*\mathcal{F}).
\end{eqnarray*}
Ces égalités combinent le
Lemme \ref{lemma-adjointness-tensor-restrict}
et le Lemme \ref{lemma-adjoint-push-pull-modules}.
\end{proof}

\noindent
Soit $f : X \to Y$ une application continue.
Soit $\mathcal{O}_X$ un (pré)faisceau d'anneaux sur $X$ et
soit $\mathcal{O}_Y$ un (pré)faisceau d'anneaux sur $Y$.
À ce stade, nous avons donc défini les foncteurs
\begin{eqnarray*}
f_* : \textit{PMod}(\mathcal{O}_X) &
\longrightarrow &
\textit{PMod}(f_*\mathcal{O}_X) \\
f_* : \textit{Mod}(\mathcal{O}_X) &
\longrightarrow &
\textit{Mod}(f_*\mathcal{O}_X) \\
f_p : \textit{PMod}(\mathcal{O}_Y) &
\longrightarrow &
\textit{PMod}(f_p\mathcal{O}_Y) \\
f^{-1} : \textit{Mod}(\mathcal{O}_Y) &
\longrightarrow &
\textit{Mod}(f^{-1}\mathcal{O}_Y)
\end{eqnarray*}
Manifestement, en général, les deux foncteurs $(f_*, f^{-1})$
sur les faisceaux de modules ne sont pas adjoints, car leurs catégories d'arrivée
ne coïncident pas. En effet, comme nous l'avons vu plus haut, cela ne fonctionne que si, par miracle, les
faisceaux d'anneaux $\mathcal{O}_X, \mathcal{O}_Y$ vérifient les
relations $\mathcal{O}_X = f^{-1}\mathcal{O}_Y$ et
$\mathcal{O}_Y = f_*\mathcal{O}_X$. Ce n'est presque jamais
le cas en pratique. Nous interrompons la discussion pour définir
la notion correcte de morphisme pour laquelle il existe un couple
de foncteurs adjoints approprié sur les faisceaux de modules.

\section{Espaces annelés}
\label{section-ringed-spaces}

\noindent
Soit $X$ un espace topologique et soit $\mathcal{O}_X$
un faisceau d'anneaux sur $X$. Nous devons
considérer le faisceau d'anneaux $\mathcal{O}_X$ comme un faisceau de fonctions
sur $X$. Et si $f : X \to Y$ est une application « convenable », alors, par composition,
une fonction sur $Y$ devient une fonction sur $X$.
Il devrait donc exister un
$f$-morphisme naturel de $\mathcal{O}_Y$ vers $\mathcal{O}_X$ ;
voir la définition \ref{definition-f-map} et le
lemme \ref{lemma-f-map}.
Pour un exemple précis, voir
l'exemple \ref{example-continuous-map-ringed} ci-dessous. Voici la
définition abstraite pertinente.

\begin{definition}
\label{definition-ringed-space}
Un {\it espace annelé} est un couple $(X, \mathcal{O}_X)$ formé
d'un espace topologique $X$ et d'un faisceau d'anneaux $\mathcal{O}_X$
sur $X$. Un {\it morphisme d'espaces annelés}
$(X, \mathcal{O}_X) \to (Y, \mathcal{O}_Y)$ est un couple
formé d'une application continue $f : X \to Y$ et d'un
$f$-morphisme de faisceaux d'anneaux
$f^\sharp : \mathcal{O}_Y \to \mathcal{O}_X$.
\end{definition}

\begin{example}
\label{example-continuous-map-ringed}
Soit $f : X \to Y$ une application continue d'espaces topologiques.
Considérons les faisceaux de fonctions continues à valeurs réelles
$\mathcal{C}^0_X$ sur $X$ et $\mathcal{C}^0_Y$ sur $Y$ ;
voir l'exemple \ref{example-C0-sheaf-rings}. Nous affirmons qu'il
existe un $f$-morphisme naturel $f^\sharp : \mathcal{C}^0_Y \to \mathcal{C}^0_X$
associé à $f$. Nous le définissons simplement par la
règle
\begin{eqnarray*}
\mathcal{C}^0_Y(V) & \longrightarrow & \mathcal{C}^0_X(f^{-1}V) \\
h & \longmapsto & h \circ f
\end{eqnarray*}
En toute rigueur, nous devrions écrire $f^\sharp(h) = h \circ f|_{f^{-1}(V)}$.
Il est clair qu'il s'agit d'une famille d'applications comme dans la
définition \ref{definition-f-map}, compatible avec les structures de $\mathbf{R}$-algèbres.
C'est donc un $f$-morphisme de faisceaux de $\mathbf{R}$-algèbres ; voir le
lemme \ref{lemma-f-map-sets-algebraic-structures}.

\medskip\noindent
Il existe bien sûr beaucoup d'autres situations dans lesquelles un
morphisme canonique d'espaces annelés est associé à un type
géométrique de morphisme. Par exemple, si $M$, $N$ sont des variétés $\mathcal{C}^\infty$
et si $f : M \to N$ est une application indéfiniment différentiable, alors
$f$ induit un morphisme canonique d'espaces annelés
$(M, \mathcal{C}_M^\infty) \to (N, \mathcal{C}^\infty_N)$.
La construction (qui est identique à celle ci-dessus) est laissée
au lecteur.
\end{example}

\noindent
La manière de composer les morphismes d'espaces
annelés n'étant peut-être pas tout à fait évidente, nous l'explicitons ici.

\begin{definition}
\label{definition-composition-maps-ringed-spaces}
Soient
$(f, f^\sharp) : (X, \mathcal{O}_X) \to (Y, \mathcal{O}_Y)$ et
$(g, g^\sharp) : (Y, \mathcal{O}_Y) \to (Z, \mathcal{O}_Z)$
des morphismes d'espaces annelés. Nous définissons alors
la {\it composée de morphismes d'espaces annelés}
par la règle
$$
(g, g^\sharp) \circ (f, f^\sharp) = (g \circ f, f^\sharp \circ g^\sharp).
$$
Nous utilisons ici la composition des $f$-morphismes définie dans la
définition \ref{definition-composition-f-maps}.
\end{definition}


\section{Morphismes d'espaces annelés et modules}
\label{section-ringed-spaces-functoriality-modules}

\noindent
Nous avons maintenant introduit assez de notations pour pouvoir
définir l'image inverse et l'image directe des modules le long d'un morphisme
d'espaces annelés.

\begin{definition}
\label{definition-pushforward}
Soit $(f, f^\sharp) : (X, \mathcal{O}_X) \to (Y, \mathcal{O}_Y)$
un morphisme d'espaces annelés.
\begin{enumerate}
\item Soit $\mathcal{F}$ un faisceau de $\mathcal{O}_X$-modules.
Nous définissons l'{\it image directe} de $\mathcal{F}$ comme le
faisceau de $\mathcal{O}_Y$-modules qui, comme faisceau
de groupes abéliens, est égal à $f_*\mathcal{F}$ et dont la
structure de module est obtenue par restriction des scalaires
via $f^\sharp : \mathcal{O}_Y \to f_*\mathcal{O}_X$
de la structure de module donnée
dans le lemme \ref{lemma-pushforward-module}.
\item Soit $\mathcal{G}$ un faisceau de $\mathcal{O}_Y$-modules.
Nous définissons l'{\it image inverse} $f^*\mathcal{G}$ comme le
faisceau de $\mathcal{O}_X$-modules défini par la formule
$$
f^*\mathcal{G}
=
\mathcal{O}_X \otimes_{f^{-1}\mathcal{O}_Y} f^{-1}\mathcal{G}
$$
où l'homomorphisme d'anneaux $f^{-1}\mathcal{O}_Y \to \mathcal{O}_X$
est l'homomorphisme correspondant à $f^\sharp$, et où la structure de module
est donnée par le lemme \ref{lemma-pullback-module}.
\end{enumerate}
\end{definition}

\noindent
Nous avons ainsi défini des foncteurs
\begin{eqnarray*}
f_* : \textit{Mod}(\mathcal{O}_X)
& \longrightarrow &
\textit{Mod}(\mathcal{O}_Y) \\
f^* : \textit{Mod}(\mathcal{O}_Y)
& \longrightarrow &
\textit{Mod}(\mathcal{O}_X)
\end{eqnarray*}
Le dernier résultat concernant ces foncteurs affirme qu'ils sont bien
adjoints, comme prévu.

\begin{lemma}
\label{lemma-adjoint-pullback-pushforward-modules}
Soit $(f, f^\sharp) : (X, \mathcal{O}_X) \to (Y, \mathcal{O}_Y)$
un morphisme d'espaces annelés.
Soit $\mathcal{F}$ un faisceau de $\mathcal{O}_X$-modules.
Soit $\mathcal{G}$ un faisceau de $\mathcal{O}_Y$-modules.
Il existe une bijection canonique
$$
\Hom_{\mathcal{O}_X}(f^*\mathcal{G}, \mathcal{F})
=
\Hom_{\mathcal{O}_Y}(\mathcal{G}, f_*\mathcal{F}).
$$
Autrement dit, le foncteur $f^*$ est adjoint à gauche de
$f_*$.
\end{lemma}

\begin{proof}
Cela résulte du travail effectué précédemment :
\begin{eqnarray*}
\Hom_{\mathcal{O}_X}(f^*\mathcal{G}, \mathcal{F})
& = &
\Mor_{\textit{Mod}(\mathcal{O}_X)}(
\mathcal{O}_X \otimes_{f^{-1}\mathcal{O}_Y} f^{-1}\mathcal{G}, \mathcal{F}) \\
& = &
\Mor_{\textit{Mod}(f^{-1}\mathcal{O}_Y)}(
f^{-1}\mathcal{G}, \mathcal{F}_{f^{-1}\mathcal{O}_Y}) \\
& = &
\Hom_{\mathcal{O}_Y}(\mathcal{G}, f_*\mathcal{F}).
\end{eqnarray*}
Nous utilisons ici les lemmes \ref{lemma-adjointness-tensor-restrict}
et \ref{lemma-adjoint-push-pull-modules}.
\end{proof}

\begin{lemma}
\label{lemma-push-pull-composition-modules}
Soient $f : X \to Y$ et $g : Y \to Z$ des morphismes d'espaces annelés.
Les foncteurs $(g \circ f)_*$ et $g_* \circ f_*$ sont égaux.
Il existe un isomorphisme canonique de foncteurs
$(g \circ f)^* \cong f^* \circ g^*$.
\end{lemma}

\begin{proof}
Le résultat sur les images directes est une conséquence du lemme
\ref{lemma-pushforward-composition} et de nos définitions.
Le résultat sur les images inverses s'en déduit par le même
argument que dans la preuve du lemme \ref{lemma-pullback-composition}.
\end{proof}


\noindent
Étant donnés un morphisme d'espaces annelés
$(f, f^\sharp) : (X, \mathcal{O}_X) \to (Y, \mathcal{O}_Y)$,
un faisceau de $\mathcal{O}_X$-modules $\mathcal{F}$
et un faisceau de $\mathcal{O}_Y$-modules $\mathcal{G}$ sur $Y$,
la notion d'{\it $f$-morphisme $\varphi : \mathcal{G} \to \mathcal{F}$
de faisceaux de modules} a un sens. Nous pouvons simplement le définir
comme un $f$-morphisme $\varphi : \mathcal{G} \to \mathcal{F}$
de faisceaux abéliens (voir la définition \ref{definition-f-map} et le
lemme \ref{lemma-f-map}) tel que, pour tout ouvert $V \subset Y$, l'application
$$
\mathcal{G}(V) \longrightarrow \mathcal{F}(f^{-1}V)
$$
soit un morphisme de $\mathcal{O}_Y(V)$-modules. Nous considérons ici
$\mathcal{F}(f^{-1}V)$ comme un $\mathcal{O}_Y(V)$-module
au moyen de l'homomorphisme $f^\sharp_V : \mathcal{O}_Y(V) \to \mathcal{O}_X(f^{-1}V)$.
L'ensemble des $f$-morphismes entre
$\mathcal{G}$ et $\mathcal{F}$ est en bijection canonique
avec les ensembles
$\Mor_{\textit{Mod}(\mathcal{O}_X)}(f^*\mathcal{G}, \mathcal{F})$
et
$\Mor_{\textit{Mod}(\mathcal{O}_Y)}(\mathcal{G}, f_*\mathcal{F})$.
Voir ci-dessus.

\medskip\noindent
La composition des $f$-morphismes est définie exactement de la
même manière que dans le cas des $f$-morphismes de faisceaux
d'ensembles. En outre, étant donnés un $f$-morphisme $\mathcal{G} \to \mathcal{F}$
comme ci-dessus et $x \in X$, l'application induite sur les fibres
$$
\varphi_x : \mathcal{G}_{f(x)} \longrightarrow \mathcal{F}_x
$$
est un morphisme de $\mathcal{O}_{Y, f(x)}$-modules, où la
structure de $\mathcal{O}_{Y, f(x)}$-module sur $\mathcal{F}_x$
provient de la structure de $\mathcal{O}_{X, x}$-module au moyen de
l'application $f^\sharp_x : \mathcal{O}_{Y, f(x)} \to \mathcal{O}_{X, x}$.
Voici un lemme apparenté.

\begin{lemma}
\label{lemma-stalk-pullback-modules}
Soit $(f, f^\sharp) : (X, \mathcal{O}_X) \to (Y, \mathcal{O}_Y)$
un morphisme d'espaces annelés.
Soit $\mathcal{G}$ un faisceau de $\mathcal{O}_Y$-modules.
Soit $x \in X$. Alors
$$
(f^*\mathcal{G})_x =
\mathcal{G}_{f(x)}
\otimes_{\mathcal{O}_{Y, f(x)}}
\mathcal{O}_{X, x}
$$
comme $\mathcal{O}_{X, x}$-modules, le produit tensoriel de droite
utilisant $f^\sharp_x : \mathcal{O}_{Y, f(x)} \to \mathcal{O}_{X, x}$.
\end{lemma}

\begin{proof}
Cela résulte du lemme \ref{lemma-stalk-tensor-sheaf-modules}
et de l'identification des fibres des faisceaux images inverses
en $x$ avec les fibres correspondantes en $f(x)$. Voir par
exemple les formules de la section \ref{section-presheaves-structures-functorial}.
\end{proof}



\section{Faisceaux gratte-ciel et fibres}
\label{section-skyscraper-sheaves}

\begin{definition}
\label{definition-skyscraper-sheaf}
Soit $X$ un espace topologique.
\begin{enumerate}
\item Soit $x \in X$ un point. Notons $i_x : \{x\} \to X$ l'application d'inclusion.
Soit $A$ un ensemble et considérons $A$ comme un faisceau sur l'espace à un point $\{x\}$.
Nous appelons $i_{x, *}A$ le {\it faisceau gratte-ciel en $x$ de valeur $A$}.
\item Si, dans (1) ci-dessus, $A$ est un groupe abélien, nous considérons
$i_{x, *}A$ comme un faisceau de groupes abéliens sur $X$.
\item Si, dans (1) ci-dessus, $A$ est une structure algébrique, nous considérons
$i_{x, *}A$ comme un faisceau de structures algébriques.
\item Si $(X, \mathcal{O}_X)$ est un espace annelé, nous considérons
$i_x : \{x\} \to X$ comme un morphisme d'espaces annelés
$(\{x\}, \mathcal{O}_{X, x}) \to (X, \mathcal{O}_X)$
et, si $A$ est un $\mathcal{O}_{X, x}$-module, nous considérons
$i_{x, *}A$ comme un faisceau de $\mathcal{O}_X$-modules.
\item On dit qu'un faisceau d'ensembles $\mathcal{F}$ est un {\it faisceau gratte-ciel}
s'il existe un point $x$ de $X$ et un ensemble $A$ tels
que $\mathcal{F} \cong i_{x, *}A$.
\item On dit qu'un faisceau de groupes abéliens $\mathcal{F}$ est un
{\it faisceau gratte-ciel} s'il existe un point $x$ de $X$
et un groupe abélien $A$ tels que $\mathcal{F} \cong i_{x, *}A$
comme faisceaux de groupes abéliens.
\item On dit qu'un faisceau de structures algébriques $\mathcal{F}$ est un
{\it faisceau gratte-ciel} s'il existe un point $x$ de $X$
et une structure algébrique $A$ tels que $\mathcal{F} \cong i_{x, *}A$
comme faisceaux de structures algébriques.
\item Si $(X, \mathcal{O}_X)$ est un espace annelé et
si $\mathcal{F}$ est un faisceau de $\mathcal{O}_X$-modules,
on dit que $\mathcal{F}$ est un {\it faisceau gratte-ciel} s'il
existe un point $x \in X$ et un $\mathcal{O}_{X, x}$-module
$A$ tels que $\mathcal{F} \cong i_{x, *}A$
comme faisceaux de $\mathcal{O}_X$-modules.
\end{enumerate}
\end{definition}

\begin{lemma}
\label{lemma-skyscraper-stalks}
Soient $X$ un espace topologique, $x \in X$ un point et
$A$ un ensemble. Pour tout point $x' \in X$, la fibre du
faisceau gratte-ciel en $x$ de valeur $A$ au point $x'$ est
$$
(i_{x, *}A)_{x'} =
\left\{
\begin{matrix}
A & \text{si} & x' \in \overline{\{x\}} \\
\{*\} & \text{si} & x' \not\in \overline{\{x\}}
\end{matrix}
\right.
$$
Une description analogue vaut dans le cas des
groupes abéliens, des structures algébriques et des
faisceaux de modules.
\end{lemma}

\begin{proof}
Démonstration omise.
\end{proof}

\begin{lemma}
\label{lemma-stalk-skyscraper-adjoint}
Soit $X$ un espace topologique et soit $x \in X$ un point.
Les foncteurs $\mathcal{F} \mapsto \mathcal{F}_x$ et
$A \mapsto i_{x, *}A$ sont adjoints. Plus précisément,
$$
\Mor_{\textit{Sets}}(\mathcal{F}_x, A)
=
\Mor_{\Sh(X)}(\mathcal{F}, i_{x, *}A).
$$
Un énoncé analogue vaut dans le cas des
groupes abéliens et des structures algébriques. Dans le cas des
faisceaux de modules, nous avons
$$
\Hom_{\mathcal{O}_{X, x}}(\mathcal{F}_x, A)
=
\Hom_{\mathcal{O}_X}(\mathcal{F}, i_{x, *}A).
$$
\end{lemma}

\begin{proof}
Démonstration omise. Indication : le foncteur fibre peut être
considéré comme le foncteur image inverse pour le morphisme $i_x : \{x\} \to X$.
L'adjonction résulte alors de celle de
$i_x^{-1}$ et $i_{x, *}$ (resp.\ $i_x^*$ et $i_{x, *}$
dans le cas des faisceaux de modules).
\end{proof}






\section{Limites projectives et inductives de préfaisceaux}
\label{section-limits-presheaves}

\noindent
Soit $X$ un espace topologique.
Soit $\mathcal{I} \to \textit{PSh}(X)$, $i \mapsto \mathcal{F}_i$,
un diagramme.
\begin{enumerate}
\item Les objets $\lim_i \mathcal{F}_i$ et $\colim_i \mathcal{F}_i$
existent tous deux.
\item Pour tout ouvert $U \subset X$, on a
$$
(\lim_i \mathcal{F}_i)(U) =
\lim_i \mathcal{F}_i(U)
$$
et
$$
(\colim_i \mathcal{F}_i)(U) =
\colim_i \mathcal{F}_i(U).
$$
\item Soit $x \in X$ un point. En général, la fibre de
$\lim_i \mathcal{F}_i$ en $x$ n'est pas égale à
la limite projective des fibres. C'est toutefois le cas si la catégorie d'indices est finie.
Autrement dit, le foncteur fibre est
exact à gauche (voir Categories, Définition \ref{categories-definition-exact}).
\item Soit $x \in X$. On a toujours
$$
(\colim_i \mathcal{F}_i)_x =
\colim_i \mathcal{F}_{i, x}.
$$
\end{enumerate}
Toutes les démonstrations sont faciles.

\section{Limites projectives et inductives de faisceaux}
\label{section-limits-sheaves}

\noindent
Soit $X$ un espace topologique.
Soit $\mathcal{I} \to \Sh(X)$, $i \mapsto \mathcal{F}_i$,
un diagramme.
\begin{enumerate}
\item Les objets $\lim_i \mathcal{F}_i$ et $\colim_i \mathcal{F}_i$
existent tous deux.
\item Le foncteur d'inclusion $i : \Sh(X) \to \textit{PSh}(X)$
commute aux limites projectives. Autrement dit, on peut calculer la limite projective
dans la catégorie des faisceaux comme la limite projective dans la catégorie des
préfaisceaux. En particulier, pour tout ouvert $U \subset X$, on a
$$
(\lim_i \mathcal{F}_i)(U) = \lim_i \mathcal{F}_i(U).
$$
\item Le foncteur d'inclusion $i : \Sh(X) \to \textit{PSh}(X)$
ne commute pas en général aux limites inductives (pas même
aux limites inductives finies --- penser aux surjections). La limite inductive
s'obtient en prenant le faisceau associé à la limite inductive dans la
catégorie des préfaisceaux :
$$
\colim_i \mathcal{F}_i =
\Big(U \mapsto \colim_i \mathcal{F}_i(U)\Big)^\#.
$$
\item Soit $x \in X$ un point. En général, la fibre de
$\lim_i \mathcal{F}_i$ en $x$ n'est pas égale à
la limite projective des fibres. C'est toutefois le cas si la catégorie d'indices est finie.
Autrement dit, le foncteur fibre est
exact à gauche.
\item Soit $x \in X$. On a toujours
$$
(\colim_i \mathcal{F}_i)_x = \colim_i \mathcal{F}_{i, x}.
$$
\item Le foncteur faisceau associé
${}^\# : \textit{PSh}(X) \to \Sh(X)$ commute à toutes les
limites inductives et aux limites projectives finies. Il ne commute cependant pas
à toutes les limites projectives.
\end{enumerate}
Toutes les démonstrations sont faciles. Voici un exemple de ce qui vaut pour les limites
inductives filtrantes de faisceaux.

\begin{lemma}
\label{lemma-directed-colimits-sections}
Soit $X$ un espace topologique. Soit $I$ un ensemble ordonné filtrant.
Soit $(\mathcal{F}_i, \varphi_{ii'})$ un système inductif de faisceaux d'ensembles
indexé par $I$, voir
Categories, Section \ref{categories-section-posets-limits}.
Soit $U \subset X$ un ouvert.
Considérons l'application canonique
$$
\Psi :
\colim_i \mathcal{F}_i(U)
\longrightarrow
\left(\colim_i \mathcal{F}_i\right)(U)
$$
\begin{enumerate}
\item Si toutes les applications de transition sont injectives, alors
$\Psi$ est injective pour tout ouvert $U$.
\item Si $U$ est quasi-compact, alors $\Psi$ est injective.
\item Si $U$ est quasi-compact et si toutes les applications de transition sont injectives,
alors $\Psi$ est un isomorphisme.
\item Si $U$ admet un système cofinal de recouvrements ouverts
$\mathcal{U} : U = \bigcup_{j\in J} U_j$ with
$J$ étant fini, tels que $U_j \cap U_{j'}$ soit quasi-compact
pour tous $j, j' \in J$, alors $\Psi$ est bijective.
\end{enumerate}
\end{lemma}

\begin{proof}
Supposons toutes les applications de transition injectives. Dans ce cas, le préfaisceau
$\mathcal{F}' : V \mapsto \colim_i \mathcal{F}_i(V)$ est
séparé (voir la Définition \ref{definition-separated}).
D'après ce qui précède, on a
$(\mathcal{F}')^\# = \colim_i \mathcal{F}_i$.
Le Lemme \ref{lemma-separated-presheaf-into-sheaf} montre que
$\mathcal{F}' \to (\mathcal{F}')^\#$ est injective. Cela démontre (1).

\medskip\noindent
Supposons $U$ quasi-compact.
Supposons que $s \in \mathcal{F}_i(U)$ et
$s' \in \mathcal{F}_{i'}(U)$ déterminent au membre de gauche
des éléments ayant la même image par $\Psi$.
Comme $U$ est quasi-compact, cela signifie qu'il existe
un recouvrement ouvert fini $U = \bigcup_{j = 1, \ldots, m} U_j$
et, pour chaque $j$, un indice $i_j \in I$, avec $i_j \geq i$, $i_j \geq i'$,
tels que $\varphi_{ii_j}(s)$ et $\varphi_{i'i_j}(s')$ aient la
même restriction à $U_j$.
Soit $i''\in I$ supérieur ou égal ($\geq$) à chacun des $i_j$.
On en conclut que $\varphi_{ii''}(s)$ et $\varphi_{i'i''}(s')$
coïncident sur les ouverts $U_j$ pour tout $j$, et donc que
$\varphi_{ii''}(s) = \varphi_{i'i''}(s')$. Cela démontre (2).

\medskip\noindent
Supposons $U$ quasi-compact et toutes les applications de transition injectives.
Soit $s$ un élément du but de $\Psi$.
Comme $U$ est quasi-compact,
il existe un recouvrement ouvert fini $U = \bigcup_{j = 1, \ldots, m} U_j$,
ainsi que, pour chaque $j$, un indice $i_j \in I$ et un $s_j \in \mathcal{F}_{i_j}(U_j)$,
tels que $s|_{U_j}$ provienne de $s_j$ pour tout $j$.
Choisissons $i \in I$ supérieur ou égal ($\geq$) à chacun des $i_j$.
D'après (1), les sections $\varphi_{i_ji}(s_j)$ coïncident sur les
intersections $U_j \cap U_{j'}$. Elles se recollent donc en une section
$s' \in \mathcal{F}_i(U)$ qui s'envoie sur $s$ par $\Psi$.
Cela démontre (3).

\medskip\noindent
Supposons l'hypothèse de (4). En particulier,
$U$ est quasi-compact et (2) assure donc l'injectivité de $\Psi$.
Soit $s$ un élément du but de $\Psi$.
Par hypothèse, il existe un recouvrement ouvert fini
$U = \bigcup_{j = 1, \ldots, m} U_j$, où $U_j \cap U_{j'}$ est
quasi-compact pour tous $j, j' \in J$, ainsi que,
pour chaque $j$, un indice $i_j \in I$ et un $s_j \in \mathcal{F}_{i_j}(U_j)$
tels que $s|_{U_j}$ soit l'image de $s_j$ pour tout $j$.
Puisque $U_j \cap U_{j'}$ est quasi-compact, on peut appliquer (2)
et voir qu'il existe un $i_{jj'} \in I$, avec
$i_{jj'} \geq i_j$, $i_{jj'} \geq i_{j'}$, tel que
$\varphi_{i_ji_{jj'}}(s_j)$ et $\varphi_{i_{j'}i_{jj'}}(s_{j'})$
coïncident sur $U_j \cap U_{j'}$. Choisissons un indice $i \in I$
majorant tous les $i_{jj'}$. On voit alors que
les sections $\varphi_{i_ji}(s_j)$ de $\mathcal{F}_i$ se recollent
en une section de $\mathcal{F}_i$ sur $U$. Cette section s'envoie
sur l'élément $s$, comme voulu.
\end{proof}

\begin{example}
\label{example-conditions-needed-colimit}
Soit $X = \{s_1, s_2, \xi_1, \xi_2, \xi_3, \ldots\}$ comme ensemble.
Déclarons qu'une partie $U \subset X$ est ouverte
si $s_1 \in U$ ou $s_2 \in U$ entraîne que $U$ contient tous les
$\xi_i$. Soit $U_n = \{\xi_n, \xi_{n + 1}, \ldots\}$, et
soit $j_n : U_n \to X$ l'application d'inclusion.
Posons $\mathcal{F}_n = j_{n, *}\underline{\mathbf{Z}}$.
On a des applications de transition $\mathcal{F}_n \to \mathcal{F}_{n + 1}$.
Soit $\mathcal{F} = \colim \mathcal{F}_n$.
Remarquons que $\mathcal{F}_{n, \xi_m} = 0$ si $m < n$,
car $\{\xi_m\}$ est une partie ouverte de $X$ disjointe de $U_n$.
On voit donc que $\mathcal{F}_{\xi_n} = 0$ pour tout $n$.
En revanche, la fibre $\mathcal{F}_{s_i}$, pour $i = 1, 2$,
est la limite inductive
$$
M = \colim_n \prod\nolimits_{m \geq n} \mathbf{Z}
$$
qui n'est pas nulle. On en déduit que le faisceau
$\mathcal{F}$ est la somme directe des
faisceaux gratte-ciel de valeur $M$ aux points fermés
$s_1$ et $s_2$. Ainsi, $\Gamma(X, \mathcal{F}) = M \oplus M$.
D'autre part, le lecteur peut vérifier que
$\colim_n \Gamma(X, \mathcal{F}_n) = M$.
Une condition est donc nécessaire dans l'assertion (4) du
Lemme \ref{lemma-directed-colimits-sections} ci-dessus.
\end{example}

\noindent
Il existe une version du lemme précédent pour les faisceaux sur un diagramme
d'espaces spectraux. Introduisons les notations nécessaires à son énoncé.
Soit $\mathcal{I}$ une catégorie d'indices cofiltrante.
Soit $i \mapsto X_i$ un diagramme d'espaces spectraux indexé par $\mathcal{I}$,
tel que, pour $a : j \to i$ dans $\mathcal{I}$, l'application correspondante
$f_a : X_j \to X_i$ soit spectrale. Posons $X = \lim X_i$ et notons
$p_i : X \to X_i$ la projection.

\begin{lemma}
\label{lemma-compute-pullback-to-limit}
Dans la situation décrite ci-dessus, soit $i \in \Ob(\mathcal{I})$ et soit
$\mathcal{G}$ un faisceau sur $X_i$. Pour tout ouvert quasi-compact $U_i \subset X_i$,
on a
$$
p_i^{-1}\mathcal{G}(p_i^{-1}(U_i)) =
\colim_{a : j \to i} f_a^{-1}\mathcal{G}(f_a^{-1}(U_i))
$$
\end{lemma}

\begin{proof}
Montrons que l'application canonique
$\colim_{a : j \to i} f_a^{-1}\mathcal{G}(f_a^{-1}(U_i)) \to
p_i^{-1}\mathcal{G}(p_i^{-1}(U_i))$ est injective.
Soient $s, s'$ des sections de $f_a^{-1}\mathcal{G}$
sur $f_a^{-1}(U_i)$ pour un certain $a : j \to i$. Pour
$b : k \to j$, soit $Z_k \subset f_{a \circ b}^{-1}(U_i)$
la partie fermée formée des points $x$ tels que les images de
$s$ et $s'$ dans la fibre $(f_{a \circ b}^{-1}\mathcal{G})_x$
soient distinctes. Si $Z_k$ est non vide pour tout $b : k \to j$, alors le
Lemme \ref{topology-lemma-inverse-limit-spectral-spaces-nonempty} de Topology
montre que $\lim_{b : k \to j} Z_k$ est également non vide.
Pour $x \in \lim_{b : k \to j} Z_k \subset X$
(remarquons que $\mathcal{I}/j \to \mathcal{I}$ est initial),
les images de $s$ et $s'$ dans la fibre de
$p_i^{-1}\mathcal{G}$ en $x$ sont alors également distinctes, puisque
$(p_i^{-1}\mathcal{G})_x = (f_{b \circ a}^{-1}\mathcal{G})_{p_k(x)}$
pour tout $b : k \to j$ comme ci-dessus. Ainsi, si les images de $s$ et $s'$
dans $p_i^{-1}\mathcal{G}(p_i^{-1}(U_i))$ sont égales, alors $Z_k$
est vide pour un certain $b : k \to j$. Cela démontre l'injectivité.

\medskip\noindent
Surjectivité. Soit $s$ une section de $p_i^{-1}\mathcal{G}$ sur
$p_i^{-1}(U_i)$. D'après le
Lemme \ref{topology-lemma-directed-inverse-limit-spectral-spaces} de Topology,
la partie $p_i^{-1}(U_i)$ est un ouvert quasi-compact de l'espace spectral $X$.
Par construction du faisceau image inverse, on peut trouver un recouvrement ouvert
$p_i^{-1}(U_i) = \bigcup_{l \in L} W_l$,
des ouverts $V_{l, i} \subset X_i$ et
des sections $s_{l, i} \in \mathcal{G}(V_{l, i})$
tels que $p_i(W_l) \subset V_{l, i}$ et $p_i^{-1}s_{l, i}|_{W_l} = s|_{W_l}$.
Comme $X$ et $X_i$ sont spectraux et que $p_i^{-1}(U_i)$ est un ouvert quasi-compact,
on peut supposer $L$ fini et les $W_l$ et $V_{l, i}$ ouverts quasi-compacts
pour tout $l$. On peut alors appliquer le
Lemme \ref{topology-lemma-descend-opens} de Topology
pour trouver un $a : j \to i$ et un recouvrement ouvert
$f_a^{-1}(U_i) = \bigcup_{l \in L} W_{l, j}$ par des ouverts
quasi-compacts dont l'image inverse sur $X$ est le recouvrement
$p_i^{-1}(U_i) = \bigcup_{l \in L} W_l$
et tels qu'en outre $W_{l, j} \subset f_a^{-1}(V_{l, i})$.
Notons $s_{l, j}$ la restriction de l'image inverse de $s_{l, i}$
par $f_a$ à $W_{l, j}$. On voit alors que $s_{l, j}$ et $s_{l', j}$
se restreignent en des éléments de $(f_a^{-1}\mathcal{G})(W_{l, j} \cap W_{l', j})$
dont les images inverses donnent le même élément $(p_i^{-1}\mathcal{G})(W_l \cap W_{l'})$,
à savoir la restriction de $s$. L'injectivité permet donc de trouver
$b : k \to j$ tel que les sections $f_b^{-1}s_{l, j}$
se recollent en une section sur $f_{a \circ b}^{-1}(U_i)$, comme voulu.
\end{proof}

\noindent
Supposons maintenant donnés, en plus du système cofiltrant $X_i$ d'espaces spectraux,
les objets suivants :
\begin{enumerate}
\item un faisceau $\mathcal{F}_i$ sur $X_i$ pour tout $i \in \Ob(\mathcal{I})$ ;
\item pour $a : j \to i$, un $f_a$-morphisme
$\varphi_a : \mathcal{F}_i \to \mathcal{F}_j$
\end{enumerate}
tels que $\varphi_c = \varphi_b \circ \varphi_a$
dès que $c = a \circ b$. Posons $\mathcal{F} = \colim p_i^{-1}\mathcal{F}_i$
sur $X$.

\begin{lemma}
\label{lemma-descend-opens}
Dans la situation décrite ci-dessus, soit $i \in \Ob(\mathcal{I})$ et soit
$U_i \subset X_i$ un ouvert quasi-compact. Alors
$$
\colim_{a : j \to i} \mathcal{F}_j(f_a^{-1}(U_i)) = \mathcal{F}(p_i^{-1}(U_i))
$$
\end{lemma}

\begin{proof}
Rappelons que $p_i^{-1}(U_i)$ est un ouvert quasi-compact de l'espace spectral
$X$, voir le
Lemme \ref{topology-lemma-directed-inverse-limit-spectral-spaces} de Topology.
Le Lemme \ref{lemma-directed-colimits-sections} s'applique donc et donne
$$
\mathcal{F}(p_i^{-1}(U_i)) =
\colim_{a : j \to i} p_j^{-1}\mathcal{F}_j(p_i^{-1}(U_i)).
$$
Un argument formel montre que
$$
\colim_{a : j \to i} \mathcal{F}_j(f_a^{-1}(U_i)) =
\colim_{a : j \to i} \colim_{b : k \to j}
f_b^{-1}\mathcal{F}_j(f_{a \circ b}^{-1}(U_i))
$$
Il suffit donc de montrer que
$$
p_j^{-1}\mathcal{F}_j(p_i^{-1}(U_i)) =
\colim_{b : k \to j} f_b^{-1}\mathcal{F}_j(f_{a \circ b}^{-1}(U_i))
$$
C'est le Lemme \ref{lemma-compute-pullback-to-limit}
appliqué à $\mathcal{F}_j$ et à l'ouvert quasi-compact $f_a^{-1}(U_i)$.
\end{proof}









\section{Bases et faisceaux}
\label{section-bases}

\noindent
Il existe parfois une base de la topologie
formée d'ouverts avec lesquels il est plus facile de travailler
qu'avec des ouverts quelconques. Pour plus de commodité, nous donnons ici
quelques définitions et lemmes simples afin de
faciliter le travail avec les (pré)faisceaux dans une telle situation.

\begin{definition}
\label{definition-presheaf-basis}
Soit $X$ un espace topologique. Soit $\mathcal{B}$ une
base de la topologie de $X$.
\begin{enumerate}
\item Un {\it préfaisceau d'ensembles $\mathcal{F}$ sur $\mathcal{B}$}
est une règle qui associe à chaque $U \in \mathcal{B}$ un ensemble
$\mathcal{F}(U)$ et à chaque inclusion $V \subset U$
d'éléments de $\mathcal{B}$ une application
$\rho^U_V : \mathcal{F}(U) \to \mathcal{F}(V)$ telle que
$\rho^U_U = \text{id}_{\mathcal{F}(U)}$ pour tout $U \in \mathcal{B}$ et que,
si $W \subset V \subset U$ dans $\mathcal{B}$, on ait
$\rho^U_W = \rho^V_W \circ \rho ^U_V$.
\item Un {\it morphisme $\varphi : \mathcal{F} \to \mathcal{G}$
de préfaisceaux d'ensembles sur $\mathcal{B}$} est une règle qui associe à chaque
élément $U \in \mathcal{B}$ une application d'ensembles $\varphi : \mathcal{F}(U)
\to \mathcal{G}(U)$ compatible avec les applications de restriction.
\end{enumerate}
\end{definition}

\noindent
Comme pour les préfaisceaux usuels, nous employons la terminologie des sections,
des restrictions de sections, etc. En particulier, nous pouvons définir la
{\it fibre} de $\mathcal{F}$ en un point $x \in X$ comme la
limite inductive
$$
\mathcal{F}_x = \colim_{U\in \mathcal{B}, x\in U} \mathcal{F}(U).
$$
Comme pour la fibre d'un préfaisceau sur $X$, cette limite inductive est
filtrante. En effet, la famille des $U\in \mathcal{B}$,
$x \in U$, est un système fondamental de voisinages ouverts de $x$.

\medskip\noindent
Il est facile de construire des exemples montrant que la notion de préfaisceau
sur $X$ est très différente de celle de préfaisceau sur une base
de la topologie de $X$. Ce phénomène ne se produit pas pour les
faisceaux. La notion suivante est donc beaucoup plus utile.

\begin{definition}
\label{definition-sheaf-basis}
Soit $X$ un espace topologique. Soit $\mathcal{B}$ une
base de la topologie de $X$.
\begin{enumerate}
\item Un {\it faisceau d'ensembles $\mathcal{F}$ sur $\mathcal{B}$} est un préfaisceau
d'ensembles sur $\mathcal{B}$ qui satisfait la propriété supplémentaire
suivante : étant donnés $U \in \mathcal{B}$, un recouvrement
$U = \bigcup_{i \in I} U_i$ avec $U_i \in \mathcal{B}$, et
des recouvrements $U_i \cap U_j = \bigcup_{k \in I_{ij}} U_{ijk}$ avec
$U_{ijk} \in \mathcal{B}$, la condition de faisceau suivante est satisfaite :
\begin{itemize}
\item[(**)] Pour toute famille
de sections $s_i \in \mathcal{F}(U_i)$, $i \in I$, telle que
$\forall i, j\in I$, $\forall k\in I_{ij}$,
$$
s_i|_{U_{ijk}} = s_j|_{U_{ijk}}
$$
il existe une unique section $s \in \mathcal{F}(U)$ telle que
$s_i = s|_{U_i}$ pour tout $i \in I$.
\end{itemize}
\item Un {\it morphisme de faisceaux d'ensembles sur $\mathcal{B}$} est simplement un
morphisme de préfaisceaux d'ensembles.
\end{enumerate}
\end{definition}

\noindent
Nous expliquons d'abord qu'il suffit de
vérifier la condition de faisceau $(**)$
sur un système cofinal de recouvrements.
Dans la situation de la définition, supposons que
$U \in \mathcal{B}$. Notons temporairement
$\text{Cov}_\mathcal{B}(U)$ l'ensemble de tous les recouvrements
de $U$ par des éléments de $\mathcal{B}$.
Remarquons que $\text{Cov}_\mathcal{B}(U)$ est préordonné par le raffinement.
Une partie $C \subset \text{Cov}_\mathcal{B}(U)$ est un
système cofinal si, pour tout $\mathcal{U} \in \text{Cov}_\mathcal{B}(U)$,
il existe un recouvrement $\mathcal{V} \in C$ qui raffine $\mathcal{U}$.

\begin{lemma}
\label{lemma-cofinal-systems-coverings}
Conservons les notations ci-dessus.
Pour chaque $U \in \mathcal{B}$, soit $C(U) \subset \text{Cov}_\mathcal{B}(U)$
un système cofinal. Pour chaque $U \in \mathcal{B}$ et chaque
$\mathcal{U} : U = \bigcup U_i$ appartenant à $C(U)$, supposons donnés des recouvrements
$\mathcal{U}_{ij} : U_i \cap U_j = \bigcup U_{ijk}$,
$U_{ijk} \in \mathcal{B}$.
Soit $\mathcal{F}$ un préfaisceau d'ensembles sur $\mathcal{B}$.
Les conditions suivantes sont équivalentes :
\begin{enumerate}
\item Le préfaisceau $\mathcal{F}$ est un faisceau sur $\mathcal{B}$.
\item Pour tout $U \in \mathcal{B}$ et tout recouvrement
$\mathcal{U} : U = \bigcup U_i$ appartenant à $C(U)$, la condition de faisceau
$(**)$ est satisfaite (pour les recouvrements $\mathcal{U}_{ij}$ donnés).
\end{enumerate}
\end{lemma}

\begin{proof}
Nous devons montrer que (2) implique (1).
Supposons que $U \in \mathcal{B}$ et que
$\mathcal{U} : U = \bigcup_{i\in I} U_i$ soit un recouvrement quelconque
par des éléments de $\mathcal{B}$. Comme le système $C(U)$ est cofinal,
nous pouvons trouver un élément $\mathcal{V} : U = \bigcup_{j \in J} V_j$
de $C(U)$ qui raffine $\mathcal{U}$. Cela signifie qu'il existe
une application $\alpha : J \to I$ telle que $V_j \subset U_{\alpha(j)}$.

\medskip\noindent
Remarquons que, si $s, s' \in \mathcal{F}(U)$ sont des sections telles
que $s|_{U_i} = s'|_{U_i}$, alors
$$
s|_{V_j}
= (s|_{U_{\alpha(j)}})|_{V_j}
= (s'|_{U_{\alpha(j)}})|_{V_j}
= s'|_{V_j}
$$
pour tout $j$. L'unicité dans $(**)$
pour le recouvrement $\mathcal{V}$ entraîne donc que $s = s'$.
Nous avons ainsi démontré la partie unicité de $(**)$
pour notre recouvrement quelconque $\mathcal{U}$.

\medskip\noindent
Supposons en outre que $U_i \cap U_{i'} = \bigcup_{k \in I_{ii'}} U_{ii'k}$
soient des recouvrements quelconques par des $U_{ii'k} \in \mathcal{B}$.
Essayons de démontrer la partie existence de $(**)$ pour le système
$(\mathcal{U}, \mathcal{U}_{ij})$. Soient donc $s_i \in \mathcal{F}(U_i)$
et supposons que
$$
s_i|_{U_{ii'k}} = s_{i'}|_{U_{ii'k}}
$$
pour tous $i, i', k$. Posons $t_j = s_{\alpha(j)}|_{V_j}$, où $\mathcal{V}$
et $\alpha$ sont comme ci-dessus.

\medskip\noindent
L'argument présente ici une petite difficulté. Soit en effet
$\mathcal{V}_{jj'} : V_j \cap V_{j'} = \bigcup_{l \in J_{jj'}} V_{jj'l}$
le recouvrement fourni par l'énoncé du lemme.
Il n'est pas clair a priori que
$$
t_j|_{V_{jj'l}} = t_{j'}|_{V_{jj'l}}
$$
pour tous $j, j', l$. Pour le voir, remarquons que nous avons bien
$$
t_j|_W = t_{j'}|_W \text{ pour tout } W \in \mathcal{B},
W \subset V_{jj'l} \cap U_{\alpha(j)\alpha(j')k}
$$
pour tout $k \in I_{\alpha(j)\alpha(j')}$, par notre hypothèse sur
la famille d'éléments $s_i$. Et puisque
$V_j \cap V_{j'} \subset U_{\alpha(j)} \cap U_{\alpha(j')}$,
nous voyons que $t_j|_{V_{jj'l}}$ et $t_{j'}|_{V_{jj'l}}$
coïncident sur les membres d'un recouvrement de $V_{jj'l}$ par des
éléments de $\mathcal{B}$. La partie unicité démontrée ci-dessus
donne donc enfin l'égalité voulue entre
$t_j|_{V_{jj'l}}$ et $t_{j'}|_{V_{jj'l}}$.
Nous obtenons alors l'existence d'un élément $t \in \mathcal{F}(U)$
par la propriété $(**)$ pour $(\mathcal{V}, \mathcal{V}_{jj'})$.

\medskip\noindent
Une petite difficulté subsiste. Nous savons que $t$ se restreint en $t_j$ sur $V_j$,
mais nous ne savons pas encore que $t$ se restreint en $s_i$ sur $U_i$. Pour conclure,
remarquons que les ensembles $U_i \cap V_j$, $j \in J$, recouvrent $U_i$. Par conséquent,
les ensembles $U_{i \alpha(j) k} \cap V_j$, $j\in J$, $k \in I_{i\alpha(j)}$,
recouvrent eux aussi $U_i$. Nous laissons au lecteur le soin de voir que $t$ et $s_i$ induisent
la même section de $\mathcal{F}$ sur tout $W \in \mathcal{B}$
contenu dans l'un des ouverts
$U_{i \alpha(j) k} \cap V_j$, $j\in J$, $k \in I_{i\alpha(j)}$.
La partie unicité établie ci-dessus permet donc de conclure.
\end{proof}

\begin{lemma}
\label{lemma-cofinal-systems-coverings-standard-case}
Soit $X$ un espace topologique. Soit $\mathcal{B}$ une base de la
topologie de $X$. Supposons que, pour tout triplet $U, U', U'' \in \mathcal{B}$
tel que $U' \subset U$ et $U'' \subset U$, on ait $U' \cap U'' \in \mathcal{B}$.
Pour chaque $U \in \mathcal{B}$, soit $C(U) \subset \text{Cov}_\mathcal{B}(U)$
un système cofinal.
Soit $\mathcal{F}$ un préfaisceau d'ensembles sur $\mathcal{B}$.
Les conditions suivantes sont équivalentes :
\begin{enumerate}
\item Le préfaisceau $\mathcal{F}$ est un faisceau sur $\mathcal{B}$.
\item Pour tout $U \in \mathcal{B}$, tout recouvrement
$\mathcal{U} : U = \bigcup U_i$ appartenant à $C(U)$ et toute
famille de sections $s_i \in \mathcal{F}(U_i)$ telle
que $s_i|_{U_i \cap U_j} = s_j|_{U_i \cap U_j}$, il
existe une unique section $s \in \mathcal{F}(U)$ qui
se restreint en $s_i$ sur $U_i$.
\end{enumerate}
\end{lemma}

\begin{proof}
C'est une reformulation du
Lemme \ref{lemma-cofinal-systems-coverings} ci-dessus
dans le cas particulier où chacun des recouvrements $\mathcal{U}_{ij}$
est formé d'un seul élément. Ce cas est d'ailleurs beaucoup
plus simple et constitue un exercice facile que l'on peut résoudre directement.
\end{proof}

\begin{lemma}
\label{lemma-condition-star-sections}
Soit $X$ un espace topologique.
Soit $\mathcal{B}$ une base de la topologie de $X$.
Soit $U \in \mathcal{B}$.
Soit $\mathcal{F}$ un faisceau d'ensembles sur $\mathcal{B}$.
L'application
$$
\mathcal{F}(U) \to \prod\nolimits_{x \in U} \mathcal{F}_x
$$
identifie $\mathcal{F}(U)$ aux éléments $(s_x)_{x\in U}$
qui vérifient la propriété
\begin{itemize}
\item[(*)] Pour tout $x \in U$, il existe $V \in \mathcal{B}$
tel que $x \in V \subset U$, ainsi qu'une section $\sigma \in \mathcal{F}(V)$,
de sorte que, pour tout $y \in V$, on ait $s_y = (V, \sigma)$ dans $\mathcal{F}_y$.
\end{itemize}
\end{lemma}

\begin{proof}
Remarquons d'abord que l'application
$\mathcal{F}(U) \to \prod\nolimits_{x \in U} \mathcal{F}_x$
est injective par l'unicité dans la condition de faisceau
de la Définition \ref{definition-sheaf-basis}. Soit $(s_x)$
un élément quelconque du membre de droite qui vérifie $(*)$.
Cela signifie clairement que l'on peut trouver un recouvrement $U = \bigcup U_i$,
$U_i \in \mathcal{B}$, tel que $(s_x)_{x \in U_i}$ provienne
de certains $\sigma_i \in \mathcal{F}(U_i)$. Pour tout $y \in U_i \cap U_j$,
les sections $\sigma_i$ et $\sigma_j$ coïncident dans la fibre
$\mathcal{F}_y$. Il existe donc un élément $V_{ijy} \in \mathcal{B}$,
$y \in V_{ijy}$, tel que $\sigma_i|_{V_{ijy}} = \sigma_j|_{V_{ijy}}$.
La condition de faisceau $(**)$ de la Définition \ref{definition-sheaf-basis}
s'applique donc au système des $\sigma_i$, et nous obtenons une section
$s \in \mathcal{F}(U)$ qui possède la propriété voulue.
\end{proof}

\noindent
Soit $X$ un espace topologique.
Soit $\mathcal{B}$ une base de la topologie de $X$.
Il existe un foncteur de restriction naturel de la catégorie
des faisceaux d'ensembles sur $X$ vers la catégorie des faisceaux
d'ensembles sur $\mathcal{B}$. Ce foncteur est en fait une équivalence
de catégories. Concrètement, cela signifie ce qui suit.

\begin{lemma}
\label{lemma-extend-off-basis}
Soit $X$ un espace topologique.
Soit $\mathcal{B}$ une base de la topologie de $X$.
Soit $\mathcal{F}$ un faisceau d'ensembles sur $\mathcal{B}$.
Il existe un unique faisceau d'ensembles $\mathcal{F}^{ext}$
sur $X$ tel que $\mathcal{F}^{ext}(U) = \mathcal{F}(U)$
pour tout $U \in \mathcal{B}$, de manière compatible avec les applications
de restriction.
\end{lemma}

\begin{proof}
Nous construisons d'abord un préfaisceau $\mathcal{F}^{ext}$ possédant la
propriété voulue. Pour un ouvert quelconque $U \subset X$, nous
définissons $\mathcal{F}^{ext}(U)$ comme l'ensemble des éléments
$(s_x)_{x \in U}$ qui vérifient la condition $(*)$ du
Lemme \ref{lemma-condition-star-sections}.
Il est clair qu'il existe des applications de restriction
qui font de $\mathcal{F}^{ext}$ un préfaisceau d'ensembles.
En outre, le Lemme \ref{lemma-condition-star-sections} montre que
$\mathcal{F}(U) = \mathcal{F}^{ext}(U)$
dès que $U$ est un élément de la base $\mathcal{B}$.
Pour voir que $\mathcal{F}^{ext}$ est un faisceau, on peut
raisonner comme dans la preuve du Lemme \ref{lemma-sheafification-sheaf}.
\end{proof}

\noindent
Remarquons que l'on a
$$
\mathcal{F}_x = \mathcal{F}_x^{ext}
$$
dans la situation du lemme. En effet, la
famille des éléments de $\mathcal{B}$ qui contiennent
$x$ forme un système fondamental de voisinages ouverts de $x$.

\begin{lemma}
\label{lemma-restrict-basis-equivalence}
Soit $X$ un espace topologique.
Soit $\mathcal{B}$ une base de la topologie de $X$.
Notons $\Sh(\mathcal{B})$ la catégorie des
faisceaux sur $\mathcal{B}$.
Il existe une équivalence de catégories
$$
\Sh(X) \longrightarrow \Sh(\mathcal{B})
$$
qui associe à un faisceau sur $X$ sa restriction aux
membres de $\mathcal{B}$.
\end{lemma}

\begin{proof}
Le foncteur inverse est donné par le Lemme \ref{lemma-extend-off-basis} ci-dessus.
La vérification des fonctorialités évidentes est laissée au
lecteur.
\end{proof}

\noindent
Ceci termine l'étude des faisceaux d'ensembles sur
une base $\mathcal{B}$. Soit $(\mathcal{C}, F)$ une
espèce de structure algébrique. Pour terminer cette section,
nous souhaitons signaler que les constructions
ci-dessus valent pour les faisceaux à valeurs dans $\mathcal{C}$.
Définissons brièvement les notions pertinentes.

\begin{definition}
\label{definition-sheaf-structures-basis}
Soit $X$ un espace topologique. Soit $\mathcal{B}$ une
base de la topologie de $X$. Soit $(\mathcal{C}, F)$
une espèce de structure algébrique.
\begin{enumerate}
\item Un {\it préfaisceau $\mathcal{F}$ à valeurs dans $\mathcal{C}$
sur $\mathcal{B}$} est une règle qui associe à chaque
$U \in \mathcal{B}$ un objet
$\mathcal{F}(U)$ de $\mathcal{C}$ et à chaque inclusion $V \subset U$
d'éléments de $\mathcal{B}$ un morphisme
$\rho^U_V : \mathcal{F}(U) \to \mathcal{F}(V)$ dans $\mathcal{C}$ tel que
$\rho^U_U = \text{id}_{\mathcal{F}(U)}$ pour tout $U \in \mathcal{B}$ et que,
si $W \subset V \subset U$ dans $\mathcal{B}$, on ait
$\rho^U_W = \rho^V_W \circ \rho ^U_V$.
\item Un {\it morphisme $\varphi : \mathcal{F} \to \mathcal{G}$
de préfaisceaux à valeurs dans $\mathcal{C}$
sur $\mathcal{B}$} est une règle qui associe à chaque
élément $U \in \mathcal{B}$ un morphisme de
structures algébriques $\varphi : \mathcal{F}(U) \to \mathcal{G}(U)$
compatible avec les applications de restriction.
\item Étant donné un préfaisceau $\mathcal{F}$ à valeurs dans $\mathcal{C}$
sur $\mathcal{B}$, nous appelons $U \mapsto F(\mathcal{F}(U))$ le
préfaisceau d'ensembles sous-jacent.
\item Un {\it faisceau $\mathcal{F}$ à valeurs dans $\mathcal{C}$
sur $\mathcal{B}$} est un préfaisceau à valeurs dans $\mathcal{C}$
sur $\mathcal{B}$ dont le préfaisceau d'ensembles sous-jacent est un faisceau.
\end{enumerate}
\end{definition}

\noindent
Nous pouvons maintenant définir la {\it fibre} en $x \in X$
d'un préfaisceau à valeurs dans $\mathcal{C}$ sur $\mathcal{B}$
comme la limite inductive filtrante
$$
\mathcal{F}_x = \colim_{U\in \mathcal{B}, x\in U} \mathcal{F}(U).
$$
Elle existe comme objet de $\mathcal{C}$
en vertu de nos hypothèses sur $\mathcal{C}$.
De plus, l'ensemble sous-jacent à $\mathcal{F}_x$
est la fibre du préfaisceau d'ensembles sous-jacent sur $\mathcal{B}$.

\medskip\noindent
Remarquons que les Lemmes \ref{lemma-cofinal-systems-coverings},
\ref{lemma-cofinal-systems-coverings-standard-case} et
\ref{lemma-condition-star-sections} portent sur la propriété de faisceau,
que nous avons définie en termes du préfaisceau
d'ensembles associé. Ils se généralisent donc sans changement à la notion
de préfaisceau à valeurs dans $\mathcal{C}$. L'analogue du
Lemme \ref{lemma-extend-off-basis} demande quelques précautions. Le voici.

\begin{lemma}
\label{lemma-extend-off-basis-structures}
Soit $X$ un espace topologique. Soit $(\mathcal{C}, F)$
une espèce de structure algébrique.
Soit $\mathcal{B}$ une base de la topologie de $X$.
Soit $\mathcal{F}$ un faisceau à valeurs dans $\mathcal{C}$
sur $\mathcal{B}$.
Il existe un unique faisceau $\mathcal{F}^{ext}$ à valeurs dans $\mathcal{C}$
sur $X$ tel que $\mathcal{F}^{ext}(U) = \mathcal{F}(U)$
pour tout $U \in \mathcal{B}$, de manière compatible avec les applications
de restriction.
\end{lemma}

\begin{proof}
En vertu des conditions imposées au couple $(\mathcal{C}, F)$, il
suffit de construire un préfaisceau $\mathcal{F}^{ext}$
qui possède la propriété voulue au niveau des préfaisceaux
d'ensembles sous-jacents. Notre première tâche consiste donc à construire
un objet convenable $\mathcal{F}^{ext}(U)$ pour tout ouvert $U \subset X$.
Nous pourrions le faire en imitant le
Lemme \ref{lemma-diagram-fibre-product} dans le cadre
des préfaisceaux sur $\mathcal{B}$. Voici toutefois une méthode légèrement différente
(mais essentiellement équivalente) :
définissons-le comme la limite inductive filtrante
$$
\mathcal{F}^{ext}(U)
:=
\colim_\mathcal{U} FIB(\mathcal{U})
$$
sur tous les recouvrements
$\mathcal{U} : U = \bigcup_{i\in I} U_i$ par des $U_i \in \mathcal{B}$
du produit fibré
$$
\xymatrix{
FIB(\mathcal{U}) \ar[r] \ar[d] &
\prod\nolimits_{x\in U} \mathcal{F}_x \ar[d] \\
\prod\nolimits_{i\in I} \mathcal{F}(U_i) \ar[r] &
\prod\nolimits_{i \in I} \prod\nolimits_{x\in U_i} \mathcal{F}_x
}
$$
Par les arguments usuels, voir le Lemme \ref{lemma-image-contained-in}
et l'Exemple \ref{example-application-lemma-image-contained-in},
il suffit de montrer que, sur les ensembles sous-jacents, cette construction
coïncide avec la définition ci-dessus au moyen de $(**)$.
Les détails sont laissés au lecteur.
\end{proof}

\noindent
Remarquons que l'on a
$$
\mathcal{F}_x = \mathcal{F}_x^{ext}
$$
comme objets de $\mathcal{C}$
dans la situation du lemme. En effet, la
famille des éléments de $\mathcal{B}$ qui contiennent
$x$ forme un système fondamental de voisinages ouverts de $x$.

\begin{lemma}
\label{lemma-restrict-basis-equivalence-structures}
Soit $X$ un espace topologique.
Soit $\mathcal{B}$ une base de la topologie de $X$.
Soit $(\mathcal{C}, F)$ une espèce de structure algébrique.
Notons $\Sh(\mathcal{B}, \mathcal{C})$ la catégorie des
faisceaux à valeurs dans $\mathcal{C}$ sur $\mathcal{B}$.
Il existe une équivalence de catégories
$$
\Sh(X, \mathcal{C})
\longrightarrow
\Sh(\mathcal{B}, \mathcal{C})
$$
qui associe à un faisceau sur $X$ sa restriction aux
membres de $\mathcal{B}$.
\end{lemma}

\begin{proof}
Le foncteur inverse est donné par le
Lemme \ref{lemma-extend-off-basis-structures} ci-dessus.
La vérification des fonctorialités évidentes est laissée au
lecteur.
\end{proof}

\noindent
Enfin, nous abordons le cas des (pré)faisceaux de modules
sur une base. Nous utiliserons le fait élémentaire que la catégorie
des préfaisceaux d'ensembles sur une base admet des produits et que
ceux-ci s'obtiennent en prenant les produits des valeurs sur les
éléments de la base.

\begin{definition}
\label{definition-sheaf-modules-basis}
Soit $X$ un espace topologique. Soit $\mathcal{B}$ une
base de la topologie de $X$. Soit $\mathcal{O}$
un préfaisceau d'anneaux sur $\mathcal{B}$.
\begin{enumerate}
\item Un {\it préfaisceau de $\mathcal{O}$-modules $\mathcal{F}$
sur $\mathcal{B}$} est un préfaisceau de groupes abéliens sur
$\mathcal{B}$ muni d'un morphisme de préfaisceaux
d'ensembles $\mathcal{O} \times \mathcal{F} \to \mathcal{F}$
tel que, pour tout $U \in \mathcal{B}$, l'application
$\mathcal{O}(U) \times \mathcal{F}(U) \to \mathcal{F}(U)$
munisse le groupe $\mathcal{F}(U)$ d'une structure de $\mathcal{O}(U)$-module.
\item Un {\it morphisme $\varphi : \mathcal{F} \to \mathcal{G}$
de préfaisceaux de $\mathcal{O}$-modules sur $\mathcal{B}$}
est un morphisme de préfaisceaux abéliens sur $\mathcal{B}$
qui induit un homomorphisme de $\mathcal{O}(U)$-modules
$\mathcal{F}(U) \to \mathcal{G}(U)$ pour tout $U \in \mathcal{B}$.
\item Supposons que $\mathcal{O}$ soit un faisceau d'anneaux
sur $\mathcal{B}$. Un {\it faisceau $\mathcal{F}$ de $\mathcal{O}$-modules
sur $\mathcal{B}$} est un préfaisceau de $\mathcal{O}$-modules
sur $\mathcal{B}$ dont le préfaisceau de groupes abéliens sous-jacent
est un faisceau.
\end{enumerate}
\end{definition}

\noindent
Nous pouvons définir la {\it fibre} en $x \in X$
d'un préfaisceau de $\mathcal{O}$-modules sur $\mathcal{B}$
comme la limite inductive filtrante
$$
\mathcal{F}_x = \colim_{U\in \mathcal{B}, x\in U} \mathcal{F}(U).
$$
C'est un $\mathcal{O}_x$-module.

\medskip\noindent
Remarquons que les Lemmes \ref{lemma-cofinal-systems-coverings},
\ref{lemma-cofinal-systems-coverings-standard-case} et
\ref{lemma-condition-star-sections} portent sur la propriété de faisceau,
que nous avons définie en termes du préfaisceau
d'ensembles associé. Ils se généralisent donc sans changement à la notion
de préfaisceau de $\mathcal{O}$-modules. L'analogue du
Lemme \ref{lemma-extend-off-basis} est le suivant.

\begin{lemma}
\label{lemma-extend-off-basis-module}
Soit $X$ un espace topologique.
Soit $\mathcal{B}$ une base de la topologie de $X$.
Soit $\mathcal{O}$ un faisceau d'anneaux sur $\mathcal{B}$.
Soit $\mathcal{F}$ un faisceau de $\mathcal{O}$-modules
sur $\mathcal{B}$. Soit $\mathcal{O}^{ext}$ le faisceau
d'anneaux sur $X$ qui prolonge $\mathcal{O}$, et soit
$\mathcal{F}^{ext}$ le faisceau abélien sur $X$ qui prolonge
$\mathcal{F}$ ; voir le Lemme \ref{lemma-extend-off-basis-structures}.
Il existe un morphisme canonique
$$
\mathcal{O}^{ext} \times \mathcal{F}^{ext}
\longrightarrow
\mathcal{F}^{ext}
$$
qui coïncide avec le morphisme donné sur les éléments de $\mathcal{B}$
et qui munit $\mathcal{F}^{ext}$ d'une structure
de $\mathcal{O}^{ext}$-module.
\end{lemma}

\begin{proof}
Il suffit de construire le morphisme de multiplication
au niveau des préfaisceaux d'ensembles. La manière la plus simple
de le voir est peut-être de démontrer directement que, si
$(f_x)_{x \in U}$, $f_x \in \mathcal{O}_x$
et
$(m_x)_{x \in U}$, $m_x \in \mathcal{F}_x$
vérifient $(*)$, alors l'élément
$(f_xm_x)_{x \in U}$ vérifie lui aussi $(*)$.
On obtient alors le résultat voulu, car, dans la preuve
du Lemme \ref{lemma-extend-off-basis}, nous construisons le prolongement
à l'aide de familles d'éléments de fibres qui vérifient $(*)$.
\end{proof}

\noindent
Remarquons que l'on a
$$
\mathcal{F}_x = \mathcal{F}_x^{ext}
$$
comme $\mathcal{O}_x$-modules dans la situation du lemme.
Cela tient au fait que la famille des éléments de $\mathcal{B}$ qui contiennent
$x$ forme un système fondamental de voisinages ouverts de $x$, ou
simplement au fait que l'assertion est vraie sur les ensembles sous-jacents.

\begin{lemma}
\label{lemma-restrict-basis-equivalence-modules}
Soit $X$ un espace topologique.
Soit $\mathcal{B}$ une base de la topologie de $X$.
Soit $\mathcal{O}$ un faisceau d'anneaux sur $X$.
Notons $\textit{Mod}(\mathcal{O}|_\mathcal{B})$ la catégorie des
faisceaux de $\mathcal{O}|_\mathcal{B}$-modules sur $\mathcal{B}$.
Il existe une équivalence de catégories
$$
\textit{Mod}(\mathcal{O})
\longrightarrow
\textit{Mod}(\mathcal{O}|_\mathcal{B})
$$
qui associe à un faisceau de $\mathcal{O}$-modules sur $X$ sa restriction aux
membres de $\mathcal{B}$.
\end{lemma}

\begin{proof}
Le foncteur inverse est donné par le
Lemme \ref{lemma-extend-off-basis-module} ci-dessus.
La vérification des fonctorialités évidentes est laissée au lecteur.
\end{proof}

\noindent
Enfin, nous étudions le lien de cette théorie avec les
applications continues. C'est maintenant très facile grâce au travail
accompli ci-dessus. Traitons d'abord le cas où une base de l'espace
d'arrivée est donnée.

\begin{lemma}
\label{lemma-f-map-basis-below-structures}
Soit $f : X \to Y$ une application continue d'espaces topologiques.
Soit $(\mathcal{C}, F)$ une espèce de structure algébrique.
Soit $\mathcal{F}$ un faisceau à valeurs dans $\mathcal{C}$ sur $X$.
Soit $\mathcal{G}$ un faisceau à valeurs dans $\mathcal{C}$ sur $Y$.
Soit $\mathcal{B}$ une base de la topologie de $Y$.
Supposons donné, pour tout $V \in \mathcal{B}$, un morphisme
$$
\varphi_V :
\mathcal{G}(V)
\longrightarrow
\mathcal{F}(f^{-1}V)
$$
de $\mathcal{C}$ compatible avec les applications de restriction.
Il existe alors un unique $f$-morphisme (voir la Définition \ref{definition-f-map}
et la discussion des $f$-morphismes dans la
Section \ref{section-presheaves-structures-functorial})
$\varphi : \mathcal{G} \to \mathcal{F}$
qui redonne $\varphi_V$ pour $V \in \mathcal{B}$.
\end{lemma}

\begin{proof}
C'est immédiat, car la famille de morphismes
revient à un morphisme entre les restrictions
de $\mathcal{G}$ et de $f_*\mathcal{F}$ à $\mathcal{B}$.
Par le Lemme \ref{lemma-restrict-basis-equivalence-structures},
cela revient à donner un morphisme de $\mathcal{G}$
vers $f_*\mathcal{F}$, ce qui, par le Lemme \ref{lemma-f-map},
équivaut à un $f$-morphisme. Voir aussi le
Lemme \ref{lemma-f-map-sets-algebraic-structures}
et la discussion qui le précède
pour le traitement du cas des faisceaux de structures algébriques.
\end{proof}

\noindent
Voici l'analogue pour les espaces annelés.

\begin{lemma}
\label{lemma-f-map-basis-below-modules}
Soit $(f, f^\sharp) : (X, \mathcal{O}_X) \to (Y, \mathcal{O}_Y)$
un morphisme d'espaces annelés.
Soit $\mathcal{F}$ un faisceau de $\mathcal{O}_X$-modules.
Soit $\mathcal{G}$ un faisceau de $\mathcal{O}_Y$-modules.
Soit $\mathcal{B}$ une base de la topologie de $Y$.
Supposons donné, pour tout $V \in \mathcal{B}$, un
morphisme de $\mathcal{O}_Y(V)$-modules
$$
\varphi_V :
\mathcal{G}(V)
\longrightarrow
\mathcal{F}(f^{-1}V)
$$
(où $\mathcal{F}(f^{-1}V)$ est muni d'une structure de module par restriction des scalaires le long de
$f^\sharp_V : \mathcal{O}_Y(V) \to \mathcal{O}_X(f^{-1}V)$)
compatible avec les applications de restriction.
Il existe alors un unique $f$-morphisme (voir la discussion des $f$-morphismes dans la
Section \ref{section-ringed-spaces-functoriality-modules})
$\varphi : \mathcal{G} \to \mathcal{F}$
qui redonne $\varphi_V$ pour $V \in \mathcal{B}$.
\end{lemma}

\begin{proof}
La preuve est la même que celle du lemme correspondant
pour les faisceaux de structures algébriques ci-dessus.
\end{proof}

\begin{lemma}
\label{lemma-f-map-basis-above-and-below-structures}
Soit $f : X \to Y$ une application continue d'espaces topologiques.
Soit $(\mathcal{C}, F)$ une espèce de structure algébrique.
Soit $\mathcal{F}$ un faisceau à valeurs dans $\mathcal{C}$ sur $X$.
Soit $\mathcal{G}$ un faisceau à valeurs dans $\mathcal{C}$ sur $Y$.
Soit $\mathcal{B}_Y$ une base de la topologie de $Y$.
Soit $\mathcal{B}_X$ une base de la topologie de $X$.
Supposons donné, pour tous $V \in \mathcal{B}_Y$ et
$U \in \mathcal{B}_X$ tels que $f(U) \subset V$, un morphisme
$$
\varphi_V^U :
\mathcal{G}(V)
\longrightarrow
\mathcal{F}(U)
$$
de $\mathcal{C}$ compatible avec les applications de restriction.
Il existe alors un unique $f$-morphisme (voir
la Définition \ref{definition-f-map} et la discussion
des $f$-morphismes dans la Section \ref{section-presheaves-structures-functorial})
$\varphi : \mathcal{G} \to \mathcal{F}$
qui redonne $\varphi_V^U$ comme la composée
$$
\mathcal{G}(V) \xrightarrow{\varphi_V}
\mathcal{F}(f^{-1}(V)) \xrightarrow{\text{rest.}}
\mathcal{F}(U)
$$
pour toute paire $(U, V)$ comme ci-dessus.
\end{lemma}

\begin{proof}
Démontrons d'abord ce résultat pour les faisceaux d'ensembles.
Fixons un ouvert $V \subset Y$. Choisissons $s \in \mathcal{G}(V)$.
Nous allons construire un élément
$\varphi_V(s) \in \mathcal{F}(f^{-1}V)$.
Nous pouvons définir une valeur $\varphi(s)_x$ dans la fibre $\mathcal{F}_x$
pour tout $x \in f^{-1}V$ en choisissant un $U \in \mathcal{B}_X$
tel que $x \in U \subset f^{-1}V$ et en prenant pour $\varphi(s)_x$
la classe d'équivalence de $(U, \varphi_V^U(s))$
dans la fibre. Il est clair que la famille $(\varphi(s)_x)_{x \in f^{-1}V}$
vérifie la condition $(*)$, car les morphismes $\varphi_V^U$,
lorsque $U$ varie,
sont compatibles avec les restrictions dans le faisceau $\mathcal{F}$.
Ainsi, par la preuve du Lemme \ref{lemma-extend-off-basis},
nous voyons que $(\varphi(s)_x)_{x \in f^{-1}V}$ correspond
à un unique élément $\varphi_V(s)$ de $\mathcal{F}(f^{-1}V)$.
Nous avons donc défini une application d'ensembles
$\varphi_V : \mathcal{G}(V) \to \mathcal{F}(f^{-1}V)$.
La compatibilité entre $\varphi_V$ et $\varphi_V^U$
résulte du Lemme \ref{lemma-condition-star-sections}.

\medskip\noindent
Nous laissons au lecteur le soin de montrer que la construction
de $\varphi_V$ est compatible avec les applications de restriction lorsque
$V \in \mathcal{B}_Y$ varie. Nous pouvons donc appliquer le Lemme
\ref{lemma-f-map-basis-below-structures} ci-dessus pour les
``recoller'' et obtenir le $f$-morphisme voulu.

\medskip\noindent
Enfin, remarquons que le morphisme de faisceaux d'ensembles ainsi construit
possède la propriété suivante : le morphisme sur les fibres
$$
\mathcal{G}_{f(x)} \longrightarrow \mathcal{F}_x
$$
est la limite inductive du système de morphismes $\varphi_V^U$ lorsque
$V \in \mathcal{B}_Y$ parcourt les éléments qui
contiennent $f(x)$ et $U \in \mathcal{B}_X$ parcourt ceux qui
contiennent $x$. En particulier, si $\mathcal{G}$ et $\mathcal{F}$
sont les faisceaux d'ensembles sous-jacents à des faisceaux de structures algébriques,
nous voyons que le morphisme sur les fibres est un morphisme de structures
algébriques. Le morphisme correspondant de
faisceaux d'ensembles sous-jacents $f^{-1}\mathcal{G} \to \mathcal{F}$
satisfait donc les hypothèses du
Lemme \ref{lemma-f-map-sets-algebraic-structures}.
Nous en concluons que $f^{-1}\mathcal{G} \to \mathcal{F}$
est un morphisme de faisceaux à valeurs dans $\mathcal{C}$.
Par adjonction, cela signifie que $\varphi$ est
un $f$-morphisme de faisceaux de structures algébriques.
\end{proof}

\begin{lemma}
\label{lemma-f-map-basis-above-and-below-modules}
Soit $(f, f^\sharp) : (X, \mathcal{O}_X) \to (Y, \mathcal{O}_Y)$
un morphisme d'espaces annelés.
Soit $\mathcal{F}$ un faisceau de $\mathcal{O}_X$-modules.
Soit $\mathcal{G}$ un faisceau de $\mathcal{O}_Y$-modules.
Soit $\mathcal{B}_Y$ une base de la topologie de $Y$.
Soit $\mathcal{B}_X$ une base de la topologie de $X$.
Supposons donné, pour tous $V \in \mathcal{B}_Y$ et
$U \in \mathcal{B}_X$ tels que $f(U) \subset V$, un
morphisme de $\mathcal{O}_Y(V)$-modules
$$
\varphi_V^U :
\mathcal{G}(V)
\longrightarrow
\mathcal{F}(U)
$$
compatible avec les applications de restriction. Ici, la
structure de $\mathcal{O}_Y(V)$-module sur $\mathcal{F}(U)$
s'obtient par restriction des scalaires à partir de la structure de $\mathcal{O}_X(U)$-module
le long du morphisme $f^\sharp_V : \mathcal{O}_Y(V)
\to \mathcal{O}_X(f^{-1}V) \to \mathcal{O}_X(U)$.
Il existe alors un unique $f$-morphisme de faisceaux de modules (voir
la Définition \ref{definition-f-map} et la discussion
des $f$-morphismes dans la Section \ref{section-ringed-spaces-functoriality-modules})
$\varphi : \mathcal{G} \to \mathcal{F}$
qui redonne $\varphi_V^U$ comme la composée
$$
\mathcal{G}(V) \xrightarrow{\varphi_V}
\mathcal{F}(f^{-1}(V)) \xrightarrow{\text{rest.}}
\mathcal{F}(U)
$$
pour toute paire $(U, V)$ comme ci-dessus.
\end{lemma}

\begin{proof}
La preuve est analogue à la précédente et est omise.
\end{proof}

\section{Immersions ouvertes et (pré)faisceaux}
\label{section-open-immersions}

\noindent
Soit $X$ un espace topologique.
Soit $j : U \to X$ l'inclusion d'un sous-ensemble ouvert $U$ dans $X$.
Dans la section \ref{section-presheaves-functorial}, nous avons défini
des foncteurs $j_*$ et $j^{-1}$ tels que $j_*$ soit adjoint à droite de
$j^{-1}$. Il se trouve que, pour une immersion ouverte, $j^{-1}$ possède
un adjoint à gauche, que nous noterons $j_!$. Commençons par remarquer que
$j^{-1}$ admet une description particulièrement simple dans le cas d'une
immersion ouverte.

\begin{lemma}
\label{lemma-j-pullback}
Soit $X$ un espace topologique.
Soit $j : U \to X$ l'inclusion d'un sous-ensemble ouvert $U$ dans $X$.
\begin{enumerate}
\item Soit $\mathcal{G}$ un préfaisceau d'ensembles sur $X$.
Le préfaisceau $j_p\mathcal{G}$
(voir la section \ref{section-presheaves-functorial}) est donné par la règle
$V \mapsto \mathcal{G}(V)$ pour tout ouvert $V \subset U$.
\item Soit $\mathcal{G}$ un faisceau d'ensembles sur $X$.
Le faisceau $j^{-1}\mathcal{G}$ est donné par la règle
$V \mapsto \mathcal{G}(V)$ pour tout ouvert $V \subset U$.
\item Pour tout point $u \in U$ et tout faisceau $\mathcal{G}$ sur $X$,
on dispose d'une identification canonique des fibres
$$
j^{-1}\mathcal{G}_u = (\mathcal{G}|_U)_u = \mathcal{G}_u.
$$
\item Dans la catégorie des préfaisceaux sur $U$, nous avons $j_pj_* = \text{id}$.
\item Dans la catégorie des faisceaux sur $U$, nous avons $j^{-1}j_* = \text{id}$.
\end{enumerate}
La même description vaut pour les (pré)faisceaux de groupes abéliens,
les (pré)faisceaux de structures algébriques et les (pré)faisceaux de modules.
\end{lemma}

\begin{proof}
La limite inductive dans la définition de $j_p\mathcal{G}(V)$
porte sur la collection de tous les ouverts $W \subset X$ tels que $V \subset W$,
ordonnée par l'inclusion opposée.
Elle possède donc un plus grand élément, à savoir $V$. Cela prouve (1).
Et (2) en résulte, car la règle $V \mapsto \mathcal{G}(V)$
pour tout ouvert $V \subset U$ est manifestement un faisceau si $\mathcal{G}$ est un
faisceau. L'assertion (3) résulte de (2), car la collection
des voisinages ouverts de $u$ contenus dans $U$ est cofinale
dans la collection de tous les voisinages ouverts de $u$ dans $X$.
Les parties (4) et (5) résultent du calcul
$j^{-1}j_*\mathcal{F}(V) = j_*\mathcal{F}(V) = \mathcal{F}(V)$.

\medskip\noindent
Les mêmes arguments s'appliquent mot pour mot aux (pré)faisceaux de groupes abéliens
et aux (pré)faisceaux de structures algébriques.
\end{proof}

\begin{definition}
\label{definition-restriction}
Soit $X$ un espace topologique.
Soit $j : U \to X$ l'inclusion d'un sous-ensemble ouvert.
\begin{enumerate}
\item Soit $\mathcal{G}$ un préfaisceau d'ensembles, de groupes abéliens ou de
structures algébriques sur $X$. Le préfaisceau $j_p\mathcal{G}$ décrit
dans le lemme \ref{lemma-j-pullback} est appelé
la {\it restriction de $\mathcal{G}$ à $U$} et noté $\mathcal{G}|_U$.
\item Soit $\mathcal{G}$ un faisceau d'ensembles sur $X$, ou un faisceau de groupes
abéliens ou de structures algébriques sur $X$. Le faisceau $j^{-1}\mathcal{G}$ est appelé
la {\it restriction de $\mathcal{G}$ à $U$} et noté $\mathcal{G}|_U$.
\item Si $(X, \mathcal{O})$ est un espace annelé, le couple
$(U, \mathcal{O}|_U)$ est appelé le
{\it sous-espace ouvert de $(X, \mathcal{O})$ associé à $U$}.
\item Si $\mathcal{G}$ est un préfaisceau de $\mathcal{O}$-modules,
alors $\mathcal{G}|_U$, muni de l'application de multiplication
$\mathcal{O}|_U \times \mathcal{G}|_U \to \mathcal{G}|_U$
(voir le lemme \ref{lemma-pullback-module}),
est appelé la {\it restriction de $\mathcal{G}$ à $U$}.
\end{enumerate}
\end{definition}

\noindent
Nous laissons au lecteur la définition de la restriction des préfaisceaux
de modules. Dans cette section, nous allons donc
étudier un adjoint à gauche du foncteur de restriction.
Voici la définition dans le cas des (pré)faisceaux
d'ensembles.

\begin{definition}
\label{definition-j-shriek}
Soit $X$ un espace topologique.
Soit $j : U \to X$ l'inclusion d'un sous-ensemble ouvert.
\begin{enumerate}
\item Soit $\mathcal{F}$ un préfaisceau d'ensembles sur $U$. Nous définissons
le {\it prolongement de $\mathcal{F}$ par l'ensemble vide $j_{p!}\mathcal{F}$}
comme le préfaisceau d'ensembles sur $X$ défini par la règle
$$
j_{p!}\mathcal{F}(V) =
\left\{
\begin{matrix}
\emptyset & \text{si} & V \not \subset U \\
\mathcal{F}(V) & \text{si} & V \subset U
\end{matrix}
\right.
$$
avec les applications de restriction évidentes.
\item Soit $\mathcal{F}$ un faisceau d'ensembles sur $U$. Nous définissons
le {\it prolongement de $\mathcal{F}$ par l'ensemble vide $j_!\mathcal{F}$}
comme le faisceau associé au préfaisceau $j_{p!}\mathcal{F}$.
\end{enumerate}
\end{definition}

\begin{lemma}
\label{lemma-j-shriek}
Soit $X$ un espace topologique.
Soit $j : U \to X$ l'inclusion d'un sous-ensemble ouvert.
\begin{enumerate}
\item Le foncteur $j_{p!}$ est adjoint à gauche du
foncteur de restriction $j_p$ (voir le lemme \ref{lemma-j-pullback}).
\item Le foncteur $j_!$ est adjoint à gauche de la restriction ;
plus précisément,
$$
\Mor_{\Sh(X)}(j_!\mathcal{F}, \mathcal{G})
=
\Mor_{\Sh(U)}(\mathcal{F}, j^{-1}\mathcal{G})
=
\Mor_{\Sh(U)}(\mathcal{F}, \mathcal{G}|_U)
$$
bifonctoriellement en $\mathcal{F}$ et $\mathcal{G}$.
\item Soit $\mathcal{F}$ un faisceau d'ensembles sur $U$.
Les fibres du faisceau $j_!\mathcal{F}$ sont décrites
comme suit :
$$
j_{!}\mathcal{F}_x =
\left\{
\begin{matrix}
\emptyset & \text{si} & x \not \in U \\
\mathcal{F}_x & \text{si} & x \in U
\end{matrix}
\right.
$$
\item Dans la catégorie des préfaisceaux sur $U$, nous avons $j_pj_{p!} = \text{id}$.
\item Dans la catégorie des faisceaux sur $U$, nous avons $j^{-1}j_! = \text{id}$.
\end{enumerate}
\end{lemma}

\begin{proof}
Pour définir un morphisme de $j_{p!}\mathcal{F}$ vers $\mathcal{G}$,
il suffit de définir des applications $\mathcal{F}(V) \to \mathcal{G}(V)$
lorsque $V \subset U$, compatibles avec les applications
de restriction. Et, d'après le lemme \ref{lemma-j-pullback},
la même description vaut pour les morphismes
$\mathcal{F} \to \mathcal{G}|_U$.
L'adjonction entre $j_!$ et la restriction résulte
de ceci et des propriétés du faisceau associé.
L'identification des fibres est évidente d'après la
définition du prolongement par l'ensemble vide
et la définition d'une fibre.
Les assertions (4) et (5) résultent du calcul de la
valeur du faisceau sur tout ouvert de $U$.
\end{proof}

\noindent
Remarquons que, si $\mathcal{F}$ est un faisceau
de groupes abéliens sur $U$, alors $j_!\mathcal{F}$, tel qu'il est
défini ci-dessus, n'est en général pas un faisceau de groupes abéliens, par exemple
parce que certaines de ses fibres sont vides (et ne sont donc certainement pas
des groupes abéliens). Nous devons donc adapter la définition de
$j_!$ au type de faisceaux considéré.
La raison du choix de l'ensemble vide dans la définition du
prolongement par l'ensemble vide est qu'il s'agit de l'objet initial
de la catégorie des ensembles. Dans le cas des groupes abéliens,
nous utilisons donc $0$ (et, plus généralement, pour les faisceaux à valeurs dans
une catégorie abélienne quelconque).

\begin{definition}
\label{definition-j-shriek-structures}
Soit $X$ un espace topologique.
Soit $j : U \to X$ l'inclusion d'un sous-ensemble ouvert.
\begin{enumerate}
\item Soit $\mathcal{F}$ un préfaisceau abélien sur $U$.
Nous définissons le {\it prolongement $j_{p!}\mathcal{F}$ de $\mathcal{F}$ par $0$}
comme le préfaisceau abélien sur $X$ défini par la règle
$$
j_{p!}\mathcal{F}(V) =
\left\{
\begin{matrix}
0 & \text{si} & V \not \subset U \\
\mathcal{F}(V) & \text{si} & V \subset U
\end{matrix}
\right.
$$
avec les applications de restriction évidentes.
\item Soit $\mathcal{F}$ un faisceau abélien sur $U$. Nous définissons
le {\it prolongement $j_!\mathcal{F}$ de $\mathcal{F}$ par $0$}
comme le faisceau associé au préfaisceau abélien $j_{p!}\mathcal{F}$.
\item Soit $\mathcal{C}$ une catégorie possédant un objet initial $e$.
Soit $\mathcal{F}$ un préfaisceau sur $U$ à valeurs dans $\mathcal{C}$.
Nous définissons le {\it prolongement $j_{p!}\mathcal{F}$ de $\mathcal{F}$ par $e$}
comme le préfaisceau sur $X$ à valeurs dans $\mathcal{C}$ défini par la
règle
$$
j_{p!}\mathcal{F}(V) =
\left\{
\begin{matrix}
e & \text{si} & V \not \subset U \\
\mathcal{F}(V) & \text{si} & V \subset U
\end{matrix}
\right.
$$
avec les applications de restriction évidentes.
\item Soit $(\mathcal{C}, F)$ un type de structure algébrique
tel que $\mathcal{C}$ possède un objet initial $e$.
Soit $\mathcal{F}$ un faisceau de structures algébriques sur $U$
(du type donné). Nous définissons le
{\it prolongement $j_!\mathcal{F}$ de $\mathcal{F}$ par $e$}
comme le faisceau associé au préfaisceau $j_{p!}\mathcal{F}$
défini ci-dessus.
\item Soit $\mathcal{O}$ un préfaisceau d'anneaux sur $X$.
Soit $\mathcal{F}$ un préfaisceau de $\mathcal{O}|_U$-modules.
Dans ce cas, nous définissons le {\it prolongement par $0$}
comme le préfaisceau de $\mathcal{O}$-modules qui est égal à
$j_{p!}\mathcal{F}$ comme préfaisceau abélien muni de
l'application de multiplication
$\mathcal{O} \times j_{p!}\mathcal{F} \to j_{p!}\mathcal{F}$.
\item Soit $\mathcal{O}$ un faisceau d'anneaux sur $X$.
Soit $\mathcal{F}$ un faisceau de $\mathcal{O}|_U$-modules.
Dans ce cas, nous définissons le {\it prolongement par $0$}
comme le $\mathcal{O}$-module qui est égal à
$j_!\mathcal{F}$ comme faisceau abélien muni de
l'application de multiplication $\mathcal{O} \times j_!\mathcal{F} \to j_!\mathcal{F}$.
\end{enumerate}
\end{definition}

\noindent
Il est vrai que l'on peut définir $j_!$ dans le cadre des faisceaux
de structures algébriques (voir ci-dessous). Toutefois, l'objet
obtenu dépend du type de structures algébriques considéré.
Par exemple, si $\mathcal{O}$ est un faisceau d'anneaux
sur $U$, alors $j_{!, rings}\mathcal{O} \not = j_{!, abelian}\mathcal{O}$,
car l'objet initial de la catégorie des anneaux
est $\mathbf{Z}$, tandis que l'objet initial de la catégorie
des groupes abéliens est $0$. En particulier, le foncteur $j_!$
{\it ne commute pas avec le passage aux faisceaux d'ensembles sous-jacents},
contrairement à ce que nous avons vu jusqu'ici ! Comme d'habitude, nous traitons séparément
les (pré)faisceaux de groupes abéliens, les (pré)faisceaux de structures algébriques
et les (pré)faisceaux de modules.

\begin{lemma}
\label{lemma-j-shriek-abelian}
Soit $X$ un espace topologique.
Soit $j : U \to X$ l'inclusion d'un sous-ensemble ouvert.
Considérons les foncteurs de restriction et de prolongement
par $0$ pour les (pré)faisceaux abéliens.
\begin{enumerate}
\item Le foncteur $j_{p!}$ est adjoint à gauche du
foncteur de restriction $j_p$ (voir le lemme \ref{lemma-j-pullback}).
\item Le foncteur $j_!$ est adjoint à gauche de la restriction ;
plus précisément,
$$
\Mor_{\textit{Ab}(X)}(j_!\mathcal{F}, \mathcal{G})
=
\Mor_{\textit{Ab}(U)}(\mathcal{F}, j^{-1}\mathcal{G})
=
\Mor_{\textit{Ab}(U)}(\mathcal{F}, \mathcal{G}|_U)
$$
bifonctoriellement en $\mathcal{F}$ et $\mathcal{G}$.
\item Soit $\mathcal{F}$ un faisceau abélien sur $U$.
Les fibres du faisceau $j_!\mathcal{F}$ sont décrites
comme suit :
$$
j_{!}\mathcal{F}_x =
\left\{
\begin{matrix}
0 & \text{si} & x \not \in U \\
\mathcal{F}_x & \text{si} & x \in U
\end{matrix}
\right.
$$
\item Dans la catégorie des préfaisceaux abéliens sur $U$,
nous avons $j_pj_{p!} = \text{id}$.
\item Dans la catégorie des faisceaux abéliens sur $U$,
nous avons $j^{-1}j_! = \text{id}$.
\end{enumerate}
\end{lemma}

\begin{proof}
Démonstration omise.
\end{proof}

\begin{lemma}
\label{lemma-j-shriek-structures}
Soit $X$ un espace topologique.
Soit $j : U \to X$ l'inclusion d'un sous-ensemble ouvert.
Soit $(\mathcal{C}, F)$ un type de structure algébrique
tel que $\mathcal{C}$ possède un objet initial $e$.
Considérons les foncteurs de restriction et de prolongement
par $e$ des (pré)faisceaux de structures algébriques définis ci-dessus.
\begin{enumerate}
\item Le foncteur $j_{p!}$ est adjoint à gauche du
foncteur de restriction $j_p$ (voir le lemme \ref{lemma-j-pullback}).
\item Le foncteur $j_!$ est adjoint à gauche de la restriction ;
plus précisément,
$$
\Mor_{\Sh(X, \mathcal{C})}(j_!\mathcal{F}, \mathcal{G})
=
\Mor_{\Sh(U, \mathcal{C})}(\mathcal{F}, j^{-1}\mathcal{G})
=
\Mor_{\Sh(U, \mathcal{C})}(\mathcal{F}, \mathcal{G}|_U)
$$
bifonctoriellement en $\mathcal{F}$ et $\mathcal{G}$.
\item Soit $\mathcal{F}$ un faisceau sur $U$.
Les fibres du faisceau $j_!\mathcal{F}$ sont décrites
comme suit :
$$
j_{!}\mathcal{F}_x =
\left\{
\begin{matrix}
e & \text{si} & x \not \in U \\
\mathcal{F}_x & \text{si} & x \in U
\end{matrix}
\right.
$$
\item Dans la catégorie des préfaisceaux de structures algébriques sur $U$,
nous avons $j_pj_{p!} = \text{id}$.
\item Dans la catégorie des faisceaux de structures algébriques sur $U$,
nous avons $j^{-1}j_! = \text{id}$.
\end{enumerate}
\end{lemma}

\begin{proof}
Démonstration omise.
\end{proof}

\begin{lemma}
\label{lemma-j-shriek-modules}
Soit $(X, \mathcal{O})$ un espace annelé.
Soit $j : (U, \mathcal{O}|_U) \to (X, \mathcal{O})$
un sous-espace ouvert.
Considérons les foncteurs de restriction et de prolongement
par $0$ des (pré)faisceaux de modules définis ci-dessus.
\begin{enumerate}
\item Le foncteur $j_{p!}$ est adjoint à gauche de la restriction ;
plus précisément,
$$
\Mor_{\textit{PMod}(\mathcal{O})}(j_{p!}\mathcal{F}, \mathcal{G})
=
\Mor_{\textit{PMod}(\mathcal{O}|_U)}(\mathcal{F}, \mathcal{G}|_U)
$$
bifonctoriellement en $\mathcal{F}$ et $\mathcal{G}$.
\item Le foncteur $j_!$ est adjoint à gauche de la restriction ;
plus précisément,
$$
\Mor_{\textit{Mod}(\mathcal{O})}(j_!\mathcal{F}, \mathcal{G})
=
\Mor_{\textit{Mod}(\mathcal{O}|_U)}(\mathcal{F}, \mathcal{G}|_U)
$$
bifonctoriellement en $\mathcal{F}$ et $\mathcal{G}$.
\item Soit $\mathcal{F}$ un faisceau de $\mathcal{O}$-modules sur $U$.
Les fibres du faisceau $j_!\mathcal{F}$ sont décrites
comme suit :
$$
j_{!}\mathcal{F}_x =
\left\{
\begin{matrix}
0 & \text{si} & x \not \in U \\
\mathcal{F}_x & \text{si} & x \in U
\end{matrix}
\right.
$$
\item Dans la catégorie des faisceaux de $\mathcal{O}|_U$-modules sur $U$,
nous avons $j^{-1}j_! = \text{id}$.
\end{enumerate}
\end{lemma}

\begin{proof}
Démonstration omise.
\end{proof}

\noindent
Remarquons que, d'après les lemmes ci-dessus, les deux foncteurs
$j_*$ et $j_!$ sont des plongements pleinement fidèles de
la catégorie des faisceaux sur $U$ dans la catégorie
des faisceaux sur $X$. Ce n'est que pour le foncteur
$j_!$ que l'on peut décrire aisément l'image
essentielle.

\begin{lemma}
\label{lemma-equivalence-categories-open}
Soit $X$ un espace topologique.
Soit $j : U \to X$ l'inclusion d'un sous-ensemble ouvert.
Le foncteur
$$
j_! : \Sh(U) \longrightarrow \Sh(X)
$$
est pleinement fidèle. Son image essentielle est formée exactement
des faisceaux $\mathcal{G}$ tels que
$\mathcal{G}_x = \emptyset$ pour tout $x \in X \setminus U$.
\end{lemma}

\begin{proof}
La pleine fidélité résulte formellement de $j^{-1} j_! = \text{id}$.
Nous avons vu que tout faisceau dans l'image du foncteur possède
la propriété sur les fibres mentionnée dans le lemme. Réciproquement, supposons
que $\mathcal{G}$ possède la propriété indiquée.
Il est alors facile de vérifier que
$$
j_! j^{-1} \mathcal{G} \to \mathcal{G}
$$
est un isomorphisme sur toutes les fibres, donc un isomorphisme.
\end{proof}

\begin{lemma}
\label{lemma-equivalence-categories-open-abelian}
Soit $X$ un espace topologique.
Soit $j : U \to X$ l'inclusion d'un sous-ensemble ouvert.
Le foncteur
$$
j_! : \textit{Ab}(U) \longrightarrow \textit{Ab}(X)
$$
est pleinement fidèle. Son image essentielle est formée exactement
des faisceaux $\mathcal{G}$ tels que
$\mathcal{G}_x = 0$ pour tout $x \in X \setminus U$.
\end{lemma}

\begin{proof}
Démonstration omise.
\end{proof}

\begin{lemma}
\label{lemma-equivalence-categories-open-structures}
Soit $X$ un espace topologique.
Soit $j : U \to X$ l'inclusion d'un sous-ensemble ouvert.
Soit $(\mathcal{C}, F)$ un type de structure algébrique
tel que $\mathcal{C}$ possède un objet initial $e$.
Le foncteur
$$
j_! : \Sh(U, \mathcal{C}) \longrightarrow \Sh(X, \mathcal{C})
$$
est pleinement fidèle. Son image essentielle est formée exactement
des faisceaux $\mathcal{G}$ tels que
$\mathcal{G}_x = e$ pour tout $x \in X \setminus U$.
\end{lemma}

\begin{proof}
Démonstration omise.
\end{proof}


\begin{lemma}
\label{lemma-equivalence-categories-open-modules}
Soit $(X, \mathcal{O})$ un espace annelé.
Soit $j : (U, \mathcal{O}|_U) \to (X, \mathcal{O})$
un sous-espace ouvert.
Le foncteur
$$
j_! : \textit{Mod}(\mathcal{O}|_U) \longrightarrow \textit{Mod}(\mathcal{O})
$$
est pleinement fidèle. Son image essentielle est formée exactement
des faisceaux $\mathcal{G}$ tels que
$\mathcal{G}_x = 0$ pour tout $x \in X \setminus U$.
\end{lemma}

\begin{proof}
Démonstration omise.
\end{proof}

\begin{remark}
\label{remark-j-shriek-not-exact}
Soit $j : U \to X$ une immersion ouverte d'espaces topologiques comme ci-dessus.
Soit $x \in X$, $x \not \in U$. Soit $\mathcal{F}$ un faisceau d'ensembles
sur $U$. Alors $j_!\mathcal{F}_x = \emptyset$ d'après le lemme \ref{lemma-j-shriek}.
Ainsi, $j_!$ ne transforme pas un objet final de $\Sh(U)$
en un objet final de $\Sh(X)$, sauf si $U = X$.
D'après nos conventions de
Catégories, section \ref{categories-section-exact-functor},
cela signifie que le foncteur $j_!$ n'est pas exact à gauche
comme foncteur entre les catégories de faisceaux d'ensembles.
Nous montrerons plus loin que $j_!$ est exact sur les faisceaux abéliens ;
voir Modules, lemme \ref{modules-lemma-j-shriek-exact}.
\end{remark}







\section{Immersions fermées et (pré)faisceaux}
\label{section-closed-immersions}

\noindent
Soit $X$ un espace topologique.
Soit $i : Z \to X$ l'inclusion d'un sous-ensemble fermé $Z$ dans $X$.
Dans la section \ref{section-presheaves-functorial}, nous avons défini
des foncteurs $i_*$ et $i^{-1}$ tels que $i_*$ soit adjoint à droite de
$i^{-1}$.

\begin{lemma}
\label{lemma-stalks-closed-pushforward}
Soit $X$ un espace topologique.
Soit $i : Z \to X$ l'inclusion d'un sous-ensemble fermé $Z$ dans $X$.
Soit $\mathcal{F}$ un faisceau d'ensembles sur $Z$.
Les fibres de $i_*\mathcal{F}$ se décrivent comme suit :
$$
i_*\mathcal{F}_x =
\left\{
\begin{matrix}
\{*\} & \text{si} & x \not \in Z \\
\mathcal{F}_x & \text{si} & x \in Z
\end{matrix}
\right.
$$
où $\{*\}$ désigne un ensemble à un élément. De plus,
$i^{-1}i_* = \text{id}$ sur la catégorie des faisceaux
d'ensembles sur $Z$. Le même énoncé vaut pour les faisceaux
abéliens sur $Z$, respectivement pour les faisceaux de structures
algébriques sur $Z$, en remplaçant $\{*\}$ par $0$, respectivement
par un objet final de la catégorie des structures algébriques.
\end{lemma}

\begin{proof}
Si $x \not \in Z$, il existe des voisinages ouverts arbitrairement petits
$U$ de $x$ qui ne rencontrent pas $Z$.
Comme $\mathcal{F}$ est un faisceau, on a
$\mathcal{F}(i^{-1}(U)) = \{*\}$ pour tout tel $U$,
voir la remarque \ref{remark-confusion}. Cela démontre le premier cas.
Le second résulte du fait que, pour $z \in Z$, tout voisinage ouvert
de $z$ est de la forme $Z \cap U$ pour un ouvert $U$ de $X$.
Pour établir l'égalité $i^{-1}i_* = \text{id}$, considérons l'application
canonique $i^{-1}i_*\mathcal{F} \to \mathcal{F}$. C'est un isomorphisme
sur les fibres (d'après ce qui précède), donc un isomorphisme.

\medskip\noindent
On raisonne de même pour les faisceaux de groupes abéliens et les
faisceaux de structures algébriques.
\end{proof}


\begin{lemma}
\label{lemma-equivalence-categories-closed}
Soit $X$ un espace topologique.
Soit $i : Z \to X$ l'inclusion d'un sous-ensemble fermé.
Le foncteur
$$
i_* : \Sh(Z) \longrightarrow \Sh(X)
$$
est pleinement fidèle. Son image essentielle est exactement constituée
des faisceaux $\mathcal{G}$ tels que
$\mathcal{G}_x = \{*\}$ pour tout $x \in X \setminus Z$.
\end{lemma}

\begin{proof}
La pleine fidélité résulte formellement de $i^{-1} i_* = \text{id}$.
Nous avons vu que tout faisceau dans l'image du foncteur possède
la propriété énoncée sur les fibres. Réciproquement, supposons
que $\mathcal{G}$ possède cette propriété.
On vérifie alors aisément que
$$
\mathcal{G} \to i_* i^{-1} \mathcal{G}
$$
est un isomorphisme sur toutes les fibres, donc un isomorphisme.
\end{proof}

\begin{lemma}
\label{lemma-equivalence-categories-closed-abelian}
Soit $X$ un espace topologique.
Soit $i : Z \to X$ l'inclusion d'un sous-ensemble fermé.
Le foncteur
$$
i_* : \textit{Ab}(Z) \longrightarrow \textit{Ab}(X)
$$
est pleinement fidèle. Son image essentielle est exactement constituée
des faisceaux $\mathcal{G}$ tels que
$\mathcal{G}_x = 0$ pour tout $x \in X \setminus Z$.
\end{lemma}

\begin{proof}
Omis.
\end{proof}

\begin{lemma}
\label{lemma-equivalence-categories-closed-structures}
Soit $X$ un espace topologique.
Soit $i : Z \to X$ l'inclusion d'un sous-ensemble fermé.
Soit $(\mathcal{C}, F)$ une espèce de structure algébrique
ayant un objet final $0$. Le foncteur
$$
i_* : \Sh(Z, \mathcal{C}) \longrightarrow \Sh(X, \mathcal{C})
$$
est pleinement fidèle. Son image essentielle est exactement constituée
des faisceaux $\mathcal{G}$ tels que
$\mathcal{G}_x = 0$ pour tout $x \in X \setminus Z$.
\end{lemma}

\begin{proof}
Omis.
\end{proof}

\begin{remark}
\label{remark-i-star-not-exact}
Soit $i : Z \to X$ une immersion fermée d'espaces topologiques comme
ci-dessus. Soit $x \in X$, $x \not \in Z$. Soit $\mathcal{F}$ un faisceau
d'ensembles sur $Z$. Alors $(i_*\mathcal{F})_x = \{ * \}$
d'après le lemme \ref{lemma-stalks-closed-pushforward}.
Ainsi, si $\mathcal{F} = * \amalg *$, où
$*$ est le faisceau à un élément, alors
$i_*\mathcal{F}_x = \{*\} \not = i_*(*)_x \amalg i_*(*)_x$,
car le membre de droite est un ensemble à deux éléments.
Selon nos conventions dans
Categories, section \ref{categories-section-exact-functor},
cela signifie que le foncteur $i_*$ n'est pas exact à droite
comme foncteur entre les catégories de faisceaux d'ensembles.
En particulier, il ne peut pas avoir d'adjoint à droite, voir
Categories, lemme \ref{categories-lemma-exact-adjoint}.

\medskip\noindent
En revanche, nous verrons plus loin (voir
Modules, lemme \ref{modules-lemma-i-star-right-adjoint})
que $i_*$ est exact sur les faisceaux abéliens et possède bien un adjoint
à droite, à savoir le foncteur qui associe à un faisceau abélien sur $X$
le faisceau de ses sections à support dans $Z$.
\end{remark}

\begin{remark}
\label{remark-closed-immersion-spaces}
Nous n'avons pas étudié le lien entre les immersions fermées et les
espaces annelés. En effet, il est préférable d'aborder la notion
d'immersion fermée d'espaces annelés dans le cadre des faisceaux
quasi-cohérents ; voir Modules, section
\ref{modules-section-closed-immersion}.
\end{remark}

\section{Recollement de faisceaux}
\label{section-glueing-sheaves}

\noindent
Dans cette section, nous recollons des faisceaux définis sur les membres
d'un recouvrement de $X$. Nous commençons par les applications.

\begin{lemma}
\label{lemma-glue-maps}
Soit $X$ un espace topologique.
Soit $X = \bigcup U_i$ un recouvrement ouvert.
Soient $\mathcal{F}$ et $\mathcal{G}$ des faisceaux d'ensembles sur $X$.
Étant donnée une famille
$$
\varphi_i :
\mathcal{F}|_{U_i}
\longrightarrow
\mathcal{G}|_{U_i}
$$
d'applications de faisceaux telle que, pour tous $i, j \in I$, les
applications $\varphi_i, \varphi_j$ induisent par restriction la même application
$\mathcal{F}|_{U_i \cap U_j} \to \mathcal{G}|_{U_i \cap U_j}$,
il existe une unique application de faisceaux
$$
\varphi :
\mathcal{F}
\longrightarrow
\mathcal{G}
$$
dont la restriction à chaque $U_i$ coïncide avec $\varphi_i$.
\end{lemma}

\begin{proof}
Pour tout ouvert $U \subset X$, définissons
$$
\varphi_U : \mathcal{F}(U) \to \mathcal{G}(U), \quad
s \mapsto \varphi_U(s)
$$
où $\varphi_U(s)$ est l'unique section vérifiant
$$
(\varphi_U(s))|_{U \cap U_i}
= (\varphi_i)_{U \cap U_i}(s|_{U \cap U_i}).
$$
L'existence et l'unicité d'une telle section résultent des axiomes
des faisceaux, puisque
\begin{align*}
((\varphi_i)_{U \cap U_i}(s|_{U \cap U_i}))|_{U \cap U_i \cap U_j}
&= (\varphi_i)_{U \cap U_i \cap U_j}(s|_{U \cap U_i \cap U_j})\\
&= (\varphi_j)_{U \cap U_i \cap U_j}(s|_{U \cap U_i \cap U_j})\\
&= ((\varphi_j)_{U \cap U_j}(s|_{U \cap U_j}))|_{U \cap U_i \cap U_j}.
\end{align*}
Cette famille d'applications définit bien une application de faisceaux.
En effet, si $V \subset U \subset X$ sont des ouverts, alors
$$
(\varphi_U(s))|_V
= \varphi_V(s|_V)
$$
car, pour tout $i \in I$, on a
\begin{align*}
(\varphi_U(s))|_{V \cap U_i}
&= ((\varphi_U(s))|_{U \cap U_i})|_{V \cap U_i}\\
&= ((\varphi_i)_{U \cap U_i}(s|_{U \cap U_i}))|_{V \cap U_i}\\
&= (\varphi_i)_{V \cap U_i}(s|_{V \cap U_i})\\
&= \varphi_V(s_{V})|_{V \cap U_i}.
\end{align*}
En outre, sa restriction à chaque $U_i$ coïncide avec $\varphi_i$.
En effet, si $U \subset X$ est ouvert et
$s \in \mathcal{F}(U \cap U_i)$, alors
\begin{align*}
\varphi_{U \cap U_i}(s)
&= \varphi_{U \cap U_i}(s)|_{U \cap U_i}\\
&= (\varphi_i)_{U \cap U_i}(s|_{U \cap U_i})\\
&= (\varphi_i)_{U \cap U_i}(s).
\end{align*}
\end{proof}

\noindent
Le lemme précédent implique que, pour deux faisceaux $\mathcal{F}$ et
$\mathcal{G}$ sur l'espace topologique $X$, la règle
$$
U \longmapsto
\Mor_{\Sh(U)}(
\mathcal{F}|_U, \mathcal{G}|_U)
$$
définit un faisceau. Il s'agit d'une sorte de {\it faisceau Hom interne}.
On l'emploie rarement pour les faisceaux d'ensembles, et plus souvent
pour les faisceaux de modules ; voir Modules, section
\ref{modules-section-internal-hom}.

\medskip\noindent
Soit $X$ un espace topologique.
Soit $X = \bigcup_{i\in I} U_i$ un recouvrement ouvert.
Pour tout $i \in I$, soit $\mathcal{F}_i$ un faisceau d'ensembles sur $U_i$.
Pour toute paire $i, j \in I$, soit
$$
\varphi_{ij} :
\mathcal{F}_i|_{U_i \cap U_j}
\longrightarrow
\mathcal{F}_j|_{U_i \cap U_j}
$$
un isomorphisme de faisceaux d'ensembles. Supposons en outre que, pour
tout triplet d'indices $i, j, k \in I$, le diagramme suivant soit
commutatif :
$$
\xymatrix{
\mathcal{F}_i|_{U_i \cap U_j \cap U_k}
\ar[rr]_{\varphi_{ik}}
\ar[rd]_{\varphi_{ij}} & &
\mathcal{F}_k|_{U_i \cap U_j \cap U_k} \\
&
\mathcal{F}_j|_{U_i \cap U_j \cap U_k}
\ar[ru]_{\varphi_{jk}}
}
$$
Nous appellerons une telle famille de données
$(\mathcal{F}_i, \varphi_{ij})$
une {\it donnée de recollement de faisceaux d'ensembles relativement
au recouvrement $X = \bigcup U_i$}.

\begin{lemma}
\label{lemma-glue-sheaves}
Soit $X$ un espace topologique.
Soit $X = \bigcup_{i\in I} U_i$ un recouvrement ouvert.
Pour toute donnée de recollement $(\mathcal{F}_i, \varphi_{ij})$
de faisceaux d'ensembles relativement au recouvrement
$X = \bigcup U_i$, il existe un faisceau d'ensembles $\mathcal{F}$ sur
$X$, muni d'isomorphismes
$$
\varphi_i : \mathcal{F}|_{U_i} \to \mathcal{F}_i
$$
tels que les diagrammes
$$
\xymatrix{
\mathcal{F}|_{U_i \cap U_j} \ar[r]_{\varphi_i} \ar[d]_{\text{id}} &
\mathcal{F}_i|_{U_i \cap U_j} \ar[d]^{\varphi_{ij}} \\
\mathcal{F}|_{U_i \cap U_j} \ar[r]^{\varphi_j} &
\mathcal{F}_j|_{U_i \cap U_j}
}
$$
soient commutatifs.
\end{lemma}

\begin{proof}
Première démonstration. Nous donnons ici une formule pour l'ensemble
des sections de $\mathcal{F}$ sur un ouvert $W \subset X$. Posons
$$
\mathcal{F}(W) =
\{
(s_i)_{i \in I} \mid
s_i \in \mathcal{F}_i(W \cap U_i),
\varphi_{ij}(s_i|_{W \cap U_i \cap U_j}) = s_j|_{W \cap U_i \cap U_j}
\}.
$$
Pour $W' \subset W$, les applications de restriction sont définies en
restreignant chacun des $s_i$ à $W' \cap U_i$. La condition de faisceau
pour $\mathcal{F}$ résulte immédiatement de celle de chacun des
$\mathcal{F}_i$.

\medskip\noindent
Il reste à démontrer que $\mathcal{F}|_{U_i}$ s'envoie
isomorphiquement sur $\mathcal{F}_i$. Soit $W \subset U_i$.
Dans ce cas, la condition de la définition de $\mathcal{F}(W)$ entraîne
$s_j = \varphi_{ij}(s_i|_{W \cap U_j})$.
La commutativité des diagrammes intervenant dans la définition d'une
donnée de recollement garantit alors que l'on peut partir d'une section
{\it quelconque} $s \in \mathcal{F}_i(W)$ et obtenir une famille compatible
en posant $s_i = s$ et
$s_j = \varphi_{ij}(s_i|_{W \cap U_j})$.

\medskip\noindent
Deuxième démonstration (esquisse). Soit $\mathcal{B}$ l'ensemble des ouverts
$U \subset X$ tels que $U \subset U_i$ pour un certain $i \in I$.
Alors $\mathcal{B}$ est une base de la topologie de $X$. Pour
$U \in \mathcal{B}$, choisissons $i \in I$ tel que $U \subset U_i$ et
posons $\mathcal{F}(U) = \mathcal{F}_i(U)$.
Les isomorphismes $\varphi_{ij}$ montrent que cette prescription est
«~indépendante du choix de $i$~». À l'aide des applications de restriction
des $\mathcal{F}_i$, on constate que $\mathcal{F}$ est un faisceau sur
$\mathcal{B}$. Enfin, le lemme \ref{lemma-extend-off-basis} permet
d'étendre $\mathcal{F}$ en un unique faisceau $\mathcal{F}$ sur $X$.
\end{proof}

\begin{lemma}
\label{lemma-glue-sheaves-structures}
Soit $X$ un espace topologique.
Soit $X = \bigcup U_i$ un recouvrement ouvert.
Soit $(\mathcal{F}_i, \varphi_{ij})$ une donnée de recollement
de faisceaux de groupes abéliens, respectivement de faisceaux de
structures algébriques, respectivement de faisceaux de
$\mathcal{O}$-modules pour un faisceau d'anneaux $\mathcal{O}$ sur $X$.
Alors la construction donnée dans la démonstration du
lemme \ref{lemma-glue-sheaves} conduit respectivement à un faisceau de
groupes abéliens, à un faisceau de structures algébriques ou à un
faisceau de $\mathcal{O}$-modules.
\end{lemma}

\begin{proof}
Cela résulte du fait que, dans cette construction, l'ensemble des sections
$\mathcal{F}(W)$ sur un ouvert $W$ est donné par l'égalisateur des
applications
$$
\xymatrix{
\prod_{i \in I} \mathcal{F}_i(W \cap U_i)
\ar@<1ex>[r]
\ar@<-1ex>[r]
&
\prod_{i, j\in I} \mathcal{F}_i(W \cap U_i \cap U_j)
}
$$
Dans chacun des cas envisagés, cet égalisateur définit un objet de la
catégorie considérée dont l'ensemble sous-jacent est celui du lemme cité.
\end{proof}

\begin{lemma}
\label{lemma-mapping-property-glue}
Soit $X$ un espace topologique.
Soit $X = \bigcup_{i\in I} U_i$ un recouvrement ouvert.
Le foncteur qui associe à un faisceau d'ensembles $\mathcal{F}$
la donnée de recollement suivante
$$
(\mathcal{F}|_{U_i},
(\mathcal{F}|_{U_i})|_{U_i \cap U_j}
\to
(\mathcal{F}|_{U_j})|_{U_i \cap U_j}
)
$$
relativement au recouvrement $X = \bigcup U_i$ définit une équivalence
de catégories entre $\Sh(X)$ et la catégorie des données de recollement.
Un énoncé analogue vaut pour les faisceaux abéliens, respectivement
les faisceaux de structures algébriques, respectivement les faisceaux
de $\mathcal{O}$-modules.
\end{lemma}

\begin{proof}
Le foncteur est pleinement fidèle d'après le
lemme \ref{lemma-glue-maps}, et essentiellement surjectif (par un
quasi-inverse donné explicitement) d'après le
lemme \ref{lemma-glue-sheaves}.
\end{proof}

\noindent
Ce lemme signifie que, si le faisceau $\mathcal{F}$ a été construit
à partir de la donnée de recollement $(\mathcal{F}_i, \varphi_{ij})$
et si $\mathcal{G}$ est un faisceau sur $X$, alors un morphisme
$f : \mathcal{F} \to \mathcal{G}$ équivaut à la donnée d'une famille
de morphismes de faisceaux
$$
f_i : \mathcal{F}_i \longrightarrow \mathcal{G}|_{U_i}
$$
compatible avec les applications de recollement $\varphi_{ij}$.
De même, se donner un morphisme de faisceaux
$g : \mathcal{G} \to \mathcal{F}$ revient à se donner une famille
de morphismes de faisceaux
$$
g_i : \mathcal{G}|_{U_i} \longrightarrow \mathcal{F}_i
$$
compatible avec les applications de recollement $\varphi_{ij}$.




\begin{multicols}{2}[\section{Autres chapitres}]
\noindent
Préliminaires
\begin{enumerate}
\item \hyperref[introduction-section-phantom]{Introduction}
\item \hyperref[conventions-section-phantom]{Conventions}
\item \hyperref[sets-section-phantom]{Théorie des ensembles}
\item \hyperref[categories-section-phantom]{Catégories}
\item \hyperref[topology-section-phantom]{Topologie}
\item \hyperref[sheaves-section-phantom]{Faisceaux sur les espaces}
\item \hyperref[sites-section-phantom]{Sites et faisceaux}
\item \hyperref[stacks-section-phantom]{Champs}
\item \hyperref[fields-section-phantom]{Corps}
\item \hyperref[algebra-section-phantom]{Algèbre commutative}
\item \hyperref[brauer-section-phantom]{Groupes de Brauer}
\item \hyperref[homology-section-phantom]{Algèbre homologique}
\item \hyperref[derived-section-phantom]{Catégories dérivées}
\item \hyperref[simplicial-section-phantom]{Méthodes simpliciales}
\item \hyperref[more-algebra-section-phantom]{Compléments d'algèbre}
\item \hyperref[smoothing-section-phantom]{Lissage des morphismes d'anneaux}
\item \hyperref[modules-section-phantom]{Faisceaux de modules}
\item \hyperref[sites-modules-section-phantom]{Modules sur les sites}
\item \hyperref[injectives-section-phantom]{Objets injectifs}
\item \hyperref[cohomology-section-phantom]{Cohomologie des faisceaux}
\item \hyperref[sites-cohomology-section-phantom]{Cohomologie sur les sites}
\item \hyperref[dga-section-phantom]{Algèbre différentielle graduée}
\item \hyperref[dpa-section-phantom]{Algèbre à puissances divisées}
\item \hyperref[sdga-section-phantom]{Faisceaux différentiels gradués}
\item \hyperref[hypercovering-section-phantom]{Hyperrecouvrements}
\end{enumerate}
Schémas
\begin{enumerate}
\setcounter{enumi}{25}
\item \hyperref[schemes-section-phantom]{Schémas}
\item \hyperref[constructions-section-phantom]{Constructions de schémas}
\item \hyperref[properties-section-phantom]{Propriétés des schémas}
\item \hyperref[morphisms-section-phantom]{Morphismes de schémas}
\item \hyperref[coherent-section-phantom]{Cohomologie des schémas}
\item \hyperref[divisors-section-phantom]{Diviseurs}
\item \hyperref[limits-section-phantom]{Limites de schémas}
\item \hyperref[varieties-section-phantom]{Variétés}
\item \hyperref[topologies-section-phantom]{Topologies sur les schémas}
\item \hyperref[descent-section-phantom]{Descente}
\item \hyperref[perfect-section-phantom]{Catégories dérivées de schémas}
\item \hyperref[more-morphisms-section-phantom]{Compléments sur les morphismes}
\item \hyperref[flat-section-phantom]{Compléments sur la platitude}
\item \hyperref[groupoids-section-phantom]{Schémas en groupoïdes}
\item \hyperref[more-groupoids-section-phantom]{Compléments sur les schémas en groupoïdes}
\item \hyperref[etale-section-phantom]{Morphismes étales de schémas}
\end{enumerate}
Thèmes de théorie des schémas
\begin{enumerate}
\setcounter{enumi}{41}
\item \hyperref[chow-section-phantom]{Homologie de Chow}
\item \hyperref[intersection-section-phantom]{Théorie de l'intersection}
\item \hyperref[pic-section-phantom]{Schémas de Picard des courbes}
\item \hyperref[weil-section-phantom]{Théories cohomologiques de Weil}
\item \hyperref[adequate-section-phantom]{Modules adéquats}
\item \hyperref[dualizing-section-phantom]{Complexes dualisants}
\item \hyperref[duality-section-phantom]{Dualité pour les schémas}
\item \hyperref[discriminant-section-phantom]{Discriminants et différentes}
\item \hyperref[derham-section-phantom]{Cohomologie de de Rham}
\item \hyperref[local-cohomology-section-phantom]{Cohomologie locale}
\item \hyperref[algebraization-section-phantom]{Géométrie algébrique et formelle}
\item \hyperref[curves-section-phantom]{Courbes algébriques}
\item \hyperref[resolve-section-phantom]{Résolution des surfaces}
\item \hyperref[models-section-phantom]{Réduction semi-stable}
\item \hyperref[functors-section-phantom]{Foncteurs et morphismes}
\item \hyperref[equiv-section-phantom]{Catégories dérivées de variétés}
\item \hyperref[pione-section-phantom]{Groupes fondamentaux de schémas}
\item \hyperref[etale-cohomology-section-phantom]{Cohomologie étale}
\item \hyperref[crystalline-section-phantom]{Cohomologie cristalline}
\item \hyperref[proetale-section-phantom]{Cohomologie pro-étale}
\item \hyperref[relative-cycles-section-phantom]{Cycles relatifs}
\item \hyperref[more-etale-section-phantom]{Compléments de cohomologie étale}
\item \hyperref[trace-section-phantom]{La formule des traces}
\end{enumerate}
Espaces algébriques
\begin{enumerate}
\setcounter{enumi}{64}
\item \hyperref[spaces-section-phantom]{Espaces algébriques}
\item \hyperref[spaces-properties-section-phantom]{Propriétés des espaces algébriques}
\item \hyperref[spaces-morphisms-section-phantom]{Morphismes d'espaces algébriques}
\item \hyperref[decent-spaces-section-phantom]{Espaces algébriques décents}
\item \hyperref[spaces-cohomology-section-phantom]{Cohomologie des espaces algébriques}
\item \hyperref[spaces-limits-section-phantom]{Limites d'espaces algébriques}
\item \hyperref[spaces-divisors-section-phantom]{Diviseurs sur les espaces algébriques}
\item \hyperref[spaces-over-fields-section-phantom]{Espaces algébriques sur des corps}
\item \hyperref[spaces-topologies-section-phantom]{Topologies sur les espaces algébriques}
\item \hyperref[spaces-descent-section-phantom]{Descente et espaces algébriques}
\item \hyperref[spaces-perfect-section-phantom]{Catégories dérivées d'espaces}
\item \hyperref[spaces-more-morphisms-section-phantom]{Compléments sur les morphismes d'espaces}
\item \hyperref[spaces-flat-section-phantom]{Platitude sur les espaces algébriques}
\item \hyperref[spaces-groupoids-section-phantom]{Groupoïdes dans les espaces algébriques}
\item \hyperref[spaces-more-groupoids-section-phantom]{Compléments sur les groupoïdes dans les espaces}
\item \hyperref[bootstrap-section-phantom]{Amorçage}
\item \hyperref[spaces-pushouts-section-phantom]{Sommes amalgamées d'espaces algébriques}
\end{enumerate}
Thèmes de géométrie
\begin{enumerate}
\setcounter{enumi}{81}
\item \hyperref[spaces-chow-section-phantom]{Groupes de Chow des espaces}
\item \hyperref[groupoids-quotients-section-phantom]{Quotients de groupoïdes}
\item \hyperref[spaces-more-cohomology-section-phantom]{Compléments sur la cohomologie des espaces}
\item \hyperref[spaces-simplicial-section-phantom]{Espaces simpliciaux}
\item \hyperref[spaces-duality-section-phantom]{Dualité pour les espaces}
\item \hyperref[formal-spaces-section-phantom]{Espaces algébriques formels}
\item \hyperref[restricted-section-phantom]{Algébrisation des espaces formels}
\item \hyperref[spaces-resolve-section-phantom]{Résolution des surfaces revisitée}
\end{enumerate}
Théorie des déformations
\begin{enumerate}
\setcounter{enumi}{89}
\item \hyperref[formal-defos-section-phantom]{Théorie formelle des déformations}
\item \hyperref[defos-section-phantom]{Théorie des déformations}
\item \hyperref[cotangent-section-phantom]{Le complexe cotangent}
\item \hyperref[examples-defos-section-phantom]{Problèmes de déformation}
\end{enumerate}
Champs algébriques
\begin{enumerate}
\setcounter{enumi}{93}
\item \hyperref[algebraic-section-phantom]{Champs algébriques}
\item \hyperref[examples-stacks-section-phantom]{Exemples de champs}
\item \hyperref[stacks-sheaves-section-phantom]{Faisceaux sur les champs algébriques}
\item \hyperref[criteria-section-phantom]{Critères de représentabilité}
\item \hyperref[artin-section-phantom]{Axiomes d'Artin}
\item \hyperref[quot-section-phantom]{Espaces de Quot et de Hilbert}
\item \hyperref[stacks-properties-section-phantom]{Propriétés des champs algébriques}
\item \hyperref[stacks-morphisms-section-phantom]{Morphismes de champs algébriques}
\item \hyperref[stacks-limits-section-phantom]{Limites de champs algébriques}
\item \hyperref[stacks-cohomology-section-phantom]{Cohomologie des champs algébriques}
\item \hyperref[stacks-perfect-section-phantom]{Catégories dérivées de champs}
\item \hyperref[stacks-introduction-section-phantom]{Introduction aux champs algébriques}
\item \hyperref[stacks-more-morphisms-section-phantom]{Compléments sur les morphismes de champs}
\item \hyperref[stacks-geometry-section-phantom]{La géométrie des champs}
\end{enumerate}
Thèmes de théorie des espaces de modules
\begin{enumerate}
\setcounter{enumi}{107}
\item \hyperref[moduli-section-phantom]{Champs de modules}
\item \hyperref[moduli-curves-section-phantom]{Espaces de modules de courbes}
\end{enumerate}
Divers
\begin{enumerate}
\setcounter{enumi}{109}
\item \hyperref[examples-section-phantom]{Exemples}
\item \hyperref[exercises-section-phantom]{Exercices}
\item \hyperref[guide-section-phantom]{Guide bibliographique}
\item \hyperref[desirables-section-phantom]{Desiderata}
\item \hyperref[coding-section-phantom]{Style de programmation}
\item \hyperref[obsolete-section-phantom]{Obsolète}
\item \hyperref[fdl-section-phantom]{Licence de documentation libre GNU}
\item \hyperref[index-section-phantom]{Index généré automatiquement}
\end{enumerate}
\end{multicols}


\bibliography{my}
\bibliographystyle{amsalpha}

\end{document}
